% concatenated from NLS_with_Memory_SEP_9_2026_rev.tex


%%%%%%%%%%%%%%%%%%%%%%%%%%%%%%%%%%%%%%%%%%%%%%%%%%%%%%%%%%%%%%%%%%%%%
% !TEX encoding = UTF-8 Unicode
\documentclass[11pt,letter]{amsart}
\usepackage{amssymb,amsmath,epsfig,graphics,mathrsfs,enumerate,verbatim}
\usepackage[pagebackref,colorlinks=true,linkcolor=blue,citecolor=blue]{hyperref}

%%%%%%%%%%%%

\DeclareFontFamily{U}{mathx}{\hyphenchar\font45}
\DeclareFontShape{U}{mathx}{m}{n}{
      <5> <6> <7> <8> <9> <10>
      <10.95> <12> <14.4> <17.28> <20.74> <24.88>
      mathx10
      }{}
\DeclareSymbolFont{mathx}{U}{mathx}{m}{n}
\DeclareFontSubstitution{U}{mathx}{m}{n}
\DeclareMathAccent{\widecheck}{0}{mathx}{"71}
%%%%%%%%%%%%%%%%%%%%%%%%%%%%%
\usepackage{fancyhdr}
\pagestyle{fancy}
\fancyhead[RO,LE]{\small\thepage}
\fancyhead[LO]{\small \emph{\nouppercase{\rightmark}}}
\fancyhead[RE]{\small \emph{\nouppercase{\rightmark}}}
\fancyfoot[L,R,C]{}
\renewcommand{\headrulewidth}{1pt}
% If you don't like the horizontal line beneath the page
% number and the section's title, change the above
% \renewcommand{\headrulewidth}{1pt}  to
% \renewcommand{\headrulewidth}{0pt}
% 0pt tells Latex not to draw a line.
\renewcommand{\footrulewidth}{0pt}

%\usepackage{amssymb,amsmath}

\usepackage{hyperref}
\hypersetup{
%Colours links instead of ugly Hereafter
 colorlinks   = true,
%Colour for external hyperlinks
 urlcolor     = blue,
%Colour of internal links
 linkcolor    = blue,
%Colour of citations
 citecolor   = red ,
%Generate book marks
 bookmarksopen=true
}



%%%%%%%%%%%%%%%%%%%%%%%%%%%%%




%%%%%%%%%%%%%%%%%%%%%%%%%%%%%%
%\documentclass[14pt]{amsart}
\usepackage{amsmath}
\usepackage{amsfonts}
\usepackage{amssymb}
\usepackage{amsthm}
\usepackage{epsfig,graphics,mathrsfs}
\usepackage{graphicx}

%%%%%%%%%%%%%%%%%%%%%%%%
\makeatletter
\@namedef{subjclassname@2020}{%
  \textup{2020} Mathematics Subject Classification}
\makeatother
%%%%%%%%%%%%%%%%%%%%%%%%%%

\usepackage[usenames, dvipsnames]{color}

\usepackage{hyperref}

\usepackage[a4paper,
bindingoffset=0in,
left=0.8in,
right=0.8in,
top=1.3in,
bottom=1.3in,
footskip=.25in]{geometry}



\newcommand*\MSC[1][1991]{\par\leavevmode\hbox{%
\textit{#1 Mathematical subject classification:\ }}}
\newcommand\blfootnote[1]{%
  \begingroup
  \renewcommand\thefootnote{}\footnote{#1}%
  \addtocounter{footnote}{-1}%
  \endgroup
}

\setcounter{tocdepth}{1}


\def \e {\epsilon}
\def \N {\mathbb{N}}
\def \O {\Omega}
\def \phi {\varphi}
\def \RNu {\mathbb{R}^{d+1}}
\def \RN {\mathbb{R}_+^{d+1}}
\def \R {\mathbb{R}}
\def \l {\lambda}
\def \Gconv {G\left((p')^{-1}\circ p\right)}
\def \Geta {G_\eta}
\def \K {\mathscr{K}}
\def \LL {\mathscr L_a}
\def \Ga{\mathscr{G}_z}
\def \G{\Gamma}
\newcommand{\Ba}{\mathscr B_z^{(a)}}
\newcommand{\BaN}{\mathscr B_z^{(a),\mathrm N}}
\newcommand{\paa}{z^a \de_z}
\def \vf{\varphi}
\def \S {\mathscr{S}(\R^{N+1})}
\def \So {\mathscr{S}(\Rn)}
\newcommand{\As}{(-\mathscr A)^s}
\newcommand{\sA}{\mathscr A}
\newcommand{\Ms}{\mathscr M^{(s)}}
\newcommand{\Bpa}{\mathfrak B^\sA_{\alpha,p}}
\newcommand{\Bps}{\mathfrak B^\sA_{s,p}}
\newcommand{\Ia}{\mathscr I_\alpha}
\newcommand{\spp}{\sigma_p(\sA)}
\newcommand{\rpp}{\rho_2(\mathscr L_2)}

%%%%%%%%%
\newcommand{\Rd}{\mathbb R^d}
\newcommand{\CO}{C^\infty_0( \Omega)}
\newcommand{\Rn}{\mathbb R^n}
%\newcommand{\Rn}{\mathbb R^n}
\newcommand{\Om}{\Omega}
\newcommand{\Hn}{\mathbb H^n}
\newcommand{\Ho}{\mathbb H^1}
\newcommand{\aB}{\alpha B}
\newcommand{\eps}{\ve}
\newcommand{\BVX}{BV_X(\Omega)}
\newcommand{\p}{\partial}
\newcommand{\IO}{\int_\Omega}
\newcommand{\bG}{\mathbb {G}}
\newcommand{\bg}{\mathfrak g}
\newcommand{\bz}{\mathfrak z}
\newcommand{\bv}{\mathfrak v}
\newcommand{\Bux}{\mbox{Box}}
%\newcommand{\e}{\ve}
\newcommand{\X}{\mathcal X}
\newcommand{\Y}{\mathcal Y}
%\newcommand{\Z}{\mathcal Z}
\newcommand{\W}{\mathcal W}
\newcommand{\la}{\lambda}
\newcommand{\La}{\mathscr L}
%\newcommand{\vf}{\varphi}
\newcommand{\rhh}{|\nabla_H \rho|}
%\newcommand{\Ba}{\mathscr{B}_a}
\newcommand{\Za}{Z_\alpha}
\newcommand{\ra}{\rho_\alpha}
\newcommand{\na}{\nabla_\alpha}
\newcommand{\vt}{\vartheta}
%\newcommand{\Ga}{\mathbb G}
\newcommand{\HHa}{\mathscr H_a}
\newcommand{\HH}{\mathscr H}
%\newcommand{\paa}{z^a \p_z}



\numberwithin{equation}{section}


%\newcommand{\RN}{ \mathbb{R}^{N}}
%\newcommand{R}[1]{ \mathbb{R}^{ #1 }}
%\newcommand{\RN} {\mathbb{R}^N}
%\newcommand{\Hn}{\mathbb{H}^n}
\newcommand{\Sob}{S^{1,p}(\Omega)}
%\newcommand{\SL}{\mathcal{L}^{1,p}(\Om)}
\newcommand{\Dxk}{\frac{\partial}{\partial x_k}}
%\newcommand{\Om}{\Omega}
%\newcommand{\IO}{\int_\Omega}
\newcommand{\Co}{C^\infty_0(\Omega)}
%\newcommand{\eps}{\ve}
\newcommand{\Je}{J_\ve}
\newcommand{\beq}{\begin{equation}}
\newcommand{\bea}[1]{\begin{array}{#1} }
\newcommand{\eeq}{ \end{equation}}
\newcommand{\ea}{ \end{array}}
\newcommand{\eh}{\ve h}
\newcommand{\Dxi}{\frac{\partial}{\partial x_{i}}}
\newcommand{\Dyi}{\frac{\partial}{\partial y_{i}}}
\newcommand{\Dt}{\frac{\partial}{\partial t}}
\newcommand{\aBa}{(\alpha+1)B}
\newcommand{\GF}{\psi^{1+\frac{1}{2\alpha}}}
\newcommand{\GS}{\psi^{\frac12}}
\newcommand{\HFF}{\frac{\psi}{\rho}}
\newcommand{\HSS}{\frac{\psi}{\rho}}
\newcommand{\HFS}{\rho\psi^{\frac12-\frac{1}{2\alpha}}}
\newcommand{\HSF}{\frac{\psi^{\frac32+\frac{1}{2\alpha}}}{\rho}}
\newcommand{\AF}{\rho}
\newcommand{\AR}{\rho{\psi}^{\frac{1}{2}+\frac{1}{2\alpha}}}
\newcommand{\PF}{\alpha\frac{\psi}{|x|}}
\newcommand{\PS}{\alpha\frac{\psi}{\rho}}
\newcommand{\ds}{\displaystyle}
\newcommand{\Zt}{{\mathcal Z}^{t}}
\newcommand{\XPSI}{2\alpha\psi \begin{pmatrix} \frac{x}{|x|^2}\\ 0 \end{pmatrix} - 2\alpha\frac{{\psi}^2}{\rho^2}\begin{pmatrix} x \\ (\alpha +1)|x|^{-\alpha}y \end{pmatrix}}
\newcommand{\Z}{ \begin{pmatrix} x \\ (\alpha + 1)|x|^{-\alpha}y \end{pmatrix} }
\newcommand{\ZZ}{ \begin{pmatrix} xx^{t} & (\alpha + 1)|x|^{-\alpha}x y^{t}\\
     (\alpha + 1)|x|^{-\alpha}x^{t} y &   (\alpha + 1)^2  |x|^{-2\alpha}yy^{t}\end{pmatrix}}
%\newcommand{\I}[1,2]{ \pmb{I}_{{#1}\times {#2}}}
\newcommand{\norm}[1]{\lVert#1 \rVert}
\newcommand{\ve}{\varepsilon}
\newcommand{\Rnn}{\mathbb R^{n+1}}
\newcommand{\Rnp}{\mathbb R^{N+1}_+}
\newcommand{\B}{\mathbb{B}}
\newcommand{\Ha}{\mathbb{H}}
\newcommand{\xx}{\mathscr X}
\newcommand{\Sa}{\mathbb{S}}
\newcommand{\x}{\nabla_\mathscr X}
\newcommand{\I}{\mathscr I_{HL}}
\newcommand{\Lo}{\mathscr L^{2s,p}}
%\newcommand{\As}{(-\mathscr A)^s}
%\newcommand{\sA}{\mathscr A}
\newcommand{\Ma}{\mathscr M}
\newcommand{\Po}{\mathscr P}
\newcommand{\Ps}{\mathfrak P_s^{\sA}}
\newcommand{\In}{\mathbf 1_E}
\newcommand{\Lp}{L^p}
\newcommand{\Li}{L^\infty}
\newcommand{\Lii}{L^\infty_0}
\newcommand{\tr}{\operatorname{tr}}
\newcommand{\ssA}{\mathscr A^\star}
\newcommand{\tA}{\tilde \sA}
\newcommand{\ue}{\mathbf 1_{(-\ve,0)}}
\newcommand{\ud}{\mathbf 1_{(0,\delta)}}
\newcommand{\uex}{\mathbf 1_{(-\ve,0)}(x)}
\newcommand{\udx}{\mathbf 1_{(0,\delta)}(x)}
\newcommand{\uE}{\mathbf 1_E}
\newcommand{\sij}{\sum_{i,j=1}^m}
\newcommand{\nh}{\nabla_H}
\newcommand{\sul}{\Delta_H}
\newcommand{\nhh}{\tilde{\nabla}_H}
\newcommand{\sulp}{\Delta_{H,p}}
%%%%%%%%%%%%%%%
\newcommand{\nhn}{|\nabla_H u|^2}
\newcommand{\nhnr}{|\tilde \nabla_H u|^2}
\newcommand{\uc}{u_{,ij}}
%\newcommand{\sij}{\sum_{i,j=1}^m}
\newcommand{\sa}{\langle}
\newcommand{\da}{\rangle}


%%%%%%%%%%%%%%% begin M defs%%%%%%%%%%%%%%%%%%
\newcommand{\QH}{\boldsymbol {G\,(\mathbb{H})}}
\newcommand{\QC}{\boldsymbol {G\,(\mathbb{C})}}
\newcommand{\QK}{\boldsymbol {G\,(\mathbb{K})}}
\newcommand{\algg}{\mathfrak g}
\newcommand{\algn}{\mathfrak n}


%\newcommand{\Hn}{\mathbb{H}^n}%{\boldsymbol {G\,(\mathbb{C})}}
\newcommand{\Kn}{\mathbb{K}^n}
\newcommand{\C}{\mathbb{C}}
\newcommand{\Quat}{\mathbb{H}}
\newcommand{\Hnn}{\mathbb{H}^{n+1}}
\newcommand{\Vone}{\mathfrak h}
\newcommand{\Vtwo}{ \mathfrak v}
\newcommand{\ad}{\text{ad}}
\newcommand{\quo}{\frac{d\sigma}{|\nabla \rho|}}
\newcommand{\qui}{\psi \frac{d\sigma}{|\nabla \rho|}}
\newcommand{\fBu}{\mathscr L_{\Phi,\Psi}(u)}
\newcommand{\fB}{\mathscr L_{\Phi,\Psi}}
\newcommand{\fBa}{\mathscr L_{\Phi,\Psi,\alpha}}
\newcommand{\Gb}{\overline{\mathscr G}(X,t)}
\newcommand{\Gi}{\mathscr G(X,t)}
\DeclareMathOperator{\csch}{csch}
\DeclareMathOperator{\cotanh}{cotanh}
\DeclareMathOperator{\sech}{sech}
\newcommand{\M}{\mathcal M}
\newcommand{\Le}{\mathcal L}
\newcommand{\Si}{\mathfrak S}
\newcommand{\ps}{P_\sigma^{\Le}}

%%%%%%%%%%%%%%%%%%%%%%%%%%%%%%%%%%

\newcommand{\wh}[1]{\widehat{#1}}
\newcommand{\dd}{\,\mathrm d}
\newcommand{\ii}{\mathrm i}
\newcommand{\RNm}{\Rd\times (-\infty,0]}
\newcommand{\Sm}{\mathscr S_-(\RNu)}
\newcommand{\Dom}{\operatorname{Dom}}

%%%%%%%%%%%%%%%%%%%%%%%%%%%%%%%%%%%%%%%%%%%

\def \dive{\mathrm{div}}
%\def \tr{\mathrm{tr}}
%%%%%%%%%%%%%%%%%%%%%%%%%%%%%%%%%%%%%%%%%%%



\newtheorem{theorem}{Theorem}[section]
\newtheorem{lemma}[theorem]{Lemma}
\newtheorem{proposition}[theorem]{Proposition}
\newtheorem{corollary}[theorem]{Corollary}
\newtheorem{remark}[theorem]{Remark}
%\theoremstyle{definition}
\newtheorem{definition}[theorem]{Definition}
\newtheorem{example}[theorem]{Example}
\newtheorem{xca}[theorem]{Exercise}

\newtheorem{examples}[theorem]{Examples}
%\newtheorem{definition}[theorem]{Definition}
%\newtheorem{example}[theorem]{Example}

\numberwithin{equation}{section}
\begin{document}









%%%%%%%%%%%%%%%%%%%%%%%%%%%%%%%%%%%%%%%%%%%%%%%%%%%%%%%%%%%%%%%%%%%%%%%%%%%%%%

	
	
\title[]{An extension theory for fully fractional Schr\"odinger equations with memory} 
	
\subjclass[2020]{Primary 35Q55, 35R11; Secondary 35A08, 35A01.}

\keywords{Fully fractional Schr\"odinger operator; extension problem;
Schr\"odinger equations with memory.}

	
\date{}
	

\begin{abstract}
We develop a Caffarelli-Silvestre extension theory for the fully fractional
Schr\"odinger operator $\Le^s=(\partial_t-i\Delta_x)^s$, $0<s<1$, nonlocal in
space and time, whose Cauchy problem prescribes a past history rather than data
at a single time. Every extension theory so far rests on positivity or
sectoriality of the generator; here the semigroup is unitary and the symbol
changes sign across the characteristic paraboloid, so neither is at hand. We
construct the extension nonetheless. Its Poisson kernel, computed explicitly, is
oscillatory rather than positive; its Dirichlet-to-Neumann map is $\Le^s$; its
normalising constant is the Caffarelli-Silvestre constant times the phase
$e^{i\frac{\pi s}{2}}$; and the theory undergoes a transition at $s=\frac12$. We
then identify the intrinsic Hilbert space of histories, on which the Poisson
lifting is an isometry. This closes a circle. A lifted history is precisely an
initial datum for the singular Schr\"odinger equation with nonlinear Neumann
interaction of our companion paper, whose well-posedness theory therefore
transfers to $\Le^su=\mu|u|^{p-1}u$ with prescribed past. With our earlier work
on the Bessel operator on a half-line, the three papers form a single programme,
of which the present one is the closing step.
\end{abstract} 

	

\author{Nicola Garofalo}
\address{School of Mathematical and Statistical Sciences\\ Arizona State University}\email[Nicola Garofalo]{nicola.garofalo@asu.edu}

\author{Gigliola Staffilani}
\address{Department of Mathematics\\ Massachusetts Institute of Technology\\
Cambridge, Massachusetts, 02141\\USA}\email[Gigliola Staffilani]{gigliola@math.mit.edu}

\thanks{G. Staffilani is funded in part by the NSF grant DMS-2306378 and the Simons Foundation through the Simons
Collaboration Grant on Wave Turbulence.}
	
\maketitle
	
\tableofcontents

\section{Introduction}\label{S:intro}

Many physical systems exhibit hereditary effects in which the future evolution
depends not only on the present state, but on the entire past history of the
system. Such phenomena arise naturally in continuum mechanics, viscoelasticity,
anomalous transport, and nonlocal diffusion, and are commonly modeled by
fractional differential operators. In contrast with classical evolution
equations, where the solution is uniquely determined by its value at a single
initial time, fully fractional evolution equations require the prescription of a
\emph{history}, namely the values of the solution on the whole past time
interval. See for example \cite{CM, Mainardi, MK, DGLZ}.

Recent years have witnessed a growing interest in fully fractional evolution equations. In the parabolic setting, fully fractional heat equations with memory have been investigated from several perspectives, including global existence, blow-up, and Fujita-type phenomena, see Section \ref{S:epnonlocal} for some relevant references.

By contrast, the present work is concerned with the nonlocal dispersive operator in $\RNu$
\[
\mathcal L^s :=(\partial_t-i\Delta_x)^s,
\qquad 0<s<1,
\]
whose oscillatory character and non-sectorial nature require substantially
different analytical techniques. Our purpose is to develop an extension theory
for $\Le^s$, and to use it to solve the nonlocal Cauchy problem
\begin{equation}\label{cpmem}
\begin{cases}
\mathcal L^s u(x,t) = \mu |u(x,t)|^{p-1} u(x,t),\ \ \ x\in \Rd, t>0,
\\
u(x,t) = u_0(x,t),\ \ \ \ \ \ \ \ \ \ \ \ \ \ \ \ \ \ \ \ \ x\in \Rd, t\le 0,
\end{cases}
\end{equation}
where $\mu\in\mathbb C$, $p>1$, and the datum
$u_0:\mathbb R^d\times(-\infty,0]\longrightarrow\mathbb C$
represents the prescribed past history.

The operator $\mathcal L=\partial_t-i\Delta_x$ is the classical Schr\"odinger
operator, with unitary group $S(t)=e^{it\Delta_x}$, and its fractional powers
admit the Balakrishnan-type representation
\begin{equation}\label{E:intro-balak}
\Le^su(x,t)
=-\frac{s}{\G(1-s)}\int_0^\infty
\frac{S(\sigma)u(\cdot,t-\sigma)(x)-u(x,t)}{\sigma^{1+s}}\,d\sigma ,
\end{equation}
in which the value of $\Le^su$ at time $t$ is a superposition of the states
$S(\sigma)u(\cdot,t-\sigma)$, weighted by $\sigma^{-1-s}$.

The weight is the same one that occurs in the Balakrishnan representation of the
fully fractional heat operator $(\p_t-\Delta_x)^s$, which is likewise an
operator with memory: there too the datum is prescribed on all of $t\le0$. The
difference lies in the semigroup. In the Fourier variable $\xi$ the heat
semigroup contributes the factor $e^{-4\pi^2\sigma|\xi|^2},$ so that the state
at lag $\sigma$ enters with amplitude $\sigma^{-1-s}e^{-4\pi^2\sigma|\xi|^2}$
and \eqref{E:intro-balak} converges absolutely, with an effective cutoff at
$\sigma\sim|\xi|^{-2}$. The Schr\"odinger group contributes
$e^{-4\pi^2i\sigma|\xi|^2}$, which has modulus one: the amplitude is $\sigma^{-1-s}$ at
every frequency, there is no cutoff, and the integrals to be controlled are
oscillatory rather than absolutely convergent. This is the source of most of the
technical work below, and the reason the estimates rest on stationary-phase and
Fresnel-type arguments rather than on kernel bounds. For time-fractional
Schr\"odinger evolutions see \cite{Naber}. We emphasize that problem
\eqref{cpmem} is
fundamentally different from the much studied fractional Schr\"odinger equation
of Laskin \cite{Laskin},
\[
\partial_t u + i(-\Delta_x)^s u = 0,
\]
in which the fractional power acts in the space variables alone and the Cauchy
problem is posed with a datum at the single time $t=0$.

\medskip
\noindent\textbf{Two structural obstructions.} The analysis of $\Le^s$ is governed by two features which have no
counterpart in the elliptic or parabolic theories. Consider the real polynomial
\[
p(\xi,\tau)=2\pi\tau+4\pi^2|\xi|^2,
\]
so that the Fourier symbol of $\Le$ is $\widehat{\Le u}=i\,p\,\widehat u$. Since
$p$ changes sign across the characteristic paraboloid
\[
P=\{(\xi,\tau)\in\mathbb R^{d+1}\mid \tau = - 2\pi|\xi|^2\},
\]
the complex power $(ip)^s$ necessarily involves two different branches,
\[
(ip(\xi,\tau))^s=
\begin{cases}
e^{\,i\frac{\pi s}{2}}\,|p(\xi,\tau)|^s, & p(\xi,\tau)>0,
\\[2mm]
e^{-\,i\frac{\pi s}{2}}\,|p(\xi,\tau)|^s, & p(\xi,\tau)<0,
\end{cases}
\]
and the multiplier distinguishes the two regions of phase space separated
by $P$. Secondly, the underlying Schr\"odinger semigroup is unitary rather
than contractive, and as a consequence the extension kernel we construct is
oscillatory rather than positive. The loss of positivity changes the analytical
framework at its foundation: comparison principles, positivity arguments and
monotonicity methods, which underlie the whole parabolic theory, are
unavailable, and must be replaced by genuinely dispersive techniques.

Fractional powers of the Schr\"odinger operator were first investigated by
Samko and his collaborators \cite{Sa,CN,NS}, following ideas originating in the
classical work of M.~Riesz on fractional powers of the wave operator and their
subsequent extension by Four\`es-Bruhat to ultra-hyperbolic operators
\cite{R2,FB}. These works provide a rigorous definition of $\mathcal L^s$, but
they do not address the corresponding extension problem, which is the object of
the present paper.

%%%%%%%%%%%%%%%%%%%

\subsection{Extension problems for nonlocal operators}\label{S:epnonlocal}

One of the most successful developments in the modern analysis of nonlocal
operators is the realization of fractional powers of differential operators as
Dirichlet-to-Neumann maps of local extension problems. Beginning with the
seminal work of Caffarelli and Silvestre \cite{CS} for the fractional Laplacian,
the extension method has profoundly influenced the study of nonlocal equations,
opening the way to powerful local techniques in regularity theory, free boundary
problems, unique continuation and nonlinear analysis.

The extension philosophy has subsequently been developed in several different
directions. For the square root of the heat operator
$(\partial_t-\Delta_x)^{1/2}$, an extension already appeared, although from a
very different perspective, in the pioneering work of F. Jones \cite{Jo} on
parabolic Lipschitz spaces. The general semigroup approach to extension problems
for fractional powers of operators was later developed by Stinga and Torrea,
and subsequently generalized by Gale, Miana and Stinga to broad classes of
generators of integrated semigroups and operators with purely imaginary symbols
\cite{GMS}. In the particular setting of the fractional heat operator
$(\partial_t-\Delta_x)^s$, the extension theory and the corresponding
regularity theory were independently established by Nystr\"om-Sande and by
Stinga-Torrea \cite{NyS,ST}. The recent work \cite{FD} on the fully fractional heat equation has further illustrated the effectiveness of the positivity-based approach in the analysis of nonlinear equations with memory, including the study of blow-up and Fujita-type phenomena.
Despite their breadth, the common feature of these theories is that they are
ultimately built upon positivity or sectoriality of the underlying generator.

The one extension theory in the literature which, like ours, must dispense with
positivity is that of Enciso, Gonz\'alez and Vergara \cite{EGV}, who realized
fractional powers of the wave operator as the Dirichlet-to-Neumann map of the
Klein-Gordon equation in anti-de Sitter spacetimes. Their symbol
$|\xi|^2-\tau^2$ changes sign across the light cone, much as ours does across
the paraboloid $P$, and the same branch ambiguity has to be resolved. The two
settings then diverge: theirs is hyperbolic and the construction is driven by
the underlying Lorentzian geometry, whereas ours is dispersive and is governed
by the oscillatory Schr\"odinger semigroup, so that the extension equation is
itself of Schr\"odinger type in the half-space. We return to this comparison in
Section \ref{S:osc}.

The present paper is the third in a series devoted to the development of a local
theory for singular Schr\"odinger equations in the upper half-space and its
application to nonlocal Schr\"odinger equations through an oscillatory extension
procedure. In \cite{GS1} we developed the linear theory for the Schr\"odinger
equation associated with the Bessel operator on the half-line, establishing the
corresponding Schr\"odinger propagator, Strichartz estimates and well-posedness for the corresponding NLS.
In the subsequent paper \cite{GS2} we considered singular Schr\"odinger
equations in the upper half-space with nonlinear Neumann boundary interactions
driven by the Bessel operator, proving a complete local well-posedness theory in both
the critical and subcritical regimes. The principal objective of the present
work is to connect those results with the nonlocal memory equation
\eqref{cpmem}, by showing that $\Le^s$ admits an extension problem whose bulk
evolution is precisely the singular Schr\"odinger equation studied in
\cite{GS2}.

%%%%%%%%%%%%%%%%%%%

\subsection{Main results}\label{S:mainres}

Throughout we set
\[
a=1-2s\in(-1,1),
\qquad
\Ba=\p_{zz}+\frac az\p_z ,
\]
the latter being the Bessel operator on the half-line $\R^+_z$, and we work in
the half-space $\RN=\Rd_x\times\R^+_z$ with the weighted measure
$d\omega_a(X)=z^adxdz$. The reader should be aware that hereafter we use the notation introduced in Section \ref{S:not}. 


Our first result solves the extension problem in phase space and identifies the
Dirichlet-to-Neumann map. It is the exact dispersive analogue of the
Caffarelli-Silvestre theorem, with one essential difference: the extension
equation is not elliptic but of Schr\"odinger type, and the multiplier carries
the unimodular phase $\theta(\xi,\tau)$ of \eqref{phase}, which is what selects
the correct branch on either side of $P$. In the next statement, $K_s$ indicates the Bessel function in \eqref{Knu}.

\medskip
\noindent\textbf{Theorem A} (Theorem \ref{T:itstheFT}).
\emph{Let $0<s<1$ and $u\in\mathscr S(\RNu)$. The extension problem
\[
\p_tU-i\bigl(\Delta_xU+\Ba U\bigr)=0
\quad\text{in }\RN\times\R,
\qquad
U(x,0,t)=u(x,t),
\]
admits the phase-space representation
\[
\widehat U(\xi,z,\tau)
=\frac{2^{1-s}}{\G(s)}
\Bigl(z\,\theta(\xi,\tau)|p(\xi,\tau)|^{\frac12}\Bigr)^{s}
K_{s}\Bigl(z\,\theta(\xi,\tau)|p(\xi,\tau)|^{\frac12}\Bigr)\,\widehat u(\xi,\tau),
\]
and its weighted Neumann derivative recovers the fully fractional Schr\"odinger
operator,
\[
-\frac{2^{2s-1}\G(s)e^{i\frac{\pi s}{2}}}{\G(1-s)}\,
\lim_{z\to0^+}z^{a}\,\p_zU(\cdot,z,\cdot)=\Le^su .
\]}

\medskip
It is worth recording that the constant in Theorem A is exactly the one of
Caffarelli and Silvestre, rotated. Writing
$d_s=\frac{2^{2s-1}\G(s)}{\G(1-s)}$ for the normalising constant in the
Dirichlet-to-Neumann identity $(-\Delta)^su=-d_s\lim_{z\to0^+}z^a\p_zU$ of
\cite{CS}, the constant appearing in Theorem A is $d_s\,e^{i\frac{\pi s}{2}}$.
The phase is the same one that relates $\Le^s$ to Samko's operator
$(\Delta_x+i\p_t)^s$, and the same one which, in Remark \ref{R:admissible-mu}
below, rotates the admissible coupling constants of the nonlinear problem off
the real axis.

\medskip
The phase-space formula is explicit, but implicit in the physical variables. Our
second result inverts it, and produces the Poisson kernel of the extension. In
the elliptic and parabolic theories the corresponding kernel is a positive
approximate identity; here it is an oscillatory kernel which tensorizes into a
one-dimensional factor in the extension variable and the free Schr\"odinger
propagator in $\Rd$, see $S(x,y,t)$ in \eqref{Scla}. 

\medskip
\noindent\textbf{Theorem B} (Theorem \ref{T:invFT}).
\emph{The extension multiplier of Theorem A is the Fourier transform, in
$\mathscr S'(\RNu)$, of the kernel
$\mathbb P_a(X,Y_0,t)=P_a(z,t)\,S(x,y,t)$, where
\[
P_a(z,t)
=\frac{e^{\,i\frac{(a-1)\pi}{4}}}{2^{1-a}\G\!\left(\frac{1-a}{2}\right)}
\frac{z^{1-a}}{t^{\frac{3-a}{2}}}\,e^{i\frac{z^2}{4t}}\,\mathbf 1_{(0,\infty)}(t) .
\]
Consequently the solution of the extension problem is given by the Poisson
representation
\begin{equation}\label{U}
U(X,t)=\int_0^\infty\int_{\Rd}\mathbb P_a(X,Y_0,\sigma)\,u(y,t-\sigma)\,dy\,d\sigma .
\end{equation}}

\medskip
Formula \eqref{U} raises at once the question around which the second half of
the paper revolves. The memory problem \eqref{cpmem} prescribes a history $u_0(x,t)$,
$t\le0$; what is the corresponding initial datum for the local extension
problem? The answer is provided by the same kernel: we define the
\emph{Poisson lifting} of a history by
\begin{equation}\label{lift}
U_0(X)=\mathcal P_su_0(X)
:=\int_0^\infty\int_{\Rd}\mathbb P_a(X,Y_0,\sigma)\,u_0(y,-\sigma)\,dy\,d\sigma .
\end{equation}
The third result computes the weighted bulk energy of the lifting exactly, in
terms of the history alone. This identifies the intrinsic Hilbert space of
admissible histories, and shows that the passage from the memory problem to the
extension problem is an isometry rather than merely a bounded map.

\medskip
\noindent\textbf{Theorem C} (Theorem \ref{T:history} and Corollary \ref{C:extends}).
\emph{There is an explicit nonlocal quadratic form $Q_s$, acting on histories
through the Fourier variable in $x$ and a Mellin variable conjugate to the past
time, such that
\[
\|\mathcal P_su_0\|^2_{L^2_a(\RN)}=Q_s(u_0)
\qquad\text{for every } u_0\in\Sm .
\]
Denoting by $\mathcal H^s$ the completion of $\Sm$ under $Q_s^{1/2}$, the
Poisson lifting extends uniquely to an isometry
$\mathcal P_s:\mathcal H^s\to L^2_a(\RN)$.}

\medskip
With Theorem C the nonlinear theory of \cite{GS2} becomes available: the lifted
history is an admissible initial datum for the singular Schr\"odinger equation
in $\RN$ with nonlinear Neumann interaction,
\begin{equation}\label{cp020}
\begin{cases}
\partial_tU-i(\Delta_xU+\Ba U)=0,
&X\in\RN,\ t>0,
\\
\displaystyle
-\frac{2^{2s-1}\G(s)e^{i\pi s/2}}{\G(1-s)}\lim_{z\to0^+}z^a\partial_zU
=\mu |U(x,0,t)|^{p-1}U(x,0,t),
&x\in\Rd,\ t>0,
\\
U(X,0)=U_0(X).
\end{cases}
\end{equation}
Solving \eqref{cp020} and taking the boundary trace of the solution produces the
evolution of the memory problem. This is the content of our final result, in
which $p_c$ denotes the critical exponent \eqref{E:critical-exponents} and
$c_s$ the normalising constant of the Dirichlet-to-Neumann identity.

\medskip
\noindent\textbf{Theorem D} (Theorem \ref{T:memory-wp}).
\emph{Let $u_0\in\mathcal H^s$ be a prescribed history and $1<p\le p_c$. Then
the memory problem \eqref{cpmem} is well posed in the extension-induced mild
sense of Definition \ref{D:memory-solution}. Precisely, when $0\le a<1$: if
$p=p_c$ there is $\ve_0>0$ such that $\|u_0\|_{\mathcal H^s}\le\ve_0$ implies
the existence of a unique mild solution on $(0,T]$ for every $T<\infty$; if
$1<p<p_c$ there is $T>0$, depending on the data only through
$\|u_0\|_{\mathcal H^s}$, for which the problem has a unique mild solution on
$(0,T]$, and this solution extends to every $T<\infty$ when
$\operatorname{Im}(c_s\mu)=0$. The corresponding statements in the anomalous
range $-1<a<0$ are given in Theorem \ref{T:memory-wp}(ii).}

\medskip
The globalization hypothesis in Theorem D deserves a word, since it is not the
familiar requirement that the coupling constant be real. As we show in
Remark \ref{R:admissible-mu}, the condition $\operatorname{Im}(c_s\mu)=0$ is
equivalent to
\[
\mu\in e^{i\frac{\pi s}{2}}\,\R ,
\]
so that no real $\mu\neq0$ qualifies: the classical focusing and defocusing
constants are excluded for every $s\in(0,1)$. The admissible coupling constants
form the real line rotated by the phase which relates $\Le^s$ to Samko's
operator $(i\p_t + \Delta_x)^s$, and the rotation disappears only in the limit
$s\to0^+$.

%%%%%%%%%%%%%%%%%%%

\subsection{Strategy, and a transition at $s=\frac12$}\label{S:des}

The proofs follow the extension philosophy of \cite{CS}: the nonlocal problem
\eqref{cpmem} is not attacked directly. One first replaces $\Le^s$ by an
equivalent local equation in one additional spatial variable (Theorems A and B),
then lifts the prescribed history to an initial datum for that local problem
(Theorem C), and finally transfers the nonlinear theory of \cite{GS1,GS2} back
to the original equation (Theorem D).

Each of the three steps carries its own difficulty, and none of them is
inherited from the elliptic or parabolic theories. The first requires resolving
the branch ambiguity of $(ip)^s$ across the characteristic paraboloid, which we
do by a limiting-absorption prescription, so that the solution is selected by the
same analytic rule on both sides of $P$, and then inverting a Bessel multiplier of complex argument
whose boundedness on the characteristic rays is itself delicate; this is where
the transition at $s=\frac12$ discussed below originates. The second is an exact
computation rather than an estimate: through a Mellin analysis in the past-time
variable, the weighted bulk energy of the lifting is identified with a nonlocal
quadratic form in the history alone, and it is precisely because this is an
identity, and not an inequality, that the passage from the memory problem to the
extension problem loses nothing. The third rests on the existence of a strong
boundary trace, which does not follow from the Strichartz estimates alone:
those control the solution in $L^\infty_z$, and an essential
supremum in $z$ carries no information at the single point $z=0$. The required
continuity up to the boundary is the content of \cite[Theorems~8.1
and~8.2]{GS2}; for the free bulk evolution we reprove it, in the form in which
it is used here, in Lemma \ref{L:free-trace}.

A structural feature emerges along the way which we wish to emphasize, since it
governs the functional setting throughout: the theory undergoes a transition at
$s=\frac12$. The Bessel multiplier $M_\nu(w)=w^\nu K_\nu(w)$ that appears in
Theorem A is bounded on the characteristic ray $\arg w=-\frac\pi2$ when
$\nu\le\frac12$, whereas for $\frac12<\nu<1$ our bounds only give the growth
$|w|^{\nu-\frac12}$ there (Lemma \ref{L:MultiplierBound}). Accordingly the
extension operator is defined on all of $L^2(\RNu)$ when $0<s\le\frac12$, while
for $\frac12<s<1$ we work on the graph space $\Dom(\Le^s)$
(Theorem \ref{T:L2Extension}). For the Dirichlet-to-Neumann convergence the
roles are reversed, since there the relevant index is $1-s$: the natural
hypothesis $u\in\Dom(\Le^s)$ suffices for $\frac12\le s<1$, whereas for
$0<s<\frac12$ we require the strictly stronger condition
$u\in\Dom(\Le^{\frac14+\frac s2})$ (Theorem \ref{T:DtNL2}). This dispersive
transition is the same one identified in \cite{GS1,GS2}.

A word is in order, finally, on the sense in which the memory problem is solved.
The equation $\Le^su=\mu|u|^{p-1}u$ for $t>0$, with $u$ prescribed for $t\le0$,
is the problem that motivates this paper, and it is the problem solved when the
history is an actual function of the past, $u_0\in\Sm$. For a general element of
the history space $\mathcal H^s$, however, $u_0$ need not be a function on
$t\le0$, and no pointwise meaning can be attached to the prescription of the
past. Our formulation therefore proceeds through the extension: the history is
lifted by $\mathcal P_s$ to a datum in the weighted half-space, the nonlinear
boundary-value problem is solved there, and the future evolution is
\emph{defined} to be the boundary trace of that solution, see
Definition \ref{D:memory-solution}. Three consequences of this choice should be
stated at the outset. First, the trace so obtained satisfies the boundary
Volterra equation of Corollary \ref{C:volterra}, which is where the memory
structure of the problem becomes visible. Second, the original identity
$\Le^su=\mu|u|^{p-1}u$ is recovered as an identity in $L^2$, for positive times,
only under the additional graph regularity recorded in
Remark \ref{R:graph-formulation}. Third, the uniqueness assertion of Theorem D
is uniqueness within the class of solutions produced by the lifting; we do not
assert that an arbitrary solution of the fractional equation, in some weaker
sense, must arise in this way.

%%%%%%%%%%%%%%%%%%%

\subsection{Notation}\label{S:not}
\begin{itemize}
\item Generic points in the bulk space $\RN$ will be indicated with $X = (x,z), Y = (y,\zeta)$, etc.;
\item points of the boundary 
$\p \RN$ will be denoted by $X_0 = (x,0), Y_0 = (y,0)$, etc.; 
\item space-time points in the bulk $\RN\times \R$ will be denoted by
$(X,t), (Y,\tau)$, etc.; 
\item as customary, $\mathscr S(\R^{d+1})$ is the Schwartz class in the variables $(x,t)\in \R^{d+1}$;
\item we  denote by $\Sm$ the restrictions to the half-space $\RNm$ of the Schwartz functions in $\RNu$;
\item we indicate by  $d\omega_a(X) = z^a dx dz$ the invariant measure for the Weinstein operator $\Delta_x +\Ba$ in $\RN$. We mean by this that given $U, V\in C^\infty_0(\RN)$, we have 
 \[
\int_{\RN} \overline V(\Delta_x U + \Ba U)d\omega_a(X) = \int_{\RN} U\ \overline{(\Delta_x V + \Ba V)}d\omega_a(X);
\]
\item finally, we denote $L^2_a(\RN) =  L^2(\RN,d\omega_a(X))$. 
\end{itemize}

\subsection{Organization of the paper}

Section \ref{S:fullyfra} introduces $\Le^s$, first through the shifted
Schr\"odinger semigroup and then as a Fourier multiplier, and establishes the
equivalence of the two definitions. Section \ref{S:osc} proves Theorem A, and
Section \ref{S:appears} proves Theorem B and derives the Poisson representation
\eqref{U}. Section \ref{S:schwartz} establishes classical solvability of the
extension problem for Schwartz boundary data, and Section \ref{S:L2poi} extends
the oscillatory Poisson operator to the natural $L^2$-based setting, where the
transition at $s=\frac12$ described above takes place. Section \ref{S:inverse}
develops the potential theory associated with $\Le^s$: we construct the inverse
fractional operator, identify its fundamental solution, compute its oscillatory
Poisson lifting, and relate the latter to the boundary Duhamel operator of
\cite{GS2}. Section \ref{S:history} is devoted to the prescribed past and proves
Theorem C. Section \ref{S:wellpos} proves Theorem D. The appendix,
Section \ref{S:app}, collects the Bessel-function and oscillatory-integral
identities used throughout.


\section{The fully fractional Schr\"odinger operator}\label{S:fullyfra}

Throughout this work the notation $S(x,y,t)$ indicates the classical integral kernel of the unitary group $e^{i t\Delta_x}$,
\begin{equation}\label{Scla}
S(x,y,t)
=
(4\pi i t)^{-\frac d2}
e^{i \frac{|x-y|^2}{4t}},
\qquad t\neq 0.
\end{equation}
The solution to the Cauchy problem $\Le u = 0$, $u(x,0) =\vf(x)$, is given by
\[
u(x,t) = \int_{\Rd} S(x,y,t) \vf(y) dy.
\]
For a function $u:\R^{d+1}\to \overline \C$, its Fourier transform is
\begin{equation}\label{fou}
\hat u(\xi,\tau) = \int_{\R^{d+1}} e^{-2\pi i (\tau t + \sa\xi,x\da)} u(x,t) dx dt.
\end{equation}
Consider the semigroup $P^{\Le}_\sigma = e^{-\sigma \Le}$ associated with the Cauchy problem
\begin{equation}\label{PsS2}
\p_\sigma f = - \Le f,\ \ \ \ \ f(x,t,0) = u(x,t).
\end{equation}
For $u\in \mathscr S(\R^{d+1})$, the solution to \eqref{PsS2} is given by
\begin{align}\label{PsS}
f(x,t,\sigma) = P^{\Le}_\sigma u(x,t) & = \int_{\Rd} S(x,y,\sigma) u(y,t-\sigma) dy
 = S(\sigma)u(\cdot,t-\sigma)(x).
\end{align}
Observe that, if  
\begin{equation}\label{p}
p(\xi,\tau)= 2\pi \tau + 4\pi^2|\xi|^2,
\end{equation}
then \eqref{PsS} gives 
\begin{equation}\label{fSF}
\widehat{P^\Le_\sigma u}(\xi,\tau)  = e^{-i \sigma p(\xi,\tau)} \hat u(\xi,\tau),
\end{equation}
and therefore
\begin{equation}\label{map}
P^\Le_\sigma : \mathscr S(\R^{d+1})\ \longrightarrow\ \mathscr S(\R^{d+1}).
\end{equation}
An elementary, yet basic consequence of \eqref{fSF}
 is the following.

\begin{proposition}\label{P:L2}
The semigroup $\{P^\Le_\sigma\}_{\sigma>0}$ extends to a unitary one on $L^2(\RNu)$, i.e., one has 
\begin{equation}\label{Suni}
||P^{\mathcal L}_\sigma u||_{L^2(\R^{d+1})} =  ||u||_{L^2(\R^{d+1})}. 
\end{equation}
\end{proposition}
%Also, from \eqref{PsS} we find for any $2\le r\le \infty$:
%\[
%||P^{\mathcal L}_\sigma u(\cdot,t)||_{L^r_x} \le C(d) |\sigma|^{-d\left(\frac 12-\frac 1r\right)} ||u(\cdot,t-\sigma)||_{L^{r'}_x}.
%\]

Henceforth we denote by  
\begin{equation}\label{Lmasomeno}
L^+ = \{(\xi,\tau)\in \R^{d+1}\mid p(\xi,\tau) >0\},\ \ \ \ L^- = \{(\xi,\tau)\in \R^{d+1}\mid p(\xi,\tau) <0\}.
\end{equation}
 For a set $E$ we will denote by $\mathbf 1_E$ its indicator function.

\begin{definition}\label{D:ff} 
On functions $u\in \mathscr S(\RNu)$, we define the \emph{fully fractional} Schr\"odinger operator as  
\begin{equation}\label{fS}
\mathcal L^s u(x,t) = - \frac{s}{\G(1-s)} \int_0^\infty \frac{1}{\sigma^{1+s}} \left[P^{\Le}_\sigma u(x,t) - u(x,t)\right] d\sigma,\ \ \ \ \ 0<s<1.
\end{equation}
\end{definition}


Next, we observe that if we presently define the parabolic dilations 
\begin{equation}\label{pardila}
\delta_\la u(x,t) = u(\la x,\la^2 t),
\end{equation}
then elementary arguments show that
\begin{equation}\label{Ptdilheat}
P^\Le_\sigma(\delta_\la u)(x,t) = P^\Le_{\la^2 \sigma} u(\la x,\la^2 t).
\end{equation}
One easily obtains from \eqref{Ptdilheat}
\begin{equation}\label{dilfracheat}
\Le^s(\delta_\la u)(x,t) = \la^{2s} \Le^s u(\la x,\la^2 t),
\end{equation}  
which shows that the fully fractional Schr\"odinger operator $\Le^s$ is of order $2s$ with respect to the anisotropic parabolic dilations \eqref{pardila}. 

The following proposition identifies the Fourier symbol of the operator introduced above. In particular, it shows that the semigroup definition coincides with the natural Fourier multiplier $(ip(\xi,\tau))^s$.

\begin{proposition}\label{P:FTfp}
Let $0<s<1$. Then for $u\in \mathscr S(\R^{d+1})$ one has in $\mathscr S'(\R^{d+1})$
\begin{equation}\label{fp1}
\widehat{\mathcal L^{s} u}(\xi,\tau) = \left\{e^{-i\frac{\pi s}2} \mathbf 1_{L^-}(\xi,\tau) + e^{i\frac{\pi s}2} \mathbf 1_{L^+}(\xi,\tau)\right\}|p(\xi,\tau)|^{s} \hat u(\xi,\tau).
\end{equation}
Equivalently, we have
\begin{equation}\label{fp2}
\widehat{\mathcal L^{s} u} = (i p)^s \hat u.
\end{equation}
\end{proposition}

\begin{proof}
 We need to show that, for every $\vf\in \mathscr S(\R^{d+1})$, we have
\begin{align}\label{ex}
& \int_{\R^{d+1}} \hat \vf(\xi,\tau) \mathcal L^{s} u(\xi,\tau) d\xi d\tau
\\
& = \int_{\R^{d+1}} \vf(\xi,\tau) \left\{e^{-i\frac{\pi s}2} \mathbf 1_{L^-}(\xi,\tau) + e^{i\frac{\pi s}2} \mathbf 1_{L^+}(\xi,\tau)\right\}|p(\xi,\tau)|^{s} \hat u(\xi,\tau) d\xi d\tau.
\notag
\end{align}
From definition \eqref{fS}  and Fubini, we have
\begin{align*}
& \int_{\R^{d+1}} \hat \vf(\xi,\tau) \mathcal L^{s} u(\xi,\tau) d\xi d\tau  
 = - \frac{s}{\G(1 - s)} \int_{\R^{d+1}} \hat \vf(\xi,\tau) \int_0^\infty \frac{1}{\sigma^{1+s}} \left[P^\mathcal L_\sigma u(\xi,\tau) - u(\xi,\tau)\right] d\sigma d\xi d\tau
\\
& = - \frac{s}{\G(1 - s)} \int_{\R^{d+1}}\vf(\xi,\tau)\int_0^\infty \frac{1}{\sigma^{1+s}} \left[e^{-i \sigma p(\xi,\tau)} - 1\right]d\sigma\  \hat u(\xi,\tau)  d\xi d\tau 
\\
& = - \frac{s}{\G(1 - s)} \int_0^\infty \frac{1}{\sigma^{1+s}}  \int_{L^-} \vf(\xi,\tau) \left[e^{i \sigma |p(\xi,\tau)|} - 1\right]\hat u(\xi,\tau)  d\xi d\tau d\sigma
\\
&   - \frac{s}{\G(1 - s)} \int_0^\infty \frac{1}{\sigma^{1+s}}  \int_{L^+} \vf(\xi,\tau)\left[e^{-i \sigma p(\xi,\tau)} - 1\right]\hat u(\xi,\tau)  d\xi d\tau d\sigma
\\
& = - \frac{s}{\G(1 - s)} \int_{L^-} \vf(\xi,\tau) \hat u(\xi,\tau) \int_0^\infty \frac{1}{\sigma^{1+s}} \left[e^{i \sigma |p(\xi,\tau)|} - 1\right] d\sigma d\xi d\tau
\notag
\\
&  - \frac{s}{\G(1 - s)}  \int_{L^+} \vf(\xi,\tau) \hat u(\xi,\tau) \int_0^\infty \frac{1}{\sigma^{1+s}}\left[e^{- i \sigma p(\xi,\tau)} - 1\right]d\sigma d\xi d\tau.
\end{align*}
Applying Lemma \ref{L:aux} we obtain
 \begin{align*}
& \int_{\R^{d+1}} \hat \vf(\xi,\tau) \mathcal L^{s} u(\xi,\tau) d\xi d\tau
 = e^{-i\frac{\pi s}2} \int_{L^-} \vf(\xi,\tau) \hat u(\xi,\tau) |p(\xi,\tau)|^{s}  d\xi d\tau
\\
& + e^{i\frac{\pi s}2} \int_{L^+} \vf(\xi,\tau) \hat u(\xi,\tau) p(\xi,\tau)^{s}  d\xi d\tau.
\end{align*}
This completes the proof.



\end{proof}

\begin{remark}\label{R:not}
We note explicitly that formula \eqref{fp2} in Proposition \ref{P:FTfp} shows that 
\[
\Le^s:\mathscr S(\R^{d+1})\ \not\longrightarrow\ \mathscr S(\R^{d+1}).
\]
\end{remark}

\begin{definition}\label{D:domain}
Motivated by Proposition~\ref{P:FTfp}, we realize $\Le^s$ as a closed
operator on $L^2(\RNu)$ by setting
\[
\Dom(\Le^s)
=
\left\{u\in L^2(\RNu): |p|^s\widehat u\in L^2(\RNu)\right\},
\qquad
\widehat{\Le^su}:=(ip)^s\widehat u,
\]
where the principal branch is used. On $\mathscr S(\RNu)$ this agrees with the
semigroup definition in Definition~\ref{D:ff}, by Proposition~\ref{P:FTfp}.
We equip $\Dom(\Le^s)$ with the graph norm
\[
\|u\|_{\Dom(\Le^s)}
=
\|u\|_{L^2(\RNu)}+\|\Le^su\|_{L^2(\RNu)}.
\]
\end{definition}
For $0<s<1$, comparison of $1+|p|^s$ with $(1+|p|^2)^{s/2}$ gives
\begin{equation}\label{equigraph}
\|u\|_{\Dom(\Le^s)}
\cong \|(1+|p|^2)^{\frac s2}\widehat u\|_{L^2(\RNu)}.
\end{equation}
%%%%%%%%%%%%%%%%%%%%%%%%%%%%%%%%%%%


%%%%%%%%%%%%%%%%%%%%%%%%%%%%
 
\subsection{Comparison with Samko's definition}\label{S:samko}

Formula \eqref{fp1} shows that our definition of the nonlocal operators \eqref{fS} coincides with that adopted in \cite[Sec.4, p.296]{Sa}
\begin{equation}\label{sam1}
\Le^s u = \mathscr F^{-1}\left\{(e^{i\frac{\pi s}2} \mathbf 1_{L^+} + e^{-i\frac{\pi s}2} \mathbf 1_{L^-})|p(\xi,\tau)|^{s}\ \hat u\right\},
\end{equation}
see also \cite{Samp} for a previous study concerning the fractional powers $(\p_t -\Delta_x)^s$.
The advantage of the semigroup representation \eqref{fS}, however, is that it provides a direct operator-theoretic definition of $\Le^s$, avoiding the need to address a priori the inversion of its Fourier symbol, which is implicit in the representation
\eqref{sam1}. Subsequently, in the same book, Samko introduces the following alternative definition:
\begin{align}\label{sam}
& (\Delta_x + i \p_t)^s u(x,t) = \frac{e^{i\frac{\pi s }{2}}}{d_{d,1}(s)}\int_0^\infty \frac{1}{\sigma^{1+s}} \int_{\Rd} \tilde \Delta^{(1)}_{y,\sigma} u(x,t) S(y,\sigma) dy d\sigma, 
\end{align} 
where for arbitrary $b>1$, he sets 
\[
d_{d,1}(s) = \G(-s)\left[C^{(1)}_0 - b^{s} C^{(1)}_1\right] = \G(-s)\left[1 - b^{s}\right] = - \frac{\G(1-s)}s \left[1 - b^{s}\right],
\]
and 
\[
\tilde \Delta^{(1)}_{y,\sigma} u(x,t) = u(x-y,t-\sigma) - u(x-\sqrt b y,t - b \sigma).
\]
Using these quantities, after a direct  computation one  obtains
\begin{align*}
& e^{-i\frac{\pi s }{2}}(\Delta_x + i \p_t)^s u(x,t) = \Le^s u(x,t).
\end{align*} 
Hence, Samko's alternative definition \eqref{sam} coincides with the semigroup definition \eqref{fS}.


\vskip 0.2in

%%%%%%%%%%%%%%%%%%%%%%%%%%%%%%%%%%%%%%%

\section{The oscillatory extension problem in phase space}\label{S:osc}

In this section we solve the extension problem for the fully fractional Schr\"odinger operator $\Le^s$  by constructing an ad hoc oscillatory \emph{Poisson kernel} in phase space. The relevant analysis serves as motivation and also justification for the material in the remainder of the paper. Our construction should also be compared with the extension theory for fractional powers of the wave operator developed in \cite{EGV} by Enciso, Gonz\'alez, and Vergara, who realized the fractional wave operator as a Dirichlet-to-Neumann map associated with the Klein-Gordon equation in anti-de Sitter space. While their work concerns a hyperbolic operator and exploits the underlying Lorentzian geometry, the present setting is dispersive and is governed by the oscillatory Schr\"odinger semigroup. Accordingly, the analytical structure of the extension problem is fundamentally different.

Our approach, based on the Fourier transform, is  constructive in nature. In Theorem \ref{T:itstheFT} we first solve problem \eqref{expbS} in phase space, see \eqref{supernicehatV0}. In \eqref{onL+0} we identify the Dirichlet-to-Neumann map as the fractional power $\mathcal L^s u$ of the boundary datum. In Theorem \ref{T:invFT} in the next section, we compute the inverse Fourier transform of the multiplier in \eqref{supernicehatV0}. This step amounts to explicitly constructing the Poisson kernel $\mathbb P_a(X,Y_0,t)$ for the problem \eqref{expbS}.

To formulate the extension problem, for a given $0<s<1$ we let $a = 1-2s$, and for $z>0$ we indicate with $X= (x,z)\in \RN$. Given a function $u\in \mathscr S(\R^{d+1}_{x,t})$, we seek a function $U:\R^{d+1}_{x,t}\times \overline{\R^+_z}\to \C$, smooth for $z>0$ and continuous up to $z=0$, which, with a slight abuse of notation, we write $U(X,t) = U(x,z,t)$ instead of $U((x,t),z)$, which solves the following singular Schr\"odinger equation in the half-space $\RN\times \R_t$ with a Dirichlet condition on the boundary $\p \RN\times \R$:
\begin{equation}\label{expbS}
\begin{cases}
z^a \p_t U- i \operatorname{div}_X(z^a \nabla_X U) = 0 \ \ \ \ \ \text{in}\ \ \R^{d+1}\times\R^+,
\\
U(x,0,t) = u(x,t),\ \ \ \ \ \ \ \ \ \ \ \ \ \ \ \ \ \ \ (x,t)\in \RNu.
\end{cases}
\end{equation} 
If we denote by 
\begin{equation}\label{Ba}
\mathscr B_z^{(a)} = \p_{zz}  + \frac az \p_z,\ \ \ \ \ \ -1<a<1,
\end{equation}
the Bessel operator on the half-line $\R^+$, then 
\[
z^a \p_t U - i \operatorname{div}_X(z^a \nabla_X U) = z^a\left\{\p_t U - i \left(\Delta_x U+ \mathscr B_z^{(a)} U\right)\right\}.
\]
Therefore, we can rewrite the problem \eqref{expbS} in the following equivalent form.

\begin{definition}[The extension problem]\label{D:ep}
Given $u\in \mathscr S(\RNu)$, find a function $U\in C^\infty(\RNu\times \R^+)\cap C(\RNu\times [0,\infty))$,  such that
\begin{equation}\label{ep}
\p_t U - i \left(\Delta_x U + \mathscr B_z^{(a)} U\right) = 0,\ \ \ \ \ U(x,0,t) = u(x,t),\ \ \ \  (x,t)\in \R^{d+1}.
\end{equation}
\end{definition}
%%%%%%%%%%%%%%%%%%%%%%%%%

\subsection{The extension problem in phase space}\label{S:poiFT}

Our first main result solves the extension problem in phase space. Besides providing an explicit formula for the extension, it identifies the associated Dirichlet-to-Neumann map, thereby recovering the fractional Schr\"odinger operator.



Given $u\in \mathscr S(\R^{d+1})$, we denote by  
\[
\widehat U(\xi,z,\tau) = \int_{\R^{d+1}} e^{-2\pi i (\tau t + \sa\xi,x\da)} U(x,z,t) dx dt
\]
the partial Fourier transform with respect to $(x,t)$ of the unknown $U$ in \eqref{expbS}. This  converts the problem into the following ordinary differential equation with initial condition:
\begin{equation}\label{Bah}
\Ba \widehat U(\xi,z,\tau) = p(\xi,\tau) \widehat U(\xi,z,\tau),\ \ \ \ \ \widehat U(\xi,0,\tau) = \hat u(\xi,\tau),\ \ \ \  (\xi,\tau)\in \R^{d+1},
\end{equation}
where $p(\xi,\tau)$ is as in \eqref{p}. The ODE \eqref{Bah} admits two linearly independent
solutions in each of the regions $L_+$ and $L_-$. The relevant solution is selected
by the same analytic prescription on both sides of the characteristic set. For
$\Re p>0$ we take the normalized solution which decays as $z\to\infty$,
\[
m_s(z;p)=\frac{2^{1-s}}{\G(s)}(z\sqrt p)^sK_s(z\sqrt p),
\qquad \Re\sqrt p>0.
\]
On the negative real axis we take its limiting-absorption boundary value from the
lower half-plane,
\begin{equation}\label{E:limiting-absorption}
m_s(z;p):=\lim_{\varepsilon\downarrow0}
\frac{2^{1-s}}{\G(s)}
\bigl(z\sqrt{p-i\varepsilon}\bigr)^s
K_s\bigl(z\sqrt{p-i\varepsilon}\bigr),
\qquad p<0,
\end{equation}
where the square root is chosen with positive real part before taking the limit.
Since $\sqrt{p-i0}=-i|p|^{1/2}$ for $p<0$, this prescription gives precisely the
phase $\theta=-i$ below. In particular, the branch on $L_-$ is fixed independently
of the Dirichlet-to-Neumann identity, which will be obtained as a consequence.


Note that $p(\xi,\tau)$ is the symbol of the modified Schr\"odinger operator
\begin{equation}\label{M}
\M = - i \p_t - \Delta = e^{-i\frac{\pi}2} (\p_t - i \Delta) = e^{-i\frac{\pi}2} \Le,
\end{equation}
and therefore
\begin{equation}\label{spec}
\widehat{\M u}(\xi,\tau) = p(\xi,\tau) \hat u(\xi,\tau).
\end{equation}




With $L^+, L^-$ as in \eqref{Lmasomeno}, we now introduce the following phase function defined on $L^+\cup L^-$:
\begin{equation}\label{phase}
\theta(\xi,\tau) := e^{-i \frac{\pi}4(1-\operatorname{sgn}(p(\xi,\tau))}.
\end{equation}
Notice that 
\[
\theta(\xi,\tau) = \begin{cases}
1,\ \ \ \ \ \ \ \ \ \text{for}\ (\xi,\tau)\in L^+,
\\
-i,\ \ \ \ \ \ \  \text{when}\ (\xi,\tau)\in L^-.
\end{cases}
\]
By
$K_s$ we indicate the Bessel function of the third kind \eqref{Knu}.




\begin{theorem}\label{T:itstheFT}
The solution $U$ of the extension problem \eqref{ep} admits the following representation in phase space
\begin{equation}\label{supernicehatV0}
\widehat U(\xi,z,\tau)  = \frac{2^{1-s}}{ \G(s)} \left(z \theta(\xi,\tau)  |p(\xi,\tau)|^{1/2}\right)^s K_{s}\left(z \theta(\xi,\tau)  |p(\xi,\tau)|^{1/2}\right) \hat u(\xi,\tau)\ \ \ (\xi,\tau)\in L^+\cup L^-, 
\end{equation}
in the sense that one has
\begin{equation}\label{pointwise}
\underset{z\to 0^+}{\lim}\ \widehat U(\xi,z,\tau) = \widehat u(\xi,\tau).
\end{equation}
Moreover, $U$ satisfies the Dirichlet-to-Neumann identity:
\begin{equation}\label{onL+0}
- \frac{2^{2s-1} \G(s)e^{i\frac{\pi s}2} }{\G(1-s)}\ \underset{z\to 0^+}{\lim} z^{1-2s} \widehat U_z(\xi,z,\tau) = (i\,p(\xi,\tau))^s\,\widehat u(\xi,\tau) = \widehat{\Le^s u}(\xi,\tau). 
\end{equation}
\end{theorem}

\begin{proof}
To solve \eqref{Bah}, we distinguish the cases $(\xi,\tau)\in L^+$ and $(\xi,\tau)\in L^-$. 

\vskip 0.2in


\noindent \underline{Case $(\xi,\tau)\in L^+$}: Keeping in mind that $a=1-2s$ we rewrite \eqref{Bah} as 
\begin{equation}\label{ppav}
z^2 \widehat U_{zz} + (1-2s) z \widehat U_{z} - z^2 |p(\xi,\tau)| \widehat U = 0,\ \ \ \ \ \widehat U(\xi,0,\tau) = \hat u(\xi,\tau).
\end{equation}
This is a modified Bessel equation such as \eqref{modgenbessel}, in which we have taken:
\begin{equation}\label{alfabeta}
\alpha = s,\ \ \gamma = 1,\ \ \nu =  \pm s,\ \ \ \beta^2 = |p(\xi,\tau)|>0.
\end{equation}  
The general solution to \eqref{ppav} is given by a linear combination of the  functions $I_s$ and $K_s$, modified according to \eqref{besselcv}. We thus find  
\begin{equation}\label{notsonicehat}
\widehat U(\xi,z,\tau) = z^{s} \left[A(\xi,\tau) I_{s}(z |p(\xi,\tau)|^{1/2})+ B(\xi,\tau) K_{s}(z |p(\xi,\tau)|^{1/2})\right].
\end{equation}
By the selection principle described above, we have
$A(\xi,\tau) = 0,$
 and therefore  
\begin{align}\label{nicerhatV}
\widehat U(\xi,z,\tau) & = B(\xi,\tau) z^{s} K_{s}\left(z |p(\xi,\tau)|^{1/2}\right). 
\end{align}
Otherwise, from the asymptotic relations \eqref{ab} and \eqref{abK}, we would infer that $\widehat U(\xi,z,\tau)\to +\infty$ exponentially in $z$, as $z\to \infty$. 
From \eqref{nicerhatV} and from the asymptotic relation \eqref{zeroKnu}, we know that
\[
\widehat U(\xi,z,\tau)\ \underset{z\to 0^+}{\longrightarrow}\  B(\xi,\tau) |p(\xi,\tau)|^{-\frac{s}2} \frac{\G(s)}{2^{1-s}}.
\]
On the other hand, the Dirichlet condition $\widehat U(\xi,z,\tau) \underset{z\to 0^+}{\longrightarrow} \hat u(\xi,\tau)$ imposes at every $(\xi,\tau) \in L^+$ that
\[
B(\xi,\tau) = \frac{2^{1-s}}{ \G(s)}|p(\xi,\tau)|^{\frac{s}2} \hat u(\xi,\tau).
\]
We thus obtain from \eqref{nicerhatV}:
\begin{equation}\label{supernicehatV}
\widehat U(\xi,z,\tau)  = \frac{2^{1-s}}{ \G(s)} M_{s}\left(z |p(\xi,\tau)|^{1/2}\right) \hat u(\xi,\tau),
\end{equation}
where we have let $M_s(w) = w^{s} K_{s}(w)$. This proves \eqref{supernicehatV0} when $(\xi,\tau)\in L^+$.

We next want to understand the limit as $z\to 0^+$ of $z^{1-2s} \widehat U_z(\xi,z,\tau)$. 
Since \eqref{recuu} gives
\begin{equation*}
M_s'(w) = - w^{2s-1} M_{1-s}(w),
\end{equation*}
we obtain from \eqref{supernicehatV}:
\begin{align}\label{DtoN}
z^{1-2s} \widehat U_z(\xi,z,\tau) & = \frac{2^{1-s}}{ \G(s)}z^{1-2s} |p(\xi,\tau)|^{1/2}  M'_{s}\left(z |p(\xi,\tau)|^{1/2}\right) \hat u(\xi,\tau)
\\
& = -  \frac{2^{1-s} }{ \G(s)} z^{1-2s}  |p(\xi,\tau)|^{1/2} (|p(\xi,\tau)|^{1/2}z)^{2s-1} M_{1-s}(z |p(\xi,\tau)|^{1/2})\hat u(\xi,\tau)
\notag\\
& = -  \frac{2^{1-s}}{ \G(s)} |p(\xi,\tau)|^{s}  M_{1-s}(z |p(\xi,\tau)|^{1/2})\hat u(\xi,\tau).
\notag
\end{align}
Keeping in mind that \eqref{zeroKnu} gives
\[
M_{1-s}(z |p(\xi,\tau)|^{1/2})\ \underset{z\to 0^+}{\longrightarrow}\ 2^{-s}\G(1-s),
\]
we obtain from \eqref{DtoN}
\begin{equation}\label{DtoN2}
z^{1-2s} \widehat U_z(\xi,z,\tau)\ \underset{z\to 0^+}{\longrightarrow}\ -  \frac{2^{1-2s}\G(1-s)}{ \G(s)} |p(\xi,\tau)|^{s}\ \hat u(\xi,\tau).
\end{equation}
Comparing the right-hand side of \eqref{DtoN2} with \eqref{fp1} in Proposition \ref{P:FTfp}, we conclude that for any $(\xi,\tau)\in L^+$ we have
\begin{equation}\label{DtoN3}
z^{1-2s} \widehat U_z(\xi,z,\tau)\ \underset{z\to 0^+}{\longrightarrow}\ -  \frac{2^{1-2s}\G(1-s)}{ \G(s)} e^{-i\frac{\pi s}2} \widehat{\Le^s u}(\xi,\tau)
\end{equation}
Equivalently, we can write \eqref{DtoN3}
\begin{equation}\label{onL+}
- \frac{ \G(s)e^{i\frac{\pi s}2} }{2^{1-2s}\G(1-s)}\ \underset{z\to 0^+}{\lim} z^{1-2s} \widehat U_z(\xi,z,\tau) = \widehat{\Le^s u}(\xi,\tau), 
\end{equation}
which proves \eqref{onL+0}.

\vskip 0.2in

\noindent \underline{Case $(\xi,\tau)\in L^-$}: Here $p(\xi,\tau)<0$. Rather than
choosing a linear combination of $J_s$ and $J_{-s}$ by prescribing its Neumann
trace, we use the limiting-absorption prescription \eqref{E:limiting-absorption}.
Since
\[
\sqrt{p(\xi,\tau)-i0}=-i|p(\xi,\tau)|^{1/2},
\]
it gives directly
\begin{equation}\label{oplada}
\widehat U(\xi,z,\tau)
=
\frac{2^{1-s}}{\G(s)}
\bigl(-iz|p(\xi,\tau)|^{1/2}\bigr)^s
K_s\bigl(-iz|p(\xi,\tau)|^{1/2}\bigr)
\widehat u(\xi,\tau).
\end{equation}
This is exactly \eqref{supernicehatV0} on $L^-$, because $\theta=-i$ there.
The small-argument asymptotic
\[
w^sK_s(w)\longrightarrow 2^{s-1}\G(s),
\qquad w\to0,
\]
along the ray $\arg w=-\pi/2$ shows at once that
$\widehat U(\xi,z,\tau)\to\widehat u(\xi,\tau)$ as $z\to0^+$.

It remains to compute the Neumann trace. Using
\[
\frac{d}{dw}\bigl(w^sK_s(w)\bigr)=-w^sK_{s-1}(w)
=-w^sK_{1-s}(w)
\]
and
\[
w^{1-s}K_{1-s}(w)\longrightarrow 2^{-s}\G(1-s),
\qquad w\to0,
\]
we obtain, with $w=-iz|p|^{1/2}$,
\[
\lim_{z\to0^+}z^{1-2s}\partial_z\widehat U(\xi,z,\tau)
=-\frac{2^{1-2s}\G(1-s)}{\G(s)}
(-i)^{2s}|p(\xi,\tau)|^s\widehat u(\xi,\tau).
\]
Since $(-i)^{2s}=e^{-i\pi s}$ and, on $L^-$,
\[
\widehat{\Le^su}(\xi,\tau)=e^{-i\pi s/2}|p(\xi,\tau)|^s\widehat u(\xi,\tau),
\]
it follows that
\[
-\frac{\G(s)e^{i\pi s/2}}{2^{1-2s}\G(1-s)}
\lim_{z\to0^+}z^{1-2s}\partial_z\widehat U(\xi,z,\tau)
=\widehat{\Le^su}(\xi,\tau).
\]
Thus the Dirichlet-to-Neumann identity follows from the independently selected
outgoing boundary value, rather than being used to choose it. This proves
\eqref{supernicehatV0} and \eqref{onL+0} on $L^-$.

\end{proof}

%%%%%%%%%%%%%%%%%%%%%%%%%%%%%%%%%%%%%%%%%%%%%%%%%
\section{The Poisson kernel appears} 
\label{S:appears}

The explicit phase-space representation obtained in Theorem \ref{T:itstheFT} naturally leads to a representation in the physical variables. In our second main result, Theorem \ref{T:invFT}, we identify the oscillatory Poisson kernel associated with the extension problem by computing the inverse Fourier transform of the multiplier in \eqref{supernicehatV0}.

We will need the following auxiliary oscillatory lemma.

\begin{lemma}\label{L:critical}
Let $0<s<1$, $z>0$, and $p\in\mathbb{R}\setminus\{0\}$. Then
\[
I_s(z,p)
:=
\int_0^\infty \frac{e^{i\frac{z^2}{4\sigma}}}{\sigma^{1+s}}
e^{-ip\sigma}\,d\sigma = \begin{cases}
2^{s+1}z^{-s} e^{i \frac{\pi s}2} p^{\frac s2} K_s(z p^{1/2}),\ \ \ \ \ \ \ \ \ \ p>0,
\\
2^{s+1}z^{-s}  |p|^{\frac s2} K_s(z e^{-i \frac{\pi }2}|p|^{1/2}),\ \ \ \ \ \ p<0.
\end{cases}
\]
\end{lemma}

\begin{proof}
The change of variable $x = \frac{z^2}{4\sigma}$ gives
\[
I_s(z,p) = 2^{2s}z^{-2s} \int_0^\infty x^{s-1} e^{ix} e^{-i\frac{pz^2}{4x}} dx = 2^{2s}z^{-2s} \underset{\delta, \ve\to 0^+}{\lim}\ \int_0^\infty x^{s-1} e^{-\frac{\beta}x - \gamma x} dx,
\]
where for $\delta, \ve>0$ we let
\[
\beta = \delta + i \frac{p z^2}{4},\ \ \ \ \ \gamma = \ve - i.
\]
Since $\Re\beta=\delta>0$ and $\Re\gamma=\varepsilon>0$, we may apply Schl\"afli's identity, see \cite[9., 3.471]{GR}
\begin{equation}\label{reg}
\int_0^\infty x^{s-1} e^{-\frac{\beta}{x} - \gamma x} dx = 2 \left(\frac{\beta}{\gamma}\right)^{\frac s2} K_s(2\sqrt{\beta \gamma}),
\end{equation}
to find
\begin{equation}\label{Isp}
I_s(z,p) =  2^{2s}z^{-2s} \underset{\delta, \ve\to 0^+}{\lim}\ 2 \left(\frac{\beta}{\gamma}\right)^{\frac s2} K_s(2\sqrt{\beta \gamma})
\end{equation}
Observe that as $\delta,\ve \to 0^+$
\[
\beta=\delta+i\frac{pz^2}{4}\longrightarrow i\frac{pz^2}{4},
\qquad
\gamma=\ve-i\longrightarrow -i.
\]
Hence
\[
\frac{\beta}{\gamma}
\longrightarrow
-\frac{pz^2}{4},
\qquad
\beta\gamma
\longrightarrow
\frac{pz^2}{4}.
\]
Since the complex powers and square roots are taken with respect to the principal branch, we find 
\begin{align*}
& \frac \beta{\gamma} = \left(-\frac{pz^2}{4(1+\ve^2)} + \frac{\delta \ve}{1+\ve^2}\right) + i \left(\frac{\delta}{1+\ve^2} + \frac{\ve p z^2}{4(1+\ve^2)}\right)\ \longrightarrow\ \begin{cases} e^{i\pi} \frac{|p| z^2}4,\ \ \ \ \text{if}\ p>0,
\\
\frac{|p| z^2}4,\ \ \ \ \ \ \ \ \text{if}\ p<0,
\end{cases}
\\
& \beta \gamma = \left(\frac{pz^2}{4} +\delta \ve\right) + i \left(\ve \frac{pz^2}{4}- \delta\right)\ \longrightarrow\ \begin{cases}  \frac{|p| z^2}4,\ \ \ \ \ \ \ \ \ \ \text{if}\ p>0,
\\
e^{-i\pi} \frac{|p| z^2}4,\ \ \ \ \ \text{if}\ p<0.
\end{cases}
\end{align*}
We thus obtain from \eqref{Isp}
\[
I_s(z,p) = \begin{cases}
2^{s+1}z^{-s} e^{i \frac{\pi s}2} p^{\frac s2} K_s(z \sqrt{p}),\ \ \ \ \ \ \ \ \ \ p>0,
\\
2^{s+1}z^{-s}  |p|^{\frac s2} K_s(z e^{-i \frac{\pi }2}\sqrt{|p|}),\ \ \ \ \ p<0.
\end{cases}
\]
This gives the desired conclusion.

\end{proof}

\begin{remark}\label{R:critical}
Lemma~\ref{L:critical} may also be read in the sense of tempered
distributions in the variable $p\in \R$.
Indeed, for $\varepsilon,\delta>0$, define
\[
I_{s,\varepsilon,\delta}(z,p)
:= \int_0^\infty x^{s-1} e^{-\frac{\beta}x - \gamma x} dx,
\]
where $\beta = \beta(\delta), \gamma = \gamma(\ve)$ are as before.
Then $I_{s,\varepsilon,\delta}(z,\cdot)\in\mathscr S'(\mathbb R)$, and
the family converges in $\mathscr S'(\mathbb R)$, as
$\varepsilon,\delta\to0^+$, to the oscillatory kernel $I_s(z,p)$.
Consequently, the identity in Lemma~\ref{L:critical} holds as an identity
in $\mathscr S'(\mathbb R)$, with the pointwise formula for $p\neq0$
as its representative.
\end{remark}

We are now ready to prove the main result of this section.

\begin{theorem}\label{T:invFT}
For $z>0$ and $(\xi,\tau)\in L^+\cup L^-$, with $\theta(\xi,\tau)$ as in \eqref{phase}, let
\begin{equation}\label{ms}
m_s(z;\xi,\tau) := \frac{2^{1-s}}{ \G(s)} \left(z \theta(\xi,\tau)  |p(\xi,\tau)|^{1/2}\right)^s K_{s}\left(z \theta(\xi,\tau)  |p(\xi,\tau)|^{1/2}\right).
\end{equation}
Then, setting $a = 1-2s$, we have in $\mathscr S'(\R^{d+1})$
\begin{equation}\label{invFTFs}
\mathscr F^{-1}_{(\xi,\tau)\to (x,t)}(m_s(z;\cdot))
=
P_a(z,t)\,S(x,t),
\end{equation}
where  
\begin{equation}\label{Poi}
P_a(z,t)
=
\frac{e^{\,i\frac{(a-1)\pi}{4}}}{2^{1-a}\Gamma\!\left(\frac{1-a}{2}\right)}
\frac{z^{1-a}}{t^{\frac{3-a}{2}}}
e^{i\frac{z^2}{4t}}\mathbf 1_{(0,\infty)}(t).
\end{equation}
\end{theorem}


\begin{proof}
In order to compute $\mathscr F^{-1}(m_s(z;\cdot))$ in $\mathscr S'(\RNu)$, given $\Phi\in\mathscr S(\mathbb R^{d+1})$, we use the duality 
\[
\langle \mathscr F^{-1}(m_s(z;\cdot)),\Phi\rangle
=
\langle m_s(z;\cdot),\mathscr  F^{-1}\Phi\rangle.
\]
Thanks to Lemma~\ref{L:critical} and Remark \ref{R:critical}, the multiplier admits the oscillatory representation
\[
m_s(z;\xi,\tau)
= \frac{e^{-i\frac{\pi s}2}}{4^s\Gamma(s)} z^{2s}
\int_0^\infty
\frac{1}{t^{1+s}}
e^{i \frac{z^2}{4t}}
e^{-i t p(\xi,\tau)}\,dt,
\]
in the sense of $\mathscr S'(\RNu)$.
Let now $\eta\in C^\infty(0,\infty)$ be such that $0\le\eta\le1$, $\eta\equiv0$ on $(0,1]$ and $\eta\equiv1$ on $[2,\infty)$, and for any $R>1$ define $\chi_R(t) = \eta(Rt)\bigl(1-\eta(\frac tR)\bigr)$. Then $0\le\chi_R\le1$, $\chi_R\equiv 1$ on $[\frac2R,R]$ and $\operatorname{supp}\chi_R\subseteq[\frac1R,2R]$; in particular, for every fixed $t>0$ we have $\chi_R(t)\to 1$ as $R\to \infty$, and $\chi_R(t) = 0$ for $t \le \frac{1}{R}$ or $t\ge 2R$. Define
\[
m_{s,R}(z;\xi,\tau)
=
\frac{e^{-i\frac{\pi s}2}}{4^s\Gamma(s)} z^{2s}
\int_0^\infty
\frac{\chi_R(t)}{
t^{1+s}}
e^{i \frac{z^2}{4t}}
e^{-i t p(\xi,\tau)}\,dt.
\]
It is not difficult to check that for every fixed $z>0$,  one has
\[
m_{s,R}(z;\cdot)\ \underset{R\to \infty}{\longrightarrow}\ m_s(z;\cdot)\ \ \ \ \ \text{in}\ \mathscr S'(\RNu).
\]
Since for each fixed $R$, the integral defining $m_{s,R}(z;\cdot)$ is absolutely convergent, we can exchange order of integration. Hence,
\begin{equation}\label{reg-cutoff}
\langle \mathscr F^{-1}(m_{s,R}(z;\cdot)),\Phi\rangle
=
\frac{e^{-i\frac{\pi s}2}}{4^s\Gamma(s)} z^{2s}
\int_0^\infty
\frac{\chi_R(t)}{
t^{1+s}}
e^{i \frac{z^2}{4t}}
\langle e^{-i t p(\xi,\tau)},\mathscr F^{-1}(\Phi)\rangle
\,dt.
\end{equation}
We now claim that for any fixed $t>0$ we have
\begin{equation}\label{beauty}
\langle e^{-i t p(\xi,\tau)},\mathscr F^{-1}(\Phi)\rangle = (4\pi i t)^{-\frac d2} \int_{\Rd} e^{i\frac{|x|^2}{4t}} \Phi(x,t) dx.
\end{equation}
To check the claim, write
\[
e^{-i t p(\xi,\tau)}
=
e^{-2\pi i t\tau}\,e^{-4\pi^2 i t |\xi|^2}.
\]
First taking the inverse Fourier transform in the time variable, we find:
\[
\langle e^{-i t p(\xi,\tau)},\mathscr F^{-1}(\Phi)\rangle = \int_{\mathbb R}
e^{-2\pi i t \tau} \mathscr F^{-1}_{t\to\tau}\int_{\Rd} e^{-4\pi^2 i t |\xi|^2}
\mathscr F^{-1}_{x\to\xi}\Phi(\xi,\tau)\,d\tau
= \int_{\Rd} e^{-4\pi^2 i t |\xi|^2}
\mathscr F^{-1}_{x\to\xi}\Phi(\xi,t)\,d\xi.
\]
Hence, denoting by $\widecheck{\Phi}(x,t) = \Phi(-x,t)$, and using the standard formulas
\[
\mathscr F(\Phi) = \mathscr F^{-1}(\widecheck \Phi),\ \ \ \ \ \ \mathscr F^{-1}(\Phi) = \mathscr F(\widecheck \Phi), 
\]
we infer
\begin{align*}
\langle e^{-i t p(\xi,\tau)},\mathscr F^{-1}(\Phi)\rangle & = \int_{\Rd} e^{-4\pi^2 i t |\xi|^2}
\mathscr F_{x\to\xi}(\widecheck \Phi(\cdot,t))(\xi)\,d\xi
 = \int_{\Rd} 
\mathscr F_{x\to\xi}(S(t)\widecheck \Phi(\cdot,t))(\xi)\,d\xi
\\
& = \sa 1,\mathscr F_{x\to\xi}(S(t)\widecheck \Phi(\cdot,t))\da = \sa \delta, (S(t)\widecheck \Phi(\cdot,t))^{\widecheck{}}\ \da
\\
& = (4\pi i t)^{-\frac d2} \int_{\Rd} e^{i\frac{|x|^2}{4t}} \Phi(x,t) dx,
\end{align*}
which proves the claim \eqref{beauty}. Substituting \eqref{beauty} in \eqref{reg}, we find
\begin{align}\label{reg2}
\langle \mathscr F^{-1}(m_{s,R}),\Phi\rangle
& =
\frac{e^{-i\frac{\pi s}2}}{4^s\Gamma(s)} z^{2s}
\int_\R \int_{\Rd}
\mathbf 1_{(0,\infty)}(t)\frac{\chi_R(t)}{
t^{1+s}}
e^{i \frac{z^2}{4t}}
(4\pi i t)^{-\frac d2}  e^{i\frac{|x|^2}{4t}} \Phi(x,t) dx
\,dt
\notag\\
& = \sa \mathbf 1_{(0,\infty)}\frac{\chi_R}{
t^{1+s}}
e^{i \frac{z^2}{4t}} S(x,t),\Phi \da.
\notag\end{align}
Since we have in $\mathscr S'(\RNu)$:
\[
\mathscr F^{-1}(m_{s,R})\ \underset{R\to \infty}{\longrightarrow}\ \mathscr F^{-1}(m_{s}),\ \ \ \text{and}\ \ \ \mathbf 1_{(0,\infty)}\frac{\chi_R}{
t^{1+s}}
e^{i \frac{z^2}{4t}} S(x,t)\ \underset{R\to \infty}{\longrightarrow}\ \frac{\mathbf 1_{(0,\infty)}}{
t^{1+s}}
e^{i \frac{z^2}{4t}} S(x,t),
\]
we conclude that, with $a=1-2s$, the equations \eqref{invFTFs} and \eqref{Poi} hold. This completes the proof.



\end{proof}

Formula \eqref{invFTFs} identifies the oscillatory Poisson kernel associated with the extension problem. Consequently, the solution of Theorem \ref{T:itstheFT} admits an explicit representation in the physical variables.
Henceforth, given points in the bulk space $X = (x,z), Y = (y,\zeta)\in \RN$, we denote by $X_0 = (x,0), Y_0 = (y,0)$ their projections on the boundary manifold $\p \RN$. We are ready to introduce the main definition of this section. 

\begin{definition}\label{D:Pa}
We define the \emph{Poisson kernel} associated with the extension problem  \eqref{expbS} as the function
\begin{equation}\label{Pa}
\mathbb P_a(X,Y_0,t) := P_a(z,t) S(x,y,t)
= \frac{e^{i\frac{(a-1)\pi}{4}}}
     {2^{1-a}\Gamma\!\left(\frac{1-a}{2}\right)}
\frac{z^{1-a}}{t^{\frac{3-a}{2}}}
e^{i\frac{z^2}{4t}}\ \mathbf 1_{(0,\infty)}(t)\ S(x,y,t).
\end{equation}
\end{definition}

Combining Theorems \ref{T:itstheFT} and \ref{T:invFT}, we obtain the representation formula
\begin{equation}\label{U2}
U(x,z,t) = \int_0^\infty \int_{\Rd} \mathbb P_a(X,Y_0,\sigma) u(y,t-\sigma) dy d\sigma.
\end{equation}


\begin{lemma}\label{L:Paone}
Assume that $-1<a<1$. Then for every $z>0$ the function in \eqref{Poi} satisfies the following identity as an improper integral:
\[
\int_0^\infty P_a(z,t) dt
=
1.
\]
\end{lemma}

\begin{proof}
Set $s=\frac{1-a}{2}$,
$b=\frac{z^2}{4}$.
Then $\frac{3-a}{2}=s+1$,
and in view of \eqref{Poi} the integral can be written as
\begin{equation}\label{Paint}
\int_0^\infty P_a(z,t) dt =
\frac{e^{-i\frac{\pi s}{2}}}
     {2^{2s}\Gamma(s)}
z^{2s}
\int_0^\infty t^{-s-1}e^{i b/t}\,dt.
\end{equation}
By the change of variable
$u=\frac{b}{t}$ we find
\begin{align*}
\int_0^\infty t^{-s-1}e^{ib/t}\,dt
&=
\frac{2^{2s}}{z^{2s}}
\int_0^\infty u^{s-1}e^{iu}\,du.
\end{align*}
Since $0<\Re(s)<1$, formula \eqref{ei} in Lemma \ref{L:cis} yields
\[
\int_0^\infty u^{s-1}e^{iu}\,du
=
e^{i\frac{\pi s}2}\Gamma(s).
\]
Substituting these identities in \eqref{Paint} we reach the desired conclusion.


\end{proof}


From Lemma \ref{L:Paone} and tensorization we obtain.
 
\begin{corollary}\label{C:Paone}
We have the following oscillatory integral
\begin{equation}\label{intPa}
\int_{\RNu} \mathbb P_a(X,Y_0,t) dy dt = 1,\ \ \ \ \ \ \ \ \ \forall X\in \RN.
\end{equation} 
\end{corollary}

We close this section with the following remarkable consequence of Lemma \ref{L:Abel}. We will need it in Section \ref{P:liftingEs}.

\begin{lemma}\label{L:abepoi}
For $0<s<1, a = 1-2s, t>0$, and $z>0$ we have:
\begin{equation}\label{abepoi}
\frac{1}{\G(s)} \int_0^\infty \frac{\mathbf 1_{(0,\infty)}(t-\sigma)}{(t-\sigma)^{1-s}} P_a(z,\sigma) d\sigma = \frac{1}{\G(s)} \frac{e^{i \frac{z^2}{4t}}}{t^{1-s}}.
\end{equation}
\end{lemma}

\begin{proof} Keeping \eqref{invFTFs} and $a = 1-2s$ in mind, we have
\begin{align*}
& \frac{1}{\G(s)} \int_0^\infty \frac{\mathbf 1_{(0,\infty)}(t-\sigma)}{(t-\sigma)^{1-s}} P_a(z,\sigma) d\sigma = \frac{e^{-i\frac{\pi s}{2}}}{4^s\G(s)^2}
z^{2s} \int_0^t 
\frac{e^{i\frac{z^2}{4\sigma}}}{(t-\sigma)^{1-s} \sigma^{1+s}} d\sigma
\\
& = \frac{e^{-i\frac{\pi s}{2}}}{4^s\G(s)^2}
z^{2s}  4^s\,\G(s)\,e^{i\frac{\pi s}{2}}\,t^{s-1}\,z^{-2s}
e^{i \frac{z^2}{4t}} = \frac{1}{\G(s)} \frac{e^{i \frac{z^2}{4t}}}{t^{1-s}},
\end{align*}
which proves the lemma.

\end{proof}


%%%%%%%%%%%%%%%%%%%%%%%%%%%%%%%%%%%%

\vskip 0.2in

\subsection{Poisson kernel and Bessel intertwining}\label{S:ep}

For the sake of completeness, in this short subsection we present, in the spirit of the work of Caffarelli and Silvestre \cite[Sections 2.2 \& 2.3]{CS}, a second approach to the extension problem based on the intertwining properties of the Bessel operators. Since it presumes the knowledge of the propagator \eqref{grandeSaS0} from \cite{GS1, GS2},  
this approach is not as constructive as the one developed in Section \ref{S:poiFT}, and it also implicitly hinges on the intertwining properties of the Bessel operators on the half-line first identified in the works of Weinstein \cite{We48, We2} and Muckenhoupt and Stein \cite{MS}. In our opinion, one remarkable aspect of this indirect approach is that the Poisson kernel - an object naturally associated with the Dirichlet problem - is almost axiomatically derived once one knows the fundamental solution of the Cauchy problem in the half-space with vanishing weighted Neumann condition:
\begin{equation}\label{eppn}
\begin{cases}
\p_t U - i \left( \Delta_x U+\mathscr B_z^{(a)} U\right) = 0,\ \ \ (X,t)\in \R_+^{d+1}\times \R
\\
\underset{z\to 0^+}{\lim} z^a \p_z U(X,t) = 0,\ \  \ \ U(X,0) = U_0(X).
\end{cases}
\end{equation}
Now, the fundamental solution of  \eqref{eppn} was constructed in \cite{GS2} and it is given by the kernel:
\begin{align}\label{grandeSaS0}
\mathbb S_a(X,Y,t)
= \frac{e^{-i\,\operatorname{sgn}(t)\,\frac{(d+a+1)\pi}{4}}}{2^{\frac{2d+a+1}{2}}\pi^{\frac d2}}  
\frac{e^{i\,\frac{z^2+\zeta^2+|x-y|^2}{4t}}}{|t|^{\frac{d+a+1}{2}}}\left(\frac{z\zeta}{2|t|}\right)^{\frac{1-a}{2}}
J_{\frac{a-1}{2}}\!\left(\frac{z\zeta}{2|t|}\right),
\qquad t\neq 0.
\end{align}
We note explicitly the tensorization property
\begin{equation}\label{Sa}
\mathbb S_a(X,Y,t) =  S_a(z,\zeta,t) S(x,y,t),
\end{equation}
where $S(x,y,t)$ is the classical propagator \eqref{Scla}, whereas we have denoted by 
\begin{equation}\label{SaS}
S_a(z,\zeta,t)
=
\frac{e^{-i\,\operatorname{sgn}(t)\,\frac{(a+1)\pi}{4}}}{(2|t|)^{\frac{a+1}{2}}}
\left(\frac{z\zeta}{2|t|}\right)^{\frac{1-a}{2}}
J_{\frac{a-1}{2}}\!\left(\frac{z\zeta}{2|t|}\right)
\,e^{i \frac{z^2+\zeta^2}{4t}}, \quad t \neq 0,
\end{equation}
the fundamental solution, constructed in \cite{GS1},  of the unitary group $e^{i t \Ba}$ on the measure space $(\R^+_z, z^a dz)$, with zero Neumann condition.
Note that \eqref{Sa} implies the Neumann condition
\begin{equation}\label{Sazeron}
\underset{z\to 0^+}{\lim} z^a \p_z \mathbb S_a(X,Y,t) = 0.
\end{equation}

We will need the following elementary, yet important fact.
\begin{lemma}\label{L:S-a}
One has for $t\not= 0$
\[
S_{-a}(z,0,t) 
=
\frac{2^a\,|t|^{\frac{a-1}{2}}}
     {\Gamma\!\left(\frac{1-a}{2}\right)}
\,e^{-i\,\operatorname{sgn}(t)\,\frac{(1-a)\pi}{4}}
\,e^{i\frac{z^2}{4t}}.
\]
\end{lemma}


\begin{proof}
From \eqref{SaS0} we have
\[
S_{-a}(z,\zeta,t)
=
\frac{e^{-i\,\operatorname{sgn}(t)\,\frac{(1-a)\pi}{4}}}
     {(2|t|)^{\frac{1-a}{2}}}
\left(\frac{z\zeta}{2|t|}\right)^{\frac{a+1}{2}}
J_{-\frac{a+1}{2}}\!\left(\frac{z\zeta}{2|t|}\right)
e^{i\frac{z^{2}+\zeta^{2}}{4t}},
\qquad t\neq 0.
\]
The asymptotic of the Bessel function,
\[
J_\nu(x)
=
\frac{(x/2)^\nu}{\Gamma(\nu+1)}
+o\!\left(x^\nu\right),
\qquad x\to0,
\]
with \(\nu=-\frac{a+1}{2}\), immediately gives
\[
\lim_{\zeta\to0^+}
\left(\frac{z\zeta}{2|t|}\right)^{\frac{a+1}{2}}
J_{-\frac{a+1}{2}}
\!\left(\frac{z\zeta}{2|t|}\right)
=
\frac{2^{\frac{a+1}{2}}}
     {\Gamma\!\left(\frac{1-a}{2}\right)},
\]
and therefore the desired conclusion follows.





\end{proof}

\begin{lemma}\label{L:Smenoa}
Let \(a>-1\) and \(t\neq 0\). Then we define the boundary kernel 
\[
P_{a}(z,t) := -i\,z^{-a}\,\partial_z S_{-a}(z,0,t)
=
\frac{2^{a-1}\operatorname{sgn}(t)}{\Gamma\!\left(\frac{1-a}{2}\right)}
\,z^{1-a}\,|t|^{\frac{a-3}{2}}
\,e^{-i\,\operatorname{sgn}(t)\,\frac{(1-a)\pi}{4}}
\,e^{i\frac{z^2}{4t}}.
\]
\end{lemma}

\begin{proof}
The claimed expression of $-i\,z^{-a}\,\partial_z S_{-a}(z,0,t)$ is an elementary consequence of Lemma \ref{L:S-a}, and we thus omit the relevant details.


\end{proof}

Once Lemma \ref{L:Smenoa} is available, one obtains the Poisson kernel $\mathbb P_a(X,Y_0,t)$ as in \eqref{Pa} exploiting the same tensorization principle which underlies \eqref{Sa}. We remark that, even if Lemma \ref{L:Smenoa} seemingly introduces \emph{in one stroke} the function which in \eqref{invFTFs} of Theorem \ref{T:invFT} was derived by a delicate computation, the simplification is only apparent. The hidden difficulty lies in the computation of the propagator \eqref{SaS} which is found in our work \cite{GS1}, which is here taken for granted. 


The following result reflects the intertwining properties of the Bessel operators.

\begin{lemma}\label{L:Pa-solves-equation}
For any fixed $Y_0\in \p \RN$, the smooth function $(X,t)\to \mathbb P_a(X,Y_0,t)$ solves the equation 
\[
\partial_t \mathbb P_a - i\bigl(\Delta_x \mathbb P_a+\mathscr B_z^{(a)}\mathbb P_a\bigr)=0
\]
in the bulk space $\RN\times (0,\infty)$.
\end{lemma}

\begin{proof}
From \eqref{Pa} we have
\[
\partial_t \mathbb P_a - i\bigl(\Delta_x \mathbb P_a+\mathscr B_z^{(a)}\mathbb P_a\bigr)
=
\bigl(\partial_t P_a-i\mathscr B_z^{(a)}P_a\bigr)S
+
P_a \bigl(\partial_t S-i\Delta_x S\bigr).
\]
Since
$\partial_t S-i\Delta_x S=0$,
 it remains to show that
\[
\partial_t P_a-i\mathscr B_z^{(a)}P_a=0.
\]
Set $C_a=\frac{e^{i\frac{(a-1)\pi}{4}}}{2^{1-a}\Gamma\!\left(\frac{1-a}{2}\right)}$, so that
$P_a(z,t)=C_a\,t^{-\frac{3-a}{2}}\,z^{1-a}e^{i\frac{z^2}{4t}}$.
A direct computation gives
\[
\partial_t P_a
=
P_a\left(-\frac{3-a}{2t}-i\frac{z^2}{4t^2}\right).
\]
Next,
\[
\mathscr B_z^{(a)}\!\left(z^{1-a}e^{i\frac{z^2}{4t}}\right)
=
\left(i\frac{3-a}{2t}-\frac{z^2}{4t^2}\right)
z^{1-a}e^{i\frac{z^2}{4t}},
\]
and it follows that
\[
\mathscr B_z^{(a)}P_a
=
P_a\left(i\frac{3-a}{2t}-\frac{z^2}{4t^2}\right).
\]
Hence
\[
i\,\mathscr B_z^{(a)}P_a
=
P_a\left(-\frac{3-a}{2t}-i\frac{z^2}{4t^2}\right)
=
\partial_t P_a.
\]
Therefore,
\[
\partial_t P_a-i\mathscr B_z^{(a)}P_a=0.
\]


\end{proof}

%%%%%%%%%%%%%%%%%


%%%%%%%%%%%%%%%%%%%%%%



\section{Classical solvability for Schwartz data}\label{S:schwartz}

The explicit representation formula obtained in the previous section
suggests that the oscillatory Poisson operator should solve the extension
problem in complete analogy with the  extension theories for
fractional elliptic and parabolic operators. The purpose of this section
is to verify that this is indeed the case.

Assuming that the boundary datum belongs to the Schwartz class, we prove
that the oscillatory Poisson integral defines a smooth solution of the
singular Schr\"odinger equation in the upper half-space, satisfies the
prescribed boundary condition, and realizes the fractional Schr\"odinger
operator through the associated Dirichlet-to-Neumann map. The rapid decay
of Schwartz functions permits a completely classical treatment, avoiding
the functional-analytic issues that arise in the $L^2$ setting. Those
questions will be addressed in the next section, where the oscillatory
Poisson operator is extended to the natural graph space of the fractional
Schr\"odinger operator.

Throughout this section we assume that $u\in\mathscr S(\RNu)$,
and with $m_s$ as in \eqref{ms}, we denote by
\[
\widehat U(\xi,z,\tau)
=
m_s(z;\xi,\tau)\widehat u(\xi,\tau)
\]
the oscillatory extension introduced in \eqref{supernicehatV0} of Theorem \ref{T:itstheFT}.
The rapid decay of $\widehat u$ allows one to differentiate under the
inverse Fourier transform and thereby recover a smooth solution of the
singular Schr\"odinger equation.

\begin{theorem}[Classical solvability]
\label{T:classical}

Let $u\in\mathscr S(\RNu)$.
Then the oscillatory extension satisfies
\[
U\in C^\infty(\RNu\times(0,\infty)).
\]
Moreover:
\begin{itemize}

\item[(2)] $U$ satisfies the extension equation
\[
\partial_tU - i(\Delta_x U+\Ba U) = 0
\]
pointwise in $\RNu\times(0,\infty)$;
\item[(3)] the boundary conditions
\[
U(\cdot,z,\cdot)\ 
\underset{z\to 0^+}{\longrightarrow}\ 
u,
\]
and
\[
-
\frac{\Gamma(s)}{2^{1-2s}\Gamma(1-s)}
e^{i\pi s/2}
 z^a\partial_zU\ 
\underset{z\to 0^+}{\longrightarrow}\ 
\Le^s u
\]
hold strongly in
$L^2(\mathbb R^{d+1})$.
\end{itemize}
\end{theorem}

\begin{proof}

Since
$\widehat u\in\mathscr S(\RNu)$,
it is sufficient to show that every mixed derivative of the multiplier
$m_s(z;\xi,\tau)$ has at most polynomial growth in
$(\xi,\tau)$ on compact $z$-intervals.
Write
\[
w = z\theta(\xi,\tau)|p(\xi,\tau)|^{1/2}.
\]
Since
$\arg w\in\{0,-\pi/2\}$,
the modified Bessel function $K_s$
and all of its derivatives satisfy the standard asymptotic estimates
\[
K_\nu(w) = O(w^{-\nu}),
\qquad
w\rightarrow0,
\]
and
\[
K_\nu(w)
=
O(|w|^{-1/2}e^{-w}),
\qquad
|w|\rightarrow\infty,
\]
along the admissible rays.

Repeated application of the recurrence identity
\[
\frac{d}{dw}
\left(
w^\nu K_\nu(w)
\right) = - w^\nu K_{\nu-1}(w)
\]
shows that every derivative
$\partial_z^\beta\partial_{\xi,\tau}^\alpha m_s$
is a finite linear combination of expressions of the form
\[
P_{\alpha,\beta}
(z,p,\partial p)
K_\nu(w),
\]
where
$P_{\alpha,\beta}$
has at most polynomial growth in
$(\xi,\tau)$.

Consequently, for every compact interval
$I\Subset(0,\infty)$
there exist constants
$C_{\alpha,\beta,I}$,
$N_{\alpha,\beta}$
such that
\[
\sup_{z\in I}
\left|
\partial_z^\beta
\partial_{\xi,\tau}^\alpha
m_s(z;\xi,\tau)
\right|
\le
C_{\alpha,\beta,I}
(1+|\xi|+|\tau|)^{N_{\alpha,\beta}}.
\]
Since
$\widehat u$
is rapidly decreasing,
all mixed derivatives
\[
\partial_z^\beta
\partial_{x,t}^\gamma U
\]
may therefore be obtained by differentiating under the inverse Fourier
transform.
This proves that
\[
U
\in
C^\infty(\mathbb R^{d+1}\times(0,\infty)).
\]

Finally, the multiplier
$m_s$
satisfies the ordinary differential equation
\[
\partial_{zz}m_s
+
\frac az
\partial_zm_s
-
p(\xi,\tau)m_s
=
0,
\]

which is precisely the Fourier transform of
\[
\partial_tU
-
i(\Delta_x+B_z^{(a)}U)
=
0.
\]
Taking inverse Fourier transforms gives the extension equation
pointwise.
For the boundary conditions, the pointwise multiplier limits from
Theorem~\ref{T:itstheFT} are upgraded to strong $L^2$ convergence by dominated
convergence: since $\widehat u\in\mathscr S$, the pointwise multiplier differences
are bounded, for $0<z\le1$, by a fixed polynomial in $(\xi,\tau)$, while
$\widehat u$ decays faster than any polynomial. This proves item~(3).





\end{proof}

Theorem \ref{T:classical} shows that the oscillatory Poisson integral provides the
natural classical solution of the extension problem for Schwartz boundary
data. The next section shows that the same construction extends beyond
the Schwartz class to the natural $L^2$ graph space of the fractional
Schr\"odinger operator. Somewhat unexpectedly, the resulting theory
exhibits a qualitative dichotomy according to whether
\[
0<s\le\frac12
\qquad\text{or}\qquad
\frac12<s<1.
\]
As will become apparent, this dichotomy is not merely a technical feature
of the extension procedure, but reflects the same change in dispersive
behavior of the Bessel Schr\"odinger propagator that was first identified
in \cite{GS1,GS2}.





%%%%%%%%%%%%%%%%%%%%%%%%%%%%%%%%

\section{Extension of the oscillatory Poisson operator to $L^2$}\label{S:L2poi}

The oscillatory Poisson kernel introduced in Definition \ref{D:Pa} provides an explicit
representation of the solution to the extension problem for Schwartz data.
The purpose of this section is to show that the corresponding extension operator
extends beyond the Schwartz class to the appropriate $L^2$-based graph spaces, and that the boundary trace and the
Dirichlet-to-Neumann map remain valid in this setting.

Throughout this section we fix
$0<s<1$, 
and continue to denote by
$m_s(z;\xi,\tau)$
the multiplier \eqref{ms}.
The first step is to establish uniform estimates for the oscillatory
multiplier. Unlike the classical elliptic and parabolic extension
theories, these estimates exhibit a qualitative change at the threshold
\(s=\frac12\), reflecting the different dispersive behavior of the
underlying Bessel Schr\"odinger propagator.


\begin{lemma}\label{L:MultiplierBound}
Let $0<\nu<1$ and set $M_\nu(w)=w^\nu K_\nu(w)$ on the rays
$\arg w\in\{0,-\pi/2\}$. There exists $C_\nu>0$ such that
\[
|M_\nu(w)|\le C_\nu,\qquad 0<\nu\le\frac12,
\]
and, if $\frac12<\nu<1$,
\[
|M_\nu(w)|\le C_\nu\bigl(1+|w|^{\nu-\frac12}\bigr).
\]
Consequently, for $(\xi,\tau)\in L^+\cup L^-$,
\[
|m_s(z;\xi,\tau)|\le C_s
\Bigl(1+(z|p(\xi,\tau)|^{1/2})^{(s-\frac12)_+}\Bigr).
\]
\end{lemma}

\begin{proof}
The small-argument estimate follows from \eqref{zeroKnu}. For large argument,
\eqref{abK} gives exponential decay on the positive real ray, while on the
ray $\arg w=-\pi/2$ it gives
\[
|M_\nu(w)|\le C_\nu |w|^{\nu-\frac12}(1+O(|w|^{-1})).
\]
Combining the small- and large-argument estimates proves the assertions.
\end{proof}



The boundedness properties of the multiplier established in Lemma~\ref{L:MultiplierBound}
allow the oscillatory Poisson operator to be extended beyond the
Schwartz class. As in Lemma~\ref{L:MultiplierBound}, the resulting theory exhibits a
qualitative dichotomy according to whether \(0<s\le\frac12\) or
\(\frac12<s<1\). In the former case the extension operator is bounded on
\(L^2(\mathbb R^{d+1})\), whereas in the latter it is naturally defined
on $\Dom(\Le^s)$, see Definition \ref{D:domain}.

\begin{theorem}\label{T:L2Extension}
Let $0<s<1$ and $z>0$. For $0<s\le\frac12$ let $u\in L^2(\RNu)$, while for $\frac12<s<1$ let $u\in\Dom(\Le^s)$. We define 
\[
\widehat{U}(\xi,z,\tau)
:=
m_s(z;\xi,\tau)\ \widehat u(\xi,\tau).
\]
If we define $T_z u(x,t) := U(x,z,t)$, then:
\begin{itemize}
\item[(i)] if $0<s\le\frac12$, $T_z:L^2(\RNu)
\longrightarrow
L^2(\RNu)$
is a bounded linear operator for every $z>0$. Moreover,  $U(\cdot,z,\cdot)\ 
\underset{z\to 0^+}{\longrightarrow}\ 
u$
\ strongly in\ $L^2(\mathbb R^{d+1})$ for every
$u\in L^2(\RNu)$;
\item[(ii)] if $\frac12<s<1$,  $T_z:\Dom(\Le^s)
\longrightarrow
L^2(\RNu)$
is a bounded linear operator for every $z>0$. Moreover, $U(\cdot,z,\cdot)\ 
\underset{z\to 0^+}{\longrightarrow}\ 
u$\ strongly in $L^2(\RNu)$ for every
$u\in\Dom(\Le^s)$.
\end{itemize}
\end{theorem}

\begin{proof}
By Plancherel's theorem, we have
\[
\|T_zu\|_{L^2(\RNu)}^2 = \|\widehat{U}(\cdot,z,\cdot)\|_{L^2(\RNu)}^2
=
\int_{\R^{d+1}}
|m_s(z;\xi,\tau)|^2\,|\widehat u(\xi,\tau)|^2\,d\xi\,d\tau.
\]

We consider separately the two regimes.

\smallskip

\noindent \underline{Case (i):} $0<s\le \frac12$.
In this range, (i) in Lemma \ref{L:MultiplierBound}
 yields a uniform bound
\[
|m_s(z;\xi,\tau)|\le C_s,
\qquad z>0,\ \ (\xi,\tau)\in L^+\cup L^-.
\]
Hence
\[
\|T_zu\|_{L^2(\R^{d+1})}^2
\le
C_s^2 \int_{\R^{d+1}} |\widehat u(\xi,\tau)|^2\,d\xi\,d\tau
=
C_s^2 \|u\|_{L^2(\R^{d+1})}^2,
\]
so $T_z$ is bounded on $L^2(\RNu)$.
Moreover, \eqref{pointwise} shows that for every fixed $(\xi,\tau)\in L^+\cup L^-$ one has 
\[
m_s(z;\xi,\tau)\longrightarrow 1
\qquad\text{as }z\to0^+,
\]
and again (i) in Lemma \ref{L:MultiplierBound} provides the integrable majorant $C_s^2|\widehat u(\xi,\tau)|^2$ for the Plancherel integrand.
Therefore dominated convergence gives
\[
\|U(\cdot,z,\cdot)-u\|_{L^2(\R^{d+1})}\to0
\qquad\text{as }z\to0^+.
\]

\smallskip

\noindent\underline{Case (ii):} $\frac12<s<1$.
In this range, (ii) in Lemma \ref{L:MultiplierBound} gives
\[
|m_s(z;\xi,\tau)|
\le
C_s\bigl(1+(z|p(\xi,\tau)|^{1/2})^{\,s-\frac12}\bigr)
\le C_s\bigl(1+|p(\xi,\tau)|^{\,\frac{2s-1}{4}}\bigr),
\qquad 0<z\le1 .
\]
We restrict to $0<z\le1$, which entails no loss since only the limit $z\to0^+$
is at issue; we note in passing that as $z\to\infty$ the operator norm of $T_z$
does grow, like $z^{s-\frac12}$, so the bound below is uniform on $(0,1]$ but
not on $(0,\infty)$. Squaring the preceding estimate, for $0<z\le1$ we obtain
\[
|m_s(z;\xi,\tau)|^2
\le C_s\bigl(1+|p(\xi,\tau)|^{\,s-\frac12}\bigr)
\le C_s\bigl(1+|p(\xi,\tau)|^{\,2s}\bigr).
\]
For an arbitrary fixed $z>1$ the same argument gives boundedness with a
constant depending on $z$, which is sufficient for the first assertion of
part~(ii).
Hence, if $u\in\Dom(\Le^s)$, then
\[
\|T_zu\|_{L^2(\RNu)}^2
\le
C_s
\int_{\RNu}
\bigl(1+|p(\xi,\tau)|^{\,2s}\bigr)\,|\widehat u(\xi,\tau)|^2\,d\xi\,d\tau \le C_s\ ||u||^2_{\Dom(\Le^s)},
\]
where in the last inequality we have used the norm equivalence \eqref{equigraph}. This shows that $T_z:\Dom(\Le^s)\to L^2(\RNu)$ is bounded. We do not optimize the graph exponent here; the preceding estimate in fact permits a weaker domain assumption.
Finally, since \(m_s(z;\xi,\tau)\to1\) pointwise as \(z\to0^+\), and
\[
|m_s(z;\xi,\tau)-1|^2
\le
C_s\bigl(1+|p(\xi,\tau)|^{\,2s}\bigr),
\]
dominated convergence yields
\[
\|U(\cdot,z,\cdot)-u\|_{L^2(\R^{d+1})}\to0
\qquad\text{as }z\to0^+.
\]
This completes the proof.

\end{proof}

\begin{theorem}[Dirichlet-to-Neumann convergence in $L^2$]\label{T:DtNL2}
Let $0<s<1$. If $\frac12\le s<1$ and $u\in \Dom(\Le^s)$, or if $0<s<\frac12$ and $u\in \Dom(\Le^{\frac14+\frac s2})$, then the strong limit in $L^2(\RNu)$
\begin{equation}\label{des}
-\ \underset{z\to 0^+}{\lim}\ \frac{\Gamma(s)e^{i\pi s/2}}{2^{1-2s}\Gamma(1-s)} \ z^a \p_z U(\cdot,z,\cdot)\ =\  
\Le^s u.
\end{equation}
\end{theorem}

\begin{proof}
By Plancherel's theorem and \eqref{onL+0} in Theorem \ref{T:itstheFT}, 
\begin{align*}
& \bigl\|\frac{\Gamma(s)e^{i\pi s/2}}{2^{1-2s}\Gamma(1-s)} \ z^a \p_z U(\cdot,z,\cdot)
+\ \Le^su
\bigr\|_{L^2(\RNu)}^2
\\
& =
\int_{\RNu}
\left|
\frac{\Gamma(s)e^{i\pi s/2}}{2^{1-2s}\Gamma(1-s)} z^a\ \p_z m_s(z;\xi,\tau)
+\ (ip(\xi,\tau))^s
\right|^2
|\widehat u(\xi,\tau)|^2\,d\xi\,d\tau,
\end{align*}
where in the last equality we have used \eqref{supernicehatV0}.
Hence it suffices to show pointwise convergence of the multiplier and to
find an \(L^2\)-majorant independent of \(z\).
From the proof of Theorem \ref{T:itstheFT} we have
\[
z^a\partial_z m_s(z;\xi,\tau)
=
-\frac{2^{\,1-s}}{\Gamma(s)}\,
\theta(\xi,\tau)^{2s}\,
|p(\xi,\tau)|^s\,
M_{1-s}\!\bigl(z\theta(\xi,\tau)|p(\xi,\tau)|^{1/2}\bigr),
\]
where \(M_\nu(w)=w^\nu K_\nu(w)\). Since
\[
M_{1-s}(w)\longrightarrow 2^{-s}\Gamma(1-s)
\qquad\text{as }w\to0,
\]
we obtain the pointwise limit
\[
\frac{\Gamma(s)e^{i\pi s/2}}{2^{1-2s}\Gamma(1-s)}\ z^a\partial_z m_s(z;\xi,\tau)
\underset{z\to 0^+}{\longrightarrow}\ 
- (ip(\xi,\tau))^s.
\]
Applying Lemma~\ref{L:MultiplierBound} with $\nu=1-s$, for $\frac12\le s<1$ the function $M_{1-s}$ is bounded on the admissible rays,
and therefore the majorant $C_s|p|^s$ follows. For $0<s<\frac12$, on $L^-$ one has
$|M_{1-s}(-iy)|\le C_s(1+y^{\frac12-s})$, and hence
\[
|z^a\partial_z m_s(z;\xi,\tau)|\le C_s|p(\xi,\tau)|^s\left(1+(z|p(\xi,\tau)|^{1/2})^{\frac12-s}\right).
\]
For $u\in\Dom(\Le^{\frac14+\frac s2})$ this yields an $L^2$ majorant for $0<z\le1$.
Thus dominated convergence applies in the two stated ranges.

\end{proof}


Theorems \ref{T:L2Extension} and~\ref{T:DtNL2} extend the oscillatory
Poisson construction beyond the Schwartz class while preserving the
Dirichlet-to-Neumann correspondence. The behavior of the $L^-$ Bessel multiplier
changes at $s=\frac12$: for $s\ge\frac12$ the natural graph assumption
$u\in\Dom(\Le^s)$ suffices for the Neumann convergence, whereas for
$0<s<\frac12$ the stronger hypothesis stated in Theorem~\ref{T:DtNL2} is
required. This is the same dispersive transition identified in
\cite{GS1,GS2}. Having established the extension theory at the
natural $L^2$-level, we next turn to the potential theory associated
with the fractional Schr\"odinger operator, beginning with the
construction of its inverse fractional powers and its fundamental
solution.


%%%%%%%%%%%%%%%%%%%%
\section{Inverse fractional Schr\"odinger operators and their fundamental solutions}\label{S:inverse}


In this section we develop the potential theory associated with the fully
fractional Schr\"odinger operator $\Le^s$.
Our starting point is the construction of the inverse fractional powers of
$\mathcal L$ by means of the shifted Schr\"odinger semigroup
$\{P_\sigma^{\mathcal L}\}_{\sigma>0}$,
thus providing the dispersive analogue of the classical Riesz potentials
associated with the heat semigroup, see \cite{R, Samp, St, Sa, NyS, ST}.
We prove that the resulting operator coincides with the inverse fractional
power $\mathcal L^{-s}$, derive an explicit formula for its integral kernel,
and identify this kernel as a fundamental solution of $\mathcal L^s$.
Finally, we compute its Poisson extension and show that, up to an explicit
multiplicative constant, it coincides with the integral kernel of the
boundary operator $\Theta_a^\star$ introduced in \cite{GS2}.


To begin, we observe that, given $u\in \mathscr S(\R^{d+1})$, the standard dispersive estimate deriving from \eqref{PsS} implies
\begin{equation}\label{dispps}
\|\ps u(\cdot,t)\|_{L^\infty(\Rd)}\le C(d) |\sigma|^{-d/2} ||u(\cdot,t-\sigma)||_{L^1(\Rd)}.
\end{equation}
For every $1\le p\le \infty$ we obtain from \eqref{dispps}
\begin{equation}\label{pp}
\|\ps u\|_{L^p_t L^\infty_x}\le C(d) |\sigma|^{-d/2} ||u||_{L^p_t L^1_x},\ \ \ \sigma>0.
\end{equation}
On the other hand, the unitary identity \eqref{Suni} in Proposition \ref{P:L2} gives
\begin{equation}\label{twotwo}
\|\ps u\|_{L^2_t L^2_x} = ||u||_{L^2_t L^2_x},\ \ \ \ \ \sigma>0.
\end{equation}
Motivated by the representation \eqref{fS} of the fractional powers $\Le^s$, we now
introduce the inverse potential associated with the Schr\"odinger semigroup.

\begin{definition}\label{D:Spot}
Let $0<s<\frac d2$. For any $u\in \mathscr S(\R^{d+1})$, we define
\[
\mathscr J_{2s}(u)(x,t) = \frac{1}{\G(s)} \int_0^\infty \sigma^{s -1} \ps u(x,t) d\sigma.
\]
\end{definition}
 
Let us notice that $(x,t)\to \mathscr J_{2s}(u)(x,t)$ is measurable and the integral converges absolutely for a.e. $(x,t)\in \R^{d+1}$. This can be seen by writing
\begin{align*}
\mathscr J_{2s}(u)(x,t) & = \frac{1}{\G(s)} \int_0^1 \sigma^{s -1} \ps u(x,t) d\sigma + \frac{1}{\G(s)} \int_1^\infty \sigma^{s -1} \ps u(x,t) d\sigma
\\
& = \mathscr J^0_{2s}(u)(x,t) + \mathscr J^\infty_{2s}(u)(x,t),
\end{align*}
and observing that
\eqref{twotwo} implies 
\[
||\mathscr J^0_{2s}(u)||_{L^2(\R^{d+1})} \le \frac{1}{\G(1+s)} ||u||_{L^2(\R^{d+1})}.
\]
Furthermore, since $0<s<\frac d2$, from \eqref{dispps} we obtain for any $1\le p\le \infty$
\[
||\mathscr J^\infty_{2s}(u)||_{L^p_t L^{\infty}_x} \le C(d) ||u||_{L^p_t L^{1}_x}\ \int_1^\infty \sigma^{-\frac d2 + s -1} d\sigma = \frac{2\ C(d)}{d-2s} ||u||_{L^p_t L^{1}_x} < \infty. 
\]

The next theorem identifies the operator introduced in Definition \ref{D:Spot} with the inverse
fractional power of the Schr\"odinger operator.


Before stating the inversion theorem we make one point explicit, since it
affects the sense in which the two operators may be composed. If
$u\in\mathscr S(\R^{d+1})$, then $\Le^su\notin\mathscr S(\R^{d+1})$ in general:
being nonlocal of positive order, $\Le^s$ destroys rapid decay and leaves only
polynomial decay. The same is true of $\mathscr J_{2s}u$. Neither composition
below is therefore a composition of maps of $\mathscr S$ into itself, and for
this reason we regularise the defining integral. For $\ve>0$ and
$v\in L^2(\R^{d+1})$ we set
\begin{equation}\label{Jeps}
\mathscr J^\ve_{2s}(v)(x,t) := \frac{1}{\G(s)} \int_0^\infty \sigma^{s -1} e^{-\ve \sigma}\,P^\Le_\sigma v(x,t)\, d\sigma,
\end{equation}
an absolutely convergent $L^2$-valued integral, since
$\|P^\Le_\sigma v\|_{L^2}=\|v\|_{L^2}$ by \eqref{twotwo} and
$\int_0^\infty \sigma^{s-1}e^{-\ve\sigma}d\sigma<\infty$. On the class of
Definition \ref{D:Spot} one has $\mathscr J^\ve_{2s}(v)\to \mathscr J_{2s}(v)$
as $\ve\to0^+$.

\begin{theorem}\label{T:inverse}
Let $0<s<\min\{1,d/2\}$ and $u\in \mathscr S(\R^{d+1})$. Then
\begin{itemize}
\item[(i)] $\mathscr J^\ve_{2s}(\Le^s u)\longrightarrow u$ in $L^2(\R^{d+1})$ as
$\ve\to0^+$; we write this as $\mathscr J_{2s}\Le^su=u$;
\item[(ii)] if moreover $|p|^{-s}\hat u\in L^2(\R^{d+1})$ --- which holds
automatically when $0<s<\frac12$ --- then $\mathscr J_{2s}u\in\Dom(\Le^s)$ and
$\Le^s\mathscr J_{2s}u=u$.
\end{itemize}
In this sense $\mathscr J_{2s}=\Le^{-s}$.
\end{theorem}

\begin{proof}
Both assertions are immediate on the Fourier side. By \eqref{fSF} and
\eqref{Jeps}, for $v\in L^2(\R^{d+1})$,
\[
\widehat{\mathscr J^\ve_{2s}(v)}
= \frac{1}{\G(s)}\int_0^\infty \sigma^{s-1}e^{-(\ve+ip)\sigma}\,d\sigma\ \hat v
= (\ve + i p)^{-s}\,\hat v ,
\]
the interchange being legitimate because the $\sigma$-integral converges
absolutely in $L^2$, and the elementary identity
$\int_0^\infty\sigma^{s-1}e^{-\la\sigma}d\sigma=\G(s)\la^{-s}$ being valid for
$\Re\la>0$ with the principal branch.

For (i) we take $v=\Le^su$, so that $\hat v=(ip)^s\hat u$ by \eqref{fp2}, and
obtain
\[
\widehat{\mathscr J^\ve_{2s}(\Le^su)}
= \left(\frac{ip}{\ve+ip}\right)^{s}\hat u .
\]
Since $p$ is real and $\ve>0$ we have $|\ve+ip|^2=\ve^2+p^2\ge p^2$, whence
$\bigl|\bigl(\tfrac{ip}{\ve+ip}\bigr)^{s}\bigr|\le1$, while for every
$(\xi,\tau)$ with $p(\xi,\tau)\ne0$ the same quantity tends to $1$ as
$\ve\to0^+$. The set $\{p=0\}$ is the paraboloid
$2\pi\tau+4\pi^2|\xi|^2=0$, a smooth hypersurface and hence a null set in
$\R^{d+1}$. Dominated convergence with majorant $|\hat u|\in L^2$ therefore
gives $\widehat{\mathscr J^\ve_{2s}(\Le^su)}\to\hat u$ in $L^2$, and (i)
follows by Plancherel. Note that this argument uses only $0<s<1$.

For (ii), $(\ve+ip)^{-s}\hat u\to (ip)^{-s}\hat u$ pointwise off the paraboloid
and $|(\ve+ip)^{-s}\hat u|\le |p|^{-s}|\hat u|$, which is in $L^2$ by
hypothesis; hence $\mathscr J_{2s}u$ has Fourier transform $(ip)^{-s}\hat u$
and $|p|^s\widehat{\mathscr J_{2s}u}=\hat u\in L^2$, so that
$\mathscr J_{2s}u\in\Dom(\Le^s)$ by Definition \ref{D:domain} and
$\widehat{\Le^s\mathscr J_{2s}u}=(ip)^s(ip)^{-s}\hat u=\hat u$. The
parenthetical sufficient condition follows because $|p|$ vanishes simply and
transversally on the paraboloid, so that $|p|^{-2s}$ is locally integrable
across it precisely when $2s<1$.
\end{proof}

\begin{remark}\label{R:inverse-real}
The identity in (i) admits the following direct real-variable derivation, which
we record because it exhibits the mechanism behind the cancellation. It is
formal to the extent that it manipulates the unregularised integrals.
Note that, since $\Le u\in \mathscr S(\R^{d+1})$,  an easy integration by parts shows that  Definition \eqref{fS} can be written in the following alternative way
\begin{equation}\label{fSalt}
\Le^s u(x,t) = - \frac{1}{\G(1-s)} \int_0^\infty \frac{1}{\sigma^s} \frac{d}{d\sigma} \ps u(x,t) d\sigma.
\end{equation}
In view of \eqref{fSalt} and of Definition \ref{D:Spot}, we find
\begin{align*}
& \mathscr J_{2s}(\Le^s u)(x,t) = - \frac{1}{\G(s)\G(1-s)} \int_0^\infty \int_0^\infty \sigma^{s-1} \frac{1}{\tau^{s}} \ps(\frac{d}{d\tau} P^{\Le}_\tau u)(x,t) d\tau d\sigma.
\end{align*}
To justify the calculation without assuming absolute convergence of the full double integral, we first truncate both semigroup parameters to a compact rectangle $\varepsilon<\sigma,\tau<R$. On this rectangle Fubini's theorem and the change of variables below are legitimate. We perform the calculation there and then let $\varepsilon\downarrow0$ and $R\uparrow\infty$. The boundary terms converge by the strong continuity of $P_\alpha^{\Le}$ at $\alpha=0$ and the dispersive decay \eqref{pp} as $\alpha\to\infty$; the beta factor is integrable because $0<s<1$. Thus the limiting identity is the one displayed below.


By the semigroup property of $P^\Le_\tau$ we have $P^\Le_\sigma \circ P^\Le_\tau = P^\Le_{\sigma+\tau}$, and since $P^\Le_\sigma$ and $\frac{d}{d\tau}$ commute, we have
\begin{align*}
& P^\Le_\sigma(\frac{d}{d\tau} P^\Le_\tau u)(x,t) = \frac{d}{d\tau} P^\Le_{\sigma+\tau} u(x,t).
\end{align*}
Substituting in the above formula we obtain
\begin{align}\label{gnam}
& \mathscr J_{2s}(\Le^s u)(x,t) = - \frac{1}{\G(s)\G(1-s)} \int_0^\infty \int_0^\infty \left(\frac \sigma{\tau}\right)^{s} \frac{d}{d\tau} P^\Le_{\sigma+\tau} u(x) d\tau \frac{d\sigma}\sigma.
\end{align}
In the integral in the right-hand side we now make the change of variable $\Phi:(0,\infty)\times(0,\infty)\to (0,\infty)\times(0,\infty)$
\begin{equation}\label{uv}
\alpha = \sigma + \tau,\ \ \ \ \ \ \beta = \frac \sigma{\tau},
\end{equation}
whose inverse $\Psi(\alpha,\beta) = (\sigma,\tau)$ is given by
\begin{equation}\label{ttau}
\sigma = \frac{\alpha\beta}{1+\beta},\ \ \ \ \ \ \tau = \frac{\alpha}{1+\beta}.
\end{equation}
The Jacobian of the diffeomorphism $\Psi$ is 
\[
d\sigma d\tau = \frac{\alpha}{(1+\beta)^2} d\alpha d\beta.
\]  
Keeping in mind that the chain rule gives 
\[
\frac{d}{d\tau} P^\Le_{\sigma+\tau} u(x,t) = \left(\frac{\p \alpha}{\p \tau} \frac{d}{d\alpha} + \frac{\p \beta}{\p\tau} \frac{d}{d\beta}\right) P^\Le_\alpha u(x,t) =  \frac{d}{d\alpha} P^\Le_\alpha u(x,t),
\]
we thus find
\begin{align*}
& \mathscr J_{2s}(\Le^s u)(x,t) = - \frac{1}{\G(s)\G(1-s)} \int_0^\infty \int_0^\infty \beta^{s} \frac{d}{d\alpha} P^\Le_{\alpha} u(x,t) \frac{\alpha}{(1+\beta)^2} \frac{1+\beta}{\alpha\beta} d\alpha d\beta
\\
& = - \frac{1}{\G(s)\G(1-s)} \int_0^\infty  \frac{\beta^{s - 1}}{1+\beta} d\beta\ \int_0^\infty \frac{d}{d\alpha} P^\Le_{\alpha} u(x,t) d\alpha = u(x,t),
\end{align*}
where we have appealed to \eqref{pp} to conclude that
\[
|P^\Le_{\alpha} u(x,t)|\le ||P^\Le_{\alpha} u||_{L^\infty(\R^{d+1})} \le C(d) |\alpha|^{-d/2} ||u||_{L_t^\infty L^1_x}\ \underset{\alpha\to \infty}{\longrightarrow}\ 0,
\]
and therefore
\[
\int_0^\infty \frac{d}{d\alpha} P^\Le_{\alpha} u(x,t) d\alpha = - u(x,t),
\]
whereas we have used the following well-known representation of the Beta function
\[
B(x,y) = \int_0^\infty \frac{t^{y-1}}{(1+t)^{x+y}} dt,\ \ \ \ \ \ x,y>0, 
\]
to infer that
\[
\frac{1}{\G(s)\G(1-s)} \int_0^\infty  \frac{\beta^{s - 1}}{1+\beta} d\beta = \frac{B(1-s,s)}{\G(s)\G(1-s)} = 1.
\]


\end{remark}

\subsection{Fundamental solution of the fractional Schr\"odinger operator $\Le^s$}\label{SS:fs}

Given $u\in \mathscr S(\R^{d+1})$, we have by \eqref{PsS} 
\begin{equation}\label{pss}
P^\Le_\sigma u(x,t)  = \int_{\Rd} S(x,y,\sigma) u(y,t-\sigma) dy = S(\sigma)(u(\cdot,t-\sigma))(x),\ \ \ \ \ \ \sigma>0.
\end{equation}
Substituting \eqref{pss} in Definition \ref{D:Spot}, we have for $0<s<\frac d2$ 
\[
\mathscr J_{2s}(u)(x,t) = \frac{1}{\G(s)} \int_0^\infty \sigma^{s -1} P^\Le_\sigma u(x,t) d\sigma.
\]
We thus find
\begin{align}\label{J2s}
\mathscr J_{2s}(u)(x,t) & = \frac{1}{\G(s)}  \int_\R \int_{\Rd} \mathbf 1_{(0,\infty)}(\sigma) \sigma^{s -1}  S(x,y,\sigma) u(y,t-\sigma) dy d\sigma  
\\
& = \frac{1}{\G(s)} \int_\R \int_{\Rd}  \frac{\mathbf 1_{(0,\infty)}(t-\sigma)}{(t-\sigma)^{1-s}}  S(x,y,t-\sigma) u(y,\sigma)  dy d\sigma. 
\notag
\end{align}
From Theorem \ref{T:inverse} it follows that the kernel 
of the operator $\mathscr J_{2s}$, 
\begin{equation}\label{Es}
E_s(x,t;y,\tau) =  \frac{1}{\G(s)} \frac{\mathbf 1_{(0,\infty)}(t-\tau)}{(t-\tau)^{1-s}}  S(x,y,t-\tau),\qquad 0<s<\min\{1,d/2\},
\end{equation}
provides a fundamental solution of the nonlocal operator $\Le^s$ with pole at $(y,\tau)$. This is equivalent to having
\begin{equation}\label{E:fundamental-inversion}
\frac{1}{\G(s)} \int_\R\int_{\Rd}
\frac{\mathbf 1_{(0,\infty)}(t-\tau)}{(t-\tau)^{1-s}}
S(x,y,t-\tau)\,\Le^s u(y,\tau)\,dy\,d\tau=u(x,t).
\end{equation}
Here the integration is over the pole variables $(y,\tau)$. This is the
causal inversion identity associated with $\mathscr J_{2s}$; no formal
self-adjointness of $\Le^s$ is being used.
It follows from \eqref{Es} that, for a suitable $F(x,t)$ defined in $\RNu$, the solution of the nonlocal dispersive equation
\begin{equation}\label{Lsu=F}
\Le^s u = F,
\end{equation}
is given for any $(x,t)\in \R^{d+1}$ by the convolution of $F$ with $\mathscr E_s$, i.e.
\begin{equation}\label{potential}
u(x,t) =  \frac{1}{\G(s)} \int_\R \int_{\Rd}  \frac{\mathbf 1_{(0,\infty)}(t-\tau)}{(t-\tau)^{1-s}}  S(x,y,t-\tau) F(y,\tau) dy d\tau.
\end{equation}
If, in particular, 
\begin{equation}\label{F}
F(x,t) = \begin{cases}
\Phi(x,t),\ \ \ \ \ \ \ \ \ t> 0,
\\
\Le^s v_0(x,t),\ \ \ \ \ t\le 0,
\end{cases}
\end{equation}
then formula \eqref{potential} gives the following result.

\begin{proposition}\label{P:fsLs}
Let $\Phi, v_0\in \mathscr S(\RNu)$. Then the solution of the equation \eqref{Lsu=F}, with $F$ given by \eqref{F}, is represented by the oscillatory integral
\begin{align}\label{pot2}
u(x,t) & = \frac{1}{\G(s)} \int_0^t \frac{1}{\sigma^{1-s}} S(\sigma)(\Phi(\cdot,t-\sigma))(x) d\sigma
+ \frac{1}{\G(s)} \int_{t}^\infty \frac{1}{\sigma^{1-s}} S(\sigma)(\Le^s v_0(\cdot,t-\sigma))(x)  d\sigma.
\end{align}
\end{proposition}

\subsection{Poisson lifting of the fundamental solution}\label{P:liftingEs}

We next compute the Poisson extension of the fundamental solution $\mathscr E_s$ in \eqref{Es}. According to \eqref{U} such extension is provided by
\begin{equation}\label{EEs}
\mathscr  E_s(x,z,t) := \int_0^\infty \int_{\Rd} \mathbb P_a(X,Y_0,\sigma) E_s(y,t-\sigma) dy d\sigma,\ \ \ \ a=1-2s,
\end{equation}
where for brevity we have written 
\[
E_s(y,t-\sigma) = E_s(y,t;0,\sigma) = \frac{1}{\G(s)} \frac{\mathbf 1_{(0,\infty)}(t-\sigma)}{(t-\sigma)^{1-s}}  S(y,0,t-\sigma). 
\]
The kernel $E_s$ is an oscillatory distribution rather than a Schwartz function. The identity below is therefore understood in the distributional sense (equivalently, after Fourier transformation in the tangential variables). 

\begin{proposition}\label{P:Es}
For every $0<s<\min\{1,d/2\}$, $t>0$ and $z>0$, we have
\[
\mathscr  E_s(x,z,t) = \frac{1}{\G(s)} \frac{e^{i \frac{z^2}{4t}}}{t^{1-s}} S(x,0,t) = e^{i \frac{z^2}{4t}} \ E_s(x,t).
\]
\end{proposition}

\begin{proof}
Keeping \eqref{Pa} in mind, we have
\[
\mathbb P_a(X,Y_0,\sigma)\ E_s(y,t-\sigma) = \frac{1}{\G(s)} P_a(z,\sigma) \frac{\mathbf 1_{(0,\infty)}(t-\sigma)}{(t-\sigma)^{1-s}}  S(x,y,\sigma)  S(y,0,t-\sigma),
\]
and we thus have from \eqref{EEs}
\begin{align*}
\mathscr  E_s(x,z,t) & = \frac{1}{\G(s)}  \int_0^\infty P_a(z,\sigma) \frac{\mathbf 1_{(0,\infty)}(t-\sigma)}{(t-\sigma)^{1-s}} \left(\int_{\Rd}  S(x,y,\sigma)  S(y,0,t-\sigma) dy\right) d\sigma
\\
& = \frac{1}{\G(s)}  \int_0^\infty P_a(z,\sigma) \frac{\mathbf 1_{(0,\infty)}(t-\sigma)}{(t-\sigma)^{1-s}} d\sigma\ S(x,0,t),
\end{align*}
where we have used the Chapman-Kolmogorov equation
\[
\int_{\Rd}  S(x,y,\sigma)  S(y,0,t-\sigma) dy = S(x,0,t).
\]
We now invoke Lemma \ref{L:abepoi} to conclude that 
\begin{equation}\label{abepoi-repeat}
\frac{1}{\G(s)} \int_0^\infty \frac{\mathbf 1_{(0,\infty)}(t-\sigma)}{(t-\sigma)^{1-s}} P_a(z,\sigma) d\sigma = \frac{1}{\G(s)} \frac{e^{i \frac{z^2}{4t}}}{t^{1-s}}.
\end{equation}
From the latter three identities the conclusion follows.

\end{proof}

Proposition \ref{P:Es} shows that the Poisson lifting of the fundamental solution is obtained simply by multiplication by the oscillatory phase $e^{i \frac{z^2}{4t}}$. 

This remarkable simplification will play a crucial role in identifying
the kernel of the boundary operator $\Theta_a^\star$ introduced in
\cite{GS2}. We recall that in \cite{GS2} we established a generalized
Duhamel representation for solutions to the problem:
\begin{equation}\label{nozeroin}
\begin{cases}
\partial_tU-i(\Delta_x U+\Ba U)=F(X,t),\\
\displaystyle \lim_{z\to0^+} z^a\partial_zU(X,t)=\Phi(x,t),\\
U(X,0)=U_0(X),
\end{cases}
\end{equation}
where $\Ba$ is given by \eqref{Ba} and $-1<a<1$. Specifically, we have proved that solutions can be represented as
\[
U= \mathbb T_a^\star(U_0)+\mathbb D_a(F)+\Theta_a^\star(\Phi),
\]
where the operators $\mathbb T_a^\star$ and $\mathbb D_a$ describe the bulk evolution, while $\Theta_a^\star$ encodes the boundary interaction. Since the bulk propagator plays an essential role in Section \ref{S:memory} below, we recall its definition explicitly. With $\mathbb S_a$ the kernel \eqref{grandeSaS0} and $d\omega_a(Y) = \zeta^a dy\, d\zeta$, one has
\begin{equation}\label{E:Tstar}
\mathbb T_a^\star(U_0)(X,t) := \int_{\RN} \mathbb S_a(X,Y,t)\, U_0(Y)\, d\omega_a(Y),
\end{equation}
\begin{equation}\label{E:Da}
\mathbb D_a(F)(X,t) := \int_0^t\!\!\int_{\RN} \mathbb S_a(X,Y,t-\tau)\, F(Y,\tau)\, d\omega_a(Y)\, d\tau .
\end{equation}

By the tensorization \eqref{Sa} of the kernel, $\mathbb T^\star_a$ is the unitary group
\begin{equation}\label{E:Tstar-group}
\mathbb T^\star_a(\cdot)(\cdot,\cdot,t) = e^{i t\left(\Delta_x + \BaN\right)}
\qquad\text{on } L^2_a(\RN),
\end{equation}
where $\BaN$ denotes the following realization of the Bessel operator. Both
end-points of $\R^+$ are singular for $\Ba$; by Frobenius, the two solutions of
$\Ba u=0$ are $u_1\equiv1$ and $u_2(z)=z^{1-a}$, and both belong to
$L^2((0,\ve),z^adz)$ when $-1<a<3$. Thus $z=0$ is in the limit-circle case in
the whole range $-1<a<1$ relevant here, so that $\Ba$ is \emph{not} essentially
self-adjoint and a boundary condition must be imposed; the point at infinity is
always limit-point. Imposing the homogeneous Neumann condition
\begin{equation}\label{E:neumann}
\underset{z\to0^+}{\lim}\ z^a\p_z u = 0,
\end{equation}
which eliminates $u_2$, singles out a self-adjoint realization, and it is this
realization that we denote by $\BaN$; see \cite[(1.9)--(1.11)]{GS1}, where the
Weyl analysis is carried out, and where it is also shown that $e^{it\BaN}$ is
the unitary group whose kernel is $S_a$ of \eqref{SaS}. Identity
\eqref{Sazeron} is precisely the statement that this group respects
\eqref{E:neumann}. We shall use the diagonalization of $\BaN$ furnished by the
modified Hankel transform of order $\nu=\frac{a-1}2$ on $(\R^+,z^adz)$,
\begin{equation}\label{E:hankel}
\mathcal H_{\nu}(\psi)(\lambda)
= \int_0^\infty (\lambda z)^{-\nu}J_{\nu}(\lambda z)\,\psi(z)\,z^a\,dz ,
\end{equation}
see \cite[(2.3)]{GS1}. By \cite[Theorem~2.1]{GS1} the transform $\mathcal
H_\nu$ is an involution, $\mathcal H_\nu\circ\mathcal H_\nu=\operatorname{Id}$,
and by \cite[(2.7)]{GS1} it satisfies the Plancherel identity on
$L^2(\R^+,z^adz)$; it is therefore a surjective isometry of that space onto
itself. Moreover, by \cite[Theorem~2.2]{GS1},
\begin{equation}\label{E:hankeldiag}
\mathcal H_{\nu}\bigl(\BaN\psi\bigr)(\lambda)=-\lambda^2\,\mathcal H_{\nu}(\psi)(\lambda)
\end{equation}
for every $\psi$ in the space $\mathscr S^+$ introduced in \cite[Section~1]{GS1},
namely the space of the $\vf\in C^\infty(\R^+)$ for which all the seminorms
\[
\gamma_{m,k}(\vf)=\sup_{z>0}\bigl|z^m(\tfrac1z\p_z)^k\vf(z)\bigr|
\]
are finite,
equivalently the restrictions to $(0,\infty)$ of the even Schwartz functions on
$\R$. As observed in \cite{GS1}, the finiteness of $\gamma_{0,1}$ alone forces
$|\p_z\vf(z)|\le\gamma_{0,1}(\vf)\,z$, so that every element of $\mathscr S^+$
automatically satisfies \eqref{E:neumann} when $a>-1$; in particular
$\mathscr S^+\subset\Dom(\BaN)$, and $\mathcal H_\nu$ diagonalizes precisely
the realization selected by that condition. Writing
\[
\widehat{v}(\xi) = \int_{\Rd} e^{-2\pi i \sa \xi,x\da} v(x)\, dx
\]
for the Fourier transform in the tangential variable alone, it follows from \eqref{E:hankeldiag} and from the classical propagator \eqref{Scla} that, on the Fourier--Hankel side, $\mathbb T^\star_a$ acts as multiplication by the unimodular symbol
\begin{equation}\label{E:multiplier}
e^{-i t\left(4\pi^2|\xi|^2+\lambda^2\right)} .
\end{equation}

Our next objective is to identify the kernel of the boundary operator $\Theta_a^\star$ in terms of the Poisson lifting of the fundamental solution constructed above.

With the Neumann convention of \cite{GS2}, see \cite[(2.5)]{GS2} and the verification in \cite[Remark~2.1]{GS2}, the boundary Duhamel operator carries the factor $-i$. We therefore use
\begin{equation}\label{Theta}
\Theta_a^\star(\Phi)(x,z,t) = -i \int_0^{t} S_a(z,0,\sigma) S(\sigma)(\Phi(\cdot,t-\sigma))(x)  d\sigma.
\end{equation}

\begin{theorem}\label{T:id}
Let $\Theta^\star_a(X,Y_0,\sigma)$ denote the integral kernel of the boundary operator $\Theta_a^\star$ in \eqref{Theta}. Then 
\[
\Theta^\star_a(X,Y_0,t) = -\frac{2^{2s-1}e^{i\frac{\pi s}2}\G(s)}{\G(1-s)}\ \mathscr E_s(X,Y_0,t),
\]
where $\mathscr E_s(X,Y_0,t)$ is the Poisson lifting of the fundamental solution $ E_s(x,y,t)$ of $\Le^s$.
\end{theorem}

\begin{proof}

Moreover, from \cite[(3.9)]{GS2}, for $a=1-2s$ we have:
\begin{equation}\label{SaS0}
S_a(z,0,\sigma)
= \frac{2^{2s-1}}{\G(1-s)}
\frac{e^{-i \frac{(1-s)\pi}{2}}}{\sigma^{1-s}}
\,e^{i \frac{z^2}{4\sigma}}.
\end{equation}


Substituting \eqref{SaS0} in \eqref{Theta}, we obtain
\begin{equation}\label{Theta2}
\Theta_a^\star(\Phi)(x,z,t) = -\frac{2^{2s-1}e^{i\frac{\pi s}2}}{\G(1-s)} \int_0^{t} \frac{e^{i\frac{z^2}{4\sigma}}}{\sigma^{1-s}} S(\sigma)(\Phi(\cdot,t-\sigma))(x)  d\sigma.
\end{equation}

The sign is also consistent with the normalization used later in the nonlinear problem. Indeed, with
\[
c_s=-\frac{2^{1-2s}\G(1-s)}{\G(s)e^{i\pi s/2}},
\]
one has
\[
c_s\left(-\frac{2^{2s-1}e^{i\pi s/2}\G(s)}{\G(1-s)}\right)=1.
\]

Comparing \eqref{Theta2} with Proposition \ref{P:Es} and with the first term in the representation  \eqref{pot2}, we conclude that the kernel of the operator $\Theta_a^\star$ is a multiple of the Poisson extension of the fundamental solution of $\Le^s$. 

\end{proof}


%%%%%%%%%%%%%%%%%%%%%%%%%%%%%%%%%%%%%%%%%%%%%%%%%

\section{The natural history space}\label{S:history}

As we have mentioned in Section \ref{S:des}, the oscillatory Poisson lifting \eqref{lift} provides the initial datum for the singular Schr\"odinger equation in the upper half-space. Consequently, in order to apply the nonlinear well-posedness theory developed in our companion paper \cite{GS2}, it is necessary to determine which histories $u_0$ produce liftings $U_0$ belonging to the weighted energy space $L^2_a(\RN)$, see Section \ref{S:not}. We call the class of such $u_0$ the \emph{natural history space}.

The next question is whether this class can be characterized directly in terms of the prescribed past evolution.
We shall prove that the weighted
bulk energy of the Poisson lifting can be expressed by an explicit quadratic
form acting only on the history itself. This identity suggests the
introduction of an intrinsic Hilbert space of Schr\"odinger histories.
We begin with a basic preparatory result which shows that the class $\Sm$, defined in Section \ref{S:not}, is contained in the natural history space.

\begin{lemma}\label{L:U0L2}
Let $0<s<1$ and set $a=1-2s\in(-1,1)$. Assume that
$u_0\in \Sm$ and define
\begin{equation}\label{lift2}
U_0(X) := \int_0^\infty\int_{\R^d}
\mathbb P_a(X,Y_0,\sigma)\,u_0(y,-\sigma)\,dy\,d\sigma.
\end{equation}
Then $U_0\in L_a^2(\R^{d+1}_+)$.
\end{lemma}

\begin{proof}
We want to show that
\begin{equation}\label{goal}
\int_0^\infty\int_{\Rd}|U_0(x,z)|^2 z^a\,dx\,dz<\infty.
\end{equation}
 Let $v_0\in \mathscr S(\RNu)$ be such that $v_0|_{\{t\le 0\}} = u_0$. Write
\[
G(\sigma)(x):=S(\sigma)v_0(\cdot,-\sigma)(x),
\qquad \sigma>0,
\]
so that, by the tensorization property of the Poisson kernel \eqref{Pa},
\[
U_0(X)
=
c_a\,z^{2s}\int_0^\infty \sigma^{-s}
e^{\,i\frac{z^2}{4\sigma}}\,G(\sigma)(x)\frac{d\sigma}{\sigma},
\]
for a constant $c_a\neq 0$ depending only on $a$.
Since $v_0\in\mathscr S(\RNu)$ and $S(\sigma)$ is unitary on $L^2(\Rd)$, we have $\|G(\sigma)\|_{L^2_x}=\|v_0(\cdot,-\sigma)\|_{L^2_x}$ and $\|G'(\sigma)\|_{L^2_x}=\|(\Le v_0)(\cdot,-\sigma)\|_{L^2_x}$, both rapidly decreasing in $\sigma$, and therefore
$\sigma\mapsto G(\sigma)\in L^2(\Rd)$ is smooth and rapidly decaying together
with all its derivatives. In particular, for every $N\ge 0$ there exists
$C_N>0$ such that
\[
\|G(\sigma)\|_{L^2_x}+\|G'(\sigma)\|_{L^2_x}
\le C_N(1+\sigma)^{-N}
\qquad\forall\,\sigma>0.
\]
Indeed,
\begin{equation}\label{Gprime}
G'(\sigma)
=
\partial_\sigma\bigl(S(\sigma)v_0(\cdot,-\sigma)\bigr)
=
i\Delta_xS(\sigma)v_0(\cdot,-\sigma)
-
S(\sigma)\partial_t v_0(\cdot,-\sigma)
=
-\,S(\sigma)\bigl(\Le v_0\bigr)(\cdot,-\sigma),
\end{equation}
and the claim follows from the unitarity of
$S(\sigma)$ on $L^2(\R^d)$ and the Schwartz decay of $v_0$.
To reach the desired conclusion, we estimate $\|U_0(\cdot,z)\|_{L^2_x}$ separately for $0<z\le 1$ and $z\ge 1$.

\medskip
\noindent \underline{Step 1}: $0<z\le 1$.
Making the change of variables
$u=\frac{z^2}{4\sigma}$,
we obtain
\begin{equation}\label{split}
U_0(x,z)
=
c_a\,z^{2s} \int_0^\infty \left(\frac{4u}{z^2}\right)^s e^{iu}\  
G\left(\frac{z^2}{4u}\right)(x) \frac{du}{u} = \tilde c_a \int_0^\infty u^{s-1} e^{iu}\ G\left(\frac{z^2}{4u}\right)(x) du,
\end{equation}
for another harmless constant $\widetilde c_a\neq 0$ depending only on $a$.
Set
\[
H_z(u)(x):=G\!\left(\frac{z^2}{4u}\right)(x).
\]
Then
\[
\|H_z(u)\|_{L^2_x}\le \sup_{\sigma>0}\|G(\sigma)\|_{L^2_x},
\]
and, for $0<z\le 1$, we have from \eqref{Gprime}
\begin{equation}\label{Hprime}
\|H_z'(u)\|_{L^2_x}
=
\frac{z^2}{4u^2}\,
\left\|G'\!\left(\frac{z^2}{4u}\right)\right\|_{L^2_x}
\le \frac{C}{u^2},
\end{equation}
with $C$ independent of $z\in(0,1]$.
We now split the integral in \eqref{split} into $(0,1)$ and $(1,\infty)$.
On $(0,1)$ we simply use the bound on $H_z$ and the fact that $s\in(0,1)$:
\[
\left\|\int_0^1 u^{s-1}e^{iu}H_z(u)\,du\right\|_{L^2_x}
\le
\sup_{\sigma>0}\|G(\sigma)\|_{L^2_x}\int_0^1 u^{s-1}\,du
<\infty.
\]
On the interval $(1,R)$, where $R>1$, we integrate by parts using
$\partial_u e^{iu}=ie^{iu}$:
\[
\begin{aligned}
\int_1^R u^{s-1}e^{iu}H_z(u)\,du
&=
\left[\frac{1}{i}u^{s-1}e^{iu}H_z(u)\right]_1^R  \\
&\qquad
-\frac1i
\int_1^R
e^{iu}
\Bigl((s-1)u^{s-2}H_z(u)
+
u^{s-1}H_z'(u)
\Bigr)\,du .
\end{aligned}
\]
Since $0<s<1$ and
\[
\sup_{u>0}\|H_z(u)\|_{L^2_x}
\le
\sup_{\sigma>0}\|G(\sigma)\|_{L^2_x},
\]
we have
\[
R^{\,s-1}\|H_z(R)\|_{L^2_x}
\longrightarrow0
\qquad\text{as }R\to\infty.
\]
Hence the boundary term at infinity vanishes in $L^2_x$, and letting
$R\to\infty$ yields
\[
\begin{aligned}
\int_1^\infty u^{s-1}e^{iu}H_z(u)\,du
&=
-\frac{e^{i}}{i}H_z(1)
\\
&\qquad
-\frac1i
\int_1^\infty
e^{iu}
\Bigl((s-1)u^{s-2}H_z(u)
+
u^{s-1}H_z'(u)
\Bigr)\,du .
\end{aligned}
\]
The first term $H_z(1)(x) = G\!\left(\frac{z^2}{4}\right)(x)$ is bounded in $L^2_x$, whereas the second one is bounded
because
\[
u^{s-2}\in L^1(1,\infty),
\]
and from \eqref{Hprime} we have
\[
u^{s-1}\|H_z'(u)\|_{L^2_x}
\le
Cu^{s-3},
\]
with
\[
u^{s-3}\in L^1(1,\infty),
\]
since $0<s<1$.
Therefore,
\[
\sup_{0<z\le 1}\|U_0(\cdot,z)\|_{L^2_x}\le C.
\]
It follows that
\[
\int_0^1 z^a\|U_0(\cdot,z)\|_{L^2_x}^2\,dz
\le
C\int_0^1 z^a\,dz<\infty,
\]
since $a>-1$.

\medskip
\noindent \underline{Step 2}: $z\ge 1$.
Starting from the representation
\[
U_0(x,z)
=
c_a\,z^{2s}\int_0^\infty \sigma^{-1-s}
e^{\,i\frac{z^2}{4\sigma}}\,G(\sigma)(x)\,d\sigma,
\]
we integrate by parts in $\sigma$, using
\[
\partial_\sigma e^{\,i\frac{z^2}{4\sigma}}
=
-\frac{i z^2}{4\sigma^2}e^{\,i\frac{z^2}{4\sigma}}.
\]
Since $\sigma^{1-s}G(\sigma)\to 0$ as $\sigma\to 0^+$ and as $\sigma\to\infty$
(the first limit because $s<1$ and $G$ is bounded near $0$, the second by the
rapid decay of $G$), we obtain
\[
U_0(x,z)
=
\frac{4c_a}{iz^2}\,z^{2s}
\int_0^\infty
\partial_\sigma\!\bigl(\sigma^{1-s}G(\sigma)\bigr)(x)\,
e^{\,i\frac{z^2}{4\sigma}}\,d\sigma.
\]
Hence
\[
\|U_0(\cdot,z)\|_{L^2_x}
\le
C z^{2s-2}
\int_0^\infty
\left\|\partial_\sigma\!\bigl(\sigma^{1-s}G(\sigma)\bigr)\right\|_{L^2_x}\,d\sigma.
\]
Now
\[
\partial_\sigma\!\bigl(\sigma^{1-s}G(\sigma)\bigr)
=
(1-s)\sigma^{-s}G(\sigma)+\sigma^{1-s}G'(\sigma),
\]
and the right-hand side is integrable in $\sigma$ with values in $L^2_x$
because $s\in(0,1)$ and $G,G'$ are rapidly decaying. Therefore
\[
\|U_0(\cdot,z)\|_{L^2_x}\le C z^{2s-2},
\qquad z\ge 1.
\]
Consequently,
\[
\int_1^\infty z^a\|U_0(\cdot,z)\|_{L^2_x}^2\,dz
\le
C\int_1^\infty z^{a+4s-4}\,dz
=
C\int_1^\infty z^{2s-3}\,dz<\infty,
\]
since $0<s<1$.
Combining the estimates on $(0,1)$ and $(1,\infty)$, we conclude that \eqref{goal} does hold. This proves the lemma.

\end{proof}

The direct computation of the history energy is most conveniently organized
in the frequency variable. We therefore do not use the Gaussian-regularized
pairing of two Poisson kernels in the proof of the history-energy identity.

Throughout what follows we write, for $u_0\in\Sm$ and $\xi\in\Rd$,
\begin{equation}\label{E:Gxi}
g_\xi(\sigma):=e^{-4\pi^2i\sigma|\xi|^2}\,\widehat{u_0}(\xi,-\sigma)\,
\mathbf 1_{(0,\infty)}(\sigma),
\qquad
\mathscr G_\xi(\omega):=\int_{\R}e^{-i\omega\sigma}g_\xi(\sigma)\,d\sigma,
\end{equation}
where $\widehat{u_0}(\xi,t)=\mathscr F_{x\to\xi}\bigl(u_0(\cdot,t)\bigr)(\xi)$.
Thus $g_\xi$ is the history transported by the free Schr\"odinger flow. For
$u_0\in\Sm$, each $g_\xi$ is supported in $[0,\infty)$, has a finite right
limit at $0$, and is rapidly decreasing as $\sigma\to\infty$.

\begin{lemma}[Mellin reduction of the Poisson energy]\label{L:MellinReduction}
Let $a=1-2s$, $0<s<1$, and let $u_0\in\Sm$. Then, for a.e. $\xi$,
\begin{equation}\label{E:MellinReduction}
\int_0^\infty|\widehat U_0(\xi,z)|^2z^a\,dz
=
\frac{2^{1-2s}}{\G(s)^2}
\int_0^\infty \nu^s|\mathfrak G_\xi(\nu)|^2\,d\nu,
\end{equation}
where
\begin{equation}\label{E:mathfrakG}
\mathfrak G_\xi(\nu):=
\int_0^\infty e^{i\nu/\sigma}\sigma^{-1-s}g_\xi(\sigma)\,d\sigma.
\end{equation}
\end{lemma}

\begin{proof}
Taking the Fourier transform in $x$ in the Poisson lifting gives
\[
\widehat U_0(\xi,z)
=c_a z^{2s}\int_0^\infty \sigma^{-1-s}
 e^{iz^2/(4\sigma)}g_\xi(\sigma)\,d\sigma
=c_a z^{2s}\mathfrak G_\xi\!\left(\frac{z^2}{4}\right),
\]
where $|c_a|=(4^s\G(s))^{-1}$. Hence, with $\nu=z^2/4$,
\[
\begin{aligned}
\int_0^\infty|\widehat U_0(\xi,z)|^2z^a\,dz
&=|c_a|^2\int_0^\infty z^{1+2s}
\left|\mathfrak G_\xi\!\left(\frac{z^2}{4}\right)\right|^2dz\\
&=\frac{2^{1-2s}}{\G(s)^2}
\int_0^\infty\nu^s|\mathfrak G_\xi(\nu)|^2\,d\nu.
\end{aligned}
\]
\end{proof}

For a function $h$ on $(0,\infty)$ define
\[
(Wh)(\alpha):=\alpha^{s-1}h(1/\alpha),\qquad \alpha>0,
\]
so that $W^2=\mathrm{Id}$. We also write
\[
E_\pm[h]:=\int_0^\infty\omega^s|\widehat h(\pm\omega)|^2\,d\omega,
\qquad
\widehat h(\omega):=\int_0^\infty e^{-i\omega\sigma}h(\sigma)\,d\sigma,
\]
whenever these quantities are finite, with the Fourier transforms understood as
Abel boundary values when necessary.

\begin{lemma}[Mellin transform under inversion]\label{L:MellinW}
For every $h$ for which the Mellin transforms below exist,
\[
\mathscr M(Wh)(\lambda)=\mathscr Mh(1-s-\lambda).
\]
Moreover,
\[
\mathfrak G_\xi(\nu)=\widehat{Wg_\xi}(-\nu).
\]
\end{lemma}

\begin{proof}
The first identity follows from $\sigma=1/\alpha$:
\[
\mathscr M(Wh)(\lambda)
=\int_0^\infty \alpha^{s-1}h(1/\alpha)\alpha^{\lambda-1}\,d\alpha
=\int_0^\infty h(\sigma)\sigma^{1-s-\lambda}\frac{d\sigma}{\sigma}.
\]
The second follows from the same substitution:
\[
\mathfrak G_\xi(\nu)
=\int_0^\infty e^{i\nu\alpha}(Wg_\xi)(\alpha)\,d\alpha
=\widehat{Wg_\xi}(-\nu).
\]
\end{proof}

\begin{lemma}[The inversion identity for the one-sided energy]\label{L:MellinSwap}
Let $0<s<1$, and let $h$ be smooth on $(0,\infty)$, have a finite right limit at $0$, and be rapidly decreasing at infinity together with all its derivatives. Then
\begin{equation}\label{E:MellinSwap}
E_-[Wh]=E_+[h].
\end{equation}
\end{lemma}

\begin{proof}
We first justify the Mellin identity by an explicit Abel regularization. For
$0<\Re\mu<1$ and $\delta>0$, set
\[
F_{\pm,\delta}(\omega):=e^{-\delta\omega}\widehat h(\pm\omega),\qquad \omega>0.
\]
Then, by Tonelli's theorem for the regularized integral,
\[
\int_0^\infty\omega^{\mu-1}e^{-\delta\omega}e^{\mp i\omega\sigma}\,d\omega
=\G(\mu)(\delta\pm i\sigma)^{-\mu}.
\]
Since $h$ is bounded near $0$ and rapidly decreasing at infinity, the right-hand
side is integrable against $h(\sigma)$ after restricting first to
$\varepsilon<\sigma<R$. Letting $\varepsilon\downarrow0$ and $R\uparrow\infty$
therefore gives
\[
\mathscr M[F_{\pm,\delta}](\mu)
=\G(\mu)\int_0^\infty h(\sigma)(\delta\pm i\sigma)^{-\mu}
\,d\sigma .
\]
For fixed $\mu$ on any compact substrip of $0<\Re\mu<1$, the integrand is
bounded by $C\,|h(\sigma)|\min\{\sigma^{-\Re\mu},1\}$, which is integrable.
Dominated convergence as $\delta\downarrow0$ thus yields
\begin{equation}\label{E:MellinFourier}
\mathscr M[\widehat h(\pm\cdot)](\mu)
=\G(\mu)e^{\mp i\pi\mu/2}\mathscr Mh(1-\mu),
\qquad 0<\Re\mu<1,
\end{equation}
where the left-hand side is the Mellin transform of the Abel boundary value.

For the particular functions occurring here, the boundary values belong to the
weighted space underlying $E_\pm$. Indeed, integration by parts gives
$\widehat h(\omega)=O(|\omega|^{-1})$ as $|\omega|\to\infty$, while the
corresponding Abel limit of $\widehat{Wh}$ is obtained from the same formula
with the Mellin transform of $Wh$. Since $0<s<1$, the weight $\omega^s$ is
compatible with these endpoint bounds. We may therefore apply Mellin--Plancherel
on the line $\Re\mu=(1+s)/2$ and obtain
\[
E_\pm[h]
=\frac1{2\pi}\int_\R
\left|\G\!\left(\frac{1+s}{2}+i\eta\right)\right|^2
 e^{\pm\pi\eta}
\left|\mathscr Mh\!\left(\frac{1-s}{2}-i\eta\right)\right|^2d\eta.
\]
For $Wh$, Lemma~\ref{L:MellinW} gives
\[
\mathscr M(Wh)(1-\mu)=\mathscr Mh(\mu-s),
\]
so that, on the same Mellin line,
\[
E_-[Wh]
=\frac1{2\pi}\int_\R
\left|\G\!\left(\frac{1+s}{2}+i\eta\right)\right|^2
 e^{-\pi\eta}
\left|\mathscr Mh\!\left(\frac{1-s}{2}+i\eta\right)\right|^2d\eta.
\]
Replacing $\eta$ by $-\eta$ and using
$|\G(c-i\eta)|=|\G(c+i\eta)|$ transforms the last expression into $E_+[h]$.
This proves \eqref{E:MellinSwap}. The only limiting operation is the Abel limit
above; it is taken before Mellin--Plancherel and is controlled by the displayed
majorant on compact Mellin substrips.
\end{proof}

\begin{theorem}[History energy identity]\label{T:history}
Let $a=1-2s\in(-1,1)$, and let $u_0\in\Sm$. Then
\begin{equation}\label{E:freq}
\|U_0\|_{L_a^2(\RN)}^2
=
\frac{2^{1-2s}}{\G(s)^2}
\int_{\Rd}\int_0^\infty
\omega^s|\mathscr G_\xi(\omega)|^2\,d\omega\,d\xi.
\end{equation}
\end{theorem}

\begin{proof}
By Plancherel in $x$ and Tonelli,
\[
\|U_0\|_{L_a^2}^2
=\int_{\Rd}\int_0^\infty|\widehat U_0(\xi,z)|^2z^a\,dz\,d\xi.
\]
Lemma~\ref{L:MellinReduction} and Lemma~\ref{L:MellinW} give the right-hand
side as
\[
\frac{2^{1-2s}}{\G(s)^2}\int_{\Rd}E_-[Wg_\xi]d\xi.
\]
Lemma~\ref{L:MellinSwap} identifies $E_-[Wg_\xi]=E_+[g_\xi]$, which is precisely
\eqref{E:freq}.
\end{proof}

The frequency representation suggests the following intrinsic quadratic form.

\begin{definition}\label{D:Qs}
For histories $u,v$ for which the integral converges, set
\[
Q_s(u,v):=
\frac{2^{1-2s}}{\G(s)^2}
\int_{\Rd}\int_0^\infty
\omega^s\mathscr G^u_\xi(\omega)\overline{\mathscr G^v_\xi(\omega)}\,d\omega\,d\xi,
\qquad Q_s(u):=Q_s(u,u).
\]
For $u_0\in\Sm$, Theorem~\ref{T:history} gives
\[
Q_s(u_0)=\|\mathcal P_su_0\|_{L_a^2(\RN)}^2.
\]
\end{definition}

\begin{lemma}[One-sided Fourier uniqueness]\label{L:onesided-Fourier}
Let $g\in L^2(\mathbb R)$ be supported in $[0,\infty)$, and let
$\mathscr G=\mathscr Fg$ be its Fourier transform. If $\mathscr G$ vanishes on a
measurable subset of $\mathbb R$ of positive measure, then $g=0$ almost everywhere.
In particular, if $\mathscr G(\omega)=0$ for almost every $\omega>0$, then $g=0$.
\end{lemma}

\begin{proof}
Because $g$ is supported in $[0,\infty)$, its Fourier transform is the boundary
value of an $H^2(\mathbb C_-)$ function. Boundary uniqueness for Hardy functions
then gives $\mathscr G=0$, and Plancherel gives $g=0$.
\end{proof}

\begin{corollary}\label{C:positive}
For $u_0\in\Sm$, $Q_s(u_0)=0$ implies $u_0=0$ almost everywhere. Hence
$Q_s^{1/2}$ is a norm on $\Sm$.
\end{corollary}

\begin{proof}
If $Q_s(u_0)=0$, then the frequency representation in Definition~\ref{D:Qs}
shows that $\mathscr G_\xi(\omega)=0$ for a.e. $(\xi,\omega)$ with $\omega>0$.
For a.e. fixed $\xi$, Lemma~\ref{L:onesided-Fourier} gives $g_\xi=0$, hence
$u_0=0$.
\end{proof}

\begin{definition}\label{D:history}
We denote by $\mathcal H^s$ the completion of $\Sm$ with respect to the norm
$Q_s^{1/2}$.
\end{definition}

\begin{corollary}\label{C:extends}
The oscillatory Poisson lifting
$\mathcal P_s:\Sm\to L_a^2(\RN)$ extends uniquely to an isometry
\[
\mathcal P_s:\mathcal H^s\longrightarrow L_a^2(\RN).
\]
Indeed, if $u_0^{(k)}\to u_0$ in $\mathcal H^s$, then
$\mathcal P_su_0^{(k)}$ is Cauchy in $L_a^2(\RN)$ by Theorem~\ref{T:history}; its limit
is independent of the approximating sequence and satisfies
\[
\|\mathcal P_su_0\|_{L_a^2(\RN)}=\|u_0\|_{\mathcal H^s}.
\]
\end{corollary}



The preceding results identify the intrinsic energy space associated with
Schr\"odinger histories. We do not assert that the abstract completion
$\mathcal H^s$ has already been identified, as a set, with a concrete space of
past-time distributions; only the isometric realization through $\mathcal P_s$
is used below.

\section{Well-posedness of the nonlinear memory equation}\label{S:memory}
\label{S:wellpos}

Having completed the linear theory of the fully fractional
Schr\"odinger operator \(\Le^s\) and identified in Section
\ref{S:history} the intrinsic Hilbert space \(\mathcal H^s\) of
admissible histories, we now return to the nonlinear problem
\begin{equation}\label{E:nonlinear-memory}
\left\{
\begin{aligned}
\Le^s u(x,t)
&=
\mu |u(x,t)|^{p-1}u(x,t),
&& x\in\Rd,\quad t>0,
\\
u(x,t)
&=
u_0(x,t),
&& x\in\Rd,\quad t\leq0,
\end{aligned}
\right.
\end{equation}
where $d\ge1$ and
\[
u_0\in\mathcal H^s,\qquad \mu\in\C,\qquad p>1.
\]
The well-posedness results below concern only the ranges $1<p\le p_c$ specified in Theorem~\ref{T:memory-wp}; no claim is made here for $p>p_c$.

The fundamental point is that the prescribed history \(u_0\) has
already been converted, in a canonical way, into an initial datum for
the local half-space evolution. Indeed, by Corollary
\ref{C:extends}, the oscillatory Poisson lifting extends uniquely to
an isometry
\[
\mathcal P_s:\mathcal H^s
\longrightarrow
L_a^2(\RN),
\qquad a=1-2s.
\]
Thus, for every \(u_0\in\mathcal H^s\), there is a uniquely determined
\[
U_0=\mathcal P_su_0\in L_a^2(\RN)
\]
such that
\begin{equation}\label{E:history-isometry}
\|U_0\|_{L_a^2(\RN)}
=
\|u_0\|_{\mathcal H^s}.
\end{equation}

We use \(U_0\) as initial datum for the  singular
nonlinear Schr\"odinger evolution in the upper half-space:
\begin{equation}\label{cp02}
\begin{cases}
\p_t U - i(\Delta_x U + \Ba U) = 0, & \ \ X=(x,z)\in \mathbb{R}^{d+1}_+,\ t>0, \\
-\frac{2^{2s-1}\Gamma(s)e^{i \frac{\pi s}2}}
{\Gamma(1-s)} \displaystyle \lim_{z\to 0^+} z^a \p_z U(X,t) = \mu |U(x,0,t)|^{p-1}U(x,0,t), & \ \ x\in \Rd,\ t>0, \\
U(X,0)=U_0(X). 
\end{cases}
\end{equation}


Up to the explicit normalization in the boundary condition,
\eqref{cp02} is precisely the nonlinear Schr\"odinger equation with
nonlinear Neumann interaction studied in \cite{GS2}, and the critical and
subcritical regimes considered below are the ones introduced there in
\cite[Definition~2.7]{GS2}.

The two ranges of $a$ reach their thresholds for altogether different reasons.
When $0\le a<1$ the exponent $1+\frac{2(1-a)}{d+a+1}$ is the scaling-critical
one: it is the value of $p$ at which the parabolic scaling
$U\mapsto\la^{\frac{d+a+1}2}U(\la x,\la z,\la^2t)$, which preserves both
$L^2_a(\RN)$ and the boundary condition, leaves the problem invariant. When
$-1<a<0$ the situation is structurally different, and this is one of the
findings of that paper: the bulk propagator $\mathbb T^\star_a$ does not obey
these scalings at all, so that no scaling-critical exponent is available to
serve as a threshold. Its role is taken over by the dispersive behaviour of the
propagator, which in that range is the one of the $a=0$ problem, together with
the weighted two-exponent Strichartz theory that this forces; the resulting
threshold is $1+\frac{2}{d+1}$. The isometry \eqref{E:history-isometry}
transfers the size of the initial history to the weighted bulk norm, but it
does not by itself determine this second threshold. Thus we set
\begin{equation}\label{E:critical-exponents}
p_c=
\begin{cases}
\displaystyle
1+\frac{2(1-a)}{d+a+1},
& 0\leq a<1,
\\[3mm]
\displaystyle
1+\frac{2}{d+1},
& -1<a<0.
\end{cases}
\end{equation}

Here and below, given a Banach space $E$, we write $C^b_z([0,\infty);E)$ for
the space of bounded maps $z\mapsto F(\cdot,z,\cdot)$ from $[0,\infty)$ into
$E$ which are continuous up to and including $z=0$, normed by the supremum
over $z\ge0$; it is a closed subspace of $L^\infty_z E$, and membership in it
is what gives a pointwise meaning to the boundary value at $z=0$.

\begin{lemma}[Boundary trace of the free bulk evolution]\label{L:free-trace}
Let $I=[0,T]$ with $T<\infty$, and let $q,r$ be as in
Theorem~\ref{T:memory-wp}. Let $U_0\in L_a^2(\RN)$. If $0\le a<1$, then
\[
\mathbb T_a^\star(U_0)\in C^b_z\bigl([0,\infty);L_t^q(I;L_x^r(\Rd))\bigr);
\]
in particular $z\mapsto\mathbb T_a^\star(U_0)(\cdot,z,\cdot)$ extends
continuously to $z=0$ as an $L_t^q(I;L_x^r)$-valued map. If $-1<a<0$, the same
conclusion holds for the weighted free evolution
$k(z)\,\mathbb T_a^\star(U_0)$, in the mixed-norm framework of
\cite[Theorem~2.5]{GS2}; since $k(z)=1$ for $0\le z\le1$, this yields the same
boundary trace in $L_t^q(I;L_x^r)$.
\end{lemma}

\medskip
This statement is the free-evolution half of \cite[Theorem~8.2]{GS2}, to which
we refer for the companion assertion concerning the inhomogeneous term
$\mathbb D_a(F)$, not needed here. We include a proof because the form in which
the result is used below --- on a compact time interval $I$, and with the
weight $k$ in the anomalous range --- is the one we shall invoke repeatedly,
and because the argument is short.

\begin{proof}
We give the argument because the $L^\infty_z$ bounds in the homogeneous
Strichartz estimates do not by themselves imply continuity at the boundary.

\smallskip
\noindent\emph{A Neumann-compatible dense class.} We first work with data in
\[
\mathcal D_a:=\operatorname{span}
\Bigl\{\varphi(x)\,\psi(z)\ \Big|\
\varphi\in\mathscr S(\Rd),\
\mathcal H_{\frac{a-1}2}(\psi)\in C_0^\infty((0,\infty))\Bigr\},
\]
where $\nu=\frac{a-1}2$ and $\mathcal H_\nu$ is the Hankel transform
\eqref{E:hankel}. Three properties of this class will be used, and we record
them in the order in which they are needed.

\emph{(a) Its elements belong to $\mathscr S^+$, and therefore satisfy the
homogeneous Neumann condition \eqref{E:neumann}.} If $\mathcal H_\nu(\psi)=g\in
C_0^\infty((0,\infty))$, then by the inversion of \cite[Theorem~2.1]{GS1} we
have $\psi=\mathcal H_\nu(g)$, that is
$\psi(z)=\int_0^\infty(\lambda z)^{-\nu}J_\nu(\lambda z)\,g(\lambda)\,\lambda^a\,d\lambda$
with $g$ smooth and compactly supported away from the origin. Since
$w\mapsto w^{-\nu}J_\nu(w)$ is an even entire function, $\psi$ is a smooth
function of $z^2$; and repeated integration by parts in $\lambda$, using the
oscillation of $J_\nu$, shows that $\psi$ and all its derivatives decay
rapidly. Hence $\psi\in\mathscr S^+$, and by the property of that space recalled
after \eqref{E:hankeldiag} the elements of $\mathcal D_a$ satisfy
\eqref{E:neumann} and lie in $\Dom(\BaN)$.

\emph{(b) It is dense in $L^2_a(\RN)$.} The transform $\mathcal H_\nu$ is a
surjective isometry of $L^2(\R^+,z^adz)$ onto $L^2(\R^+,\lambda^ad\lambda)$, so
the functions $\psi$ occurring above are dense in $L^2(\R^+,z^adz)$; since
$\mathscr S(\Rd)$ is dense in $L^2(\Rd)$ and finite sums of tensor products are
dense in the Hilbert space tensor product
$L^2(\Rd)\otimes L^2(\R^+,z^adz)=L^2_a(\RN)$, the claim follows.

\emph{(c) It is invariant under $\Delta_x$, under $\BaN$ and under the
propagator $\mathbb T^\star_a(\cdot)(\cdot,\cdot,t)$.} The first two are clear
from \eqref{E:hankeldiag} and from $\mathscr S(\Rd)$ being invariant under
$\Delta_x$. For the third, by \eqref{E:multiplier} the propagator acts on the
Fourier--Hankel side as multiplication by
$e^{-it(4\pi^2|\xi|^2+\lambda^2)}$, which factors as a function of $\xi$ times a
function of $\lambda$: it therefore preserves the tensor structure, maps
$\mathscr S(\Rd)$ into itself in the first variable, and in the second leaves
the support of $\mathcal H_\nu(\psi)$ unchanged, that symbol being smooth and
of modulus one.

\smallskip
\noindent
We work with
$\mathcal D_a$ rather than with $C_0^\infty(\RN)$ because the latter is not
invariant under the propagator, whereas the argument below uses that the
quantities entering \eqref{E:conserved} are conserved in time.

\smallskip
\noindent\emph{A conserved quantity.} Let $U_0\in\mathcal D_a$,
write $U(\cdot,\cdot,t)=\mathbb T_a^\star(U_0)(\cdot,\cdot,t)$, the operator
$\mathbb T^\star_a$ being the one defined in \eqref{E:Tstar} and identified in
\eqref{E:Tstar-group}, and fix
$m\in\mathbb N$ with $4m>d$, so that the Sobolev embedding gives
$\|f\|_{L^r_x}\le C(d,m)\|(1-\Delta_x)^mf\|_{L^2_x}$ for every $2\le r\le\infty$.
Set $V(t):=(1-\Delta_x)^mU(\cdot,\cdot,t)$, so that $V(t)\in\mathcal D_a$ for
every $t$ by property (c) above. Since $\mathcal H_\nu$ diagonalizes
$\BaN$ by \eqref{E:hankeldiag}, it also diagonalizes the associated
quadratic form: for $\psi\in\Dom(\BaN)$,
\begin{equation}\label{E:formdiag}
\int_0^\infty|\p_z\psi|^2\,z^a\,dz
=\int_0^\infty\lambda^2\,|\mathcal H_\nu(\psi)(\lambda)|^2\,\lambda^a\,d\lambda .
\end{equation}
Together with Plancherel in the tangential variable, \eqref{E:formdiag} gives,
for every $t$,
\begin{equation}\label{E:conserved}
\|\p_zV(t)\|_{L^2_a(\RN)}^2
=\int_{\Rd}\!\int_0^\infty \lambda^2\,
\bigl|\mathcal H_{\nu}(\widehat{V(t)}(\xi,\cdot))(\lambda)\bigr|^2
\lambda^a\,d\lambda\,d\xi
=\|\p_z(1-\Delta_x)^mU_0\|_{L^2_a(\RN)}^2 ,
\end{equation}
the last equality because, by \eqref{E:multiplier}, the propagator acts on the
Fourier--Hankel side as multiplication by a symbol of modulus one --- in this
last step only the unimodularity is used, so that step is independent of the
normalisations chosen for the two transforms. We stress the point that will be
used repeatedly below: \emph{the right-hand side of \eqref{E:conserved} does not
depend on $t$}, so that $\|V(t)\|_{L^2_a(\RN)}$ and $\|\p_zV(t)\|_{L^2_a(\RN)}$
are finite and independent of $t$.

\smallskip
\noindent\emph{Continuity in $z$.} We estimate directly the difference between
two values of $z$; this gives continuity at once, with no splitting of the time
interval and with no restriction on the range of $z$.
Indeed, for $0\le z<z'<\infty$ and every $t$,
\[
\|V(\cdot,z',t)-V(\cdot,z,t)\|_{L^2_x}
\le\int_z^{z'}\|\p_\zeta V(\cdot,\zeta,t)\|_{L^2_x}\,d\zeta
\le\Bigl(\int_z^{z'}\zeta^{-a}d\zeta\Bigr)^{\frac12}
\Bigl(\int_0^\infty\|\p_\zeta V(\cdot,\zeta,t)\|_{L^2_x}^2\zeta^ad\zeta\Bigr)^{\frac12},
\]
by Cauchy--Schwarz with the weight split as $\zeta^{a/2}\cdot\zeta^{-a/2}$.
Since $a<1$ the first factor equals
$\bigl((z'^{\,1-a}-z^{\,1-a})/(1-a)\bigr)^{1/2}$, which is finite for all
$0\le z<z'<\infty$, the singularity of $\zeta^{-a}$ at the origin being
integrable precisely because $a<1$; the second factor is
independent of $t$ by \eqref{E:conserved}. Applying the Sobolev embedding in
$x$ we obtain, for a constant $C_1=C_1(d,a,m,U_0)>0$,
\begin{equation}\label{E:zHolder}
\sup_{t\in\R}\ \|U(\cdot,z',t)-U(\cdot,z,t)\|_{L^r_x(\Rd)}
\ \le\ C_1\,\bigl|z'^{\,1-a}-z^{\,1-a}\bigr|^{\frac12},
\qquad 0\le z,z'<\infty .
\end{equation}
We first observe that \eqref{E:zHolder} forces the quantity it estimates to be
finite for \emph{every} $z\ge0$, once it is finite for a single one. Indeed the
homogeneous estimate of \cite[Theorem~2.4]{GS2}, respectively of
\cite[Theorem~2.5]{GS2} in the anomalous range, bounds the essential supremum
in $z$ of $\|\mathbb T^\star_a(U_0)(\cdot,z,\cdot)\|_{L^q_t(I;L^r_x)}$,
respectively of its weighted counterpart, by $C\|U_0\|_{L^2_a}<\infty$; hence
there is at least one $z_1>0$, and in fact almost every $z_1>0$, at which
$\|U(\cdot,z_1,\cdot)\|_{L^q_t(I;L^r_x)}<\infty$, the weight $k$ being strictly
positive there. Since $\|\cdot\|_{L^q_t(I)}\le T^{1/q}\|\cdot\|_{L^\infty_t}$ on
the finite interval $I$, \eqref{E:zHolder} gives
\[
\|U(\cdot,z,\cdot)-U(\cdot,z_1,\cdot)\|_{L^q_t(I;L^r_x)}
\le T^{\frac1q}C_1\bigl|z^{1-a}-z_1^{1-a}\bigr|^{\frac12}<\infty
\]
for every $z\ge0$, and therefore
$\|U(\cdot,z,\cdot)\|_{L^q_t(I;L^r_x)}<\infty$ for every $z\ge0$.
Since $1-a>0$, the right-hand side of \eqref{E:zHolder} tends to $0$ as
$z'\to z$; hence
$z\mapsto U(\cdot,z,\cdot)$ is continuous from $[0,\infty)$ into
$L^\infty_t(\R;L^r_x)$, and therefore into $L^q_t(I;L^r_x)$ as well, because
$\|\cdot\|_{L^q_t(I)}\le T^{1/q}\|\cdot\|_{L^\infty_t}$ on the finite interval
$I=[0,T]$. Note that this part of the argument uses only $-1<a<1$, and so
applies verbatim in the anomalous range.

It follows that, for $U_0\in\mathcal D_a$, the map
$z\mapsto\mathbb T_a^\star(U_0)(\cdot,z,\cdot)$ is not merely continuous at
$z=0$ but belongs to $C^b_z\bigl([0,\infty);L^q_t(I;L^r_x)\bigr)$. Indeed,
when $0\le a<1$ the homogeneous estimate of \cite[Theorem~2.4]{GS2} bounds its
essential supremum in $z$ by $C\|U_0\|_{L^2_a}$, and for a map which is
continuous on $[0,\infty)$ the essential supremum coincides with the genuine
supremum. When $-1<a<0$ the same reasoning applied to the weighted estimate of
\cite[Theorem~2.5]{GS2} gives
$k\,\mathbb T^\star_a(U_0)\in C^b_z\bigl([0,\infty);L^q_t(I;L^r_x)\bigr)$,
the map $z\mapsto k(z)$ being continuous and positive on $(0,\infty)$ and
identically $1$ on $[0,1]$.

For general $U_0\in L_a^2(\RN)$, choose
$U_0^{(n)}\in \mathcal D_a$ with $U_0^{(n)}\to U_0$ in $L_a^2$.
When $0\le a<1$, the homogeneous estimate of \cite[Theorem~2.4]{GS2}, applied
with $F=0$ and vanishing Neumann datum to the admissible triple
$(q,r,\infty)$ of \eqref{E:admissible-positive} below, gives
\[
\operatorname*{ess\,sup}_{z\ge0}\ \|\mathbb T_a^\star(U_0^{(n)}-U_0^{(j)})(\cdot,z,\cdot)\|_{L_t^q(I;L_x^r)}
\le C\|U_0^{(n)}-U_0^{(j)}\|_{L_a^2}.
\]
Since, by the first part of the proof, each of the maps
$z\mapsto \mathbb T_a^\star(U_0^{(n)}-U_0^{(j)})(\cdot,z,\cdot)$ is continuous
on $[0,\infty)$, the essential supremum coincides with the genuine
supremum, and therefore
\begin{equation}\label{E:unifCauchy}
\sup_{z\ge0}\ \|\mathbb T_a^\star(U_0^{(n)}-U_0^{(j)})(\cdot,z,\cdot)\|_{L_t^q(I;L_x^r)}
\le C\|U_0^{(n)}-U_0^{(j)}\|_{L_a^2}.
\end{equation}
When $-1<a<0$, the corresponding estimate is the weighted one of
\cite[Theorem~2.5]{GS2}: with the two exponents $q,q_\infty$ of
\eqref{E:admissible-negative} below and the weight $k(z)=\min\{1,z^{a/2}\}$, it reads
\[
\operatorname*{ess\,sup}_{z\ge0}\ \bigl\|k(z)\,\mathbb T_a^\star(U_0^{(n)}-U_0^{(j)})(\cdot,z,\cdot)\bigr\|_{\left(L_t^{q}+L^{q_\infty}_t\right)(I;L_x^r)}
\le C\|U_0^{(n)}-U_0^{(j)}\|_{L_a^2},
\]
the target being the sum space $L^q_t(I)+L^{q_\infty}_t(I)$, which on the finite
interval $I$ coincides with $L^q_t(I)$ with equivalent norms, since
$q<q_\infty$. Arguing as above, we obtain \eqref{E:unifCauchy} in this range as
well, for the weighted maps $k\,\mathbb T_a^\star(U_0^{(n)}-U_0^{(j)})$.

In either case the approximating maps are continuous and bounded on
$[0,\infty)$ and form a uniformly Cauchy sequence with values in the Banach
space $L^q_t(I;L^r_x)$. Since $C^b_z\bigl([0,\infty);L^q_t(I;L^r_x)\bigr)$ is
complete, being a closed subspace of $L^\infty_zL^q_t(I;L^r_x)$, the limit lies
in it; and the same estimate applied to $U_0^{(n)}-U_0$ identifies that limit
with $\mathbb T_a^\star(U_0)$ when $0\le a<1$, and with
$k\,\mathbb T_a^\star(U_0)$ when $-1<a<0$. This proves the assertion, the
boundary value in the anomalous range being the unweighted one because
$k\equiv1$ on $[0,1]$.
\end{proof}


\subsection{The boundary formulation}

We first recall explicitly the part of the functional framework of
\cite{GS2} that is needed here. Set
\[
r=p+1.
\]
When \(0\leq a<1\), let \(q>2\) satisfy
\begin{equation}\label{E:admissible-positive}
\frac{2}{q}+\frac{d}{r}
=
\frac{d+a+1}{2}.
\end{equation}
On a time interval \(I=[0,T]\), the mild solutions of \eqref{cp02} belong to
\begin{equation}\label{E:GS2-positive-space}
C\bigl(I;L_a^2(\RN)\bigr)
\cap
C^b_z\bigl([0,\infty);
L_t^q(I;L_x^r(\Rd))\bigr).
\end{equation}
Membership in the first factor, and the bound in $L^\infty_zL^q_t(I;L^r_x)$,
are furnished by \cite[Theorems~2.4 and~2.9]{GS2}; the upgrade of that bound to
the space $C^b_z$, which is what gives the boundary value a pointwise meaning,
comes from Theorem~8.2 of \cite{GS2}, see also Lemma~\ref{L:free-trace}, for
the free bulk term, and from Theorem~8.1 of \cite{GS2} for the boundary
Duhamel term. We note that membership in $C^b_z$ is in any case built into the
notion of mild solution of \cite[Definition~2.6]{GS2}.
For $a\ge0$ this bulk admissibility relation is also the boundary
admissibility relation used in Theorem~8.1 of \cite{GS2}.  For $a<0$ the two
relations differ, as recorded below.  The boundary Duhamel term has a strong
trace by Theorem~8.1 of \cite{GS2}. Hence the full mild solution has a well-defined boundary value
\(U(\cdot,0,\cdot)\) in $L_t^q(I;L_x^r(\Rd))$, with
\begin{equation}\label{E:positive-trace}
U(\cdot,z,\cdot)
\longrightarrow
U(\cdot,0,\cdot)
\quad\text{in }L_t^q(I;L_x^r(\Rd))
\qquad\text{as }z\to0^+.
\end{equation}
When \(-1<a<0\), the theory of \cite{GS2} involves the weight
\[
k(z)=\min\{1,z^{a/2}\}
\]
and the exponents \(q,q_\infty\) determined by the equations:
\begin{equation}\label{E:admissible-negative}
\frac{2}{q_\infty}+\frac{d}{r}
=
\frac{d+a+1}{2},
\qquad
\frac{2}{q}+\frac{d}{r}
=
\frac{d+1}{2},
\qquad q>2.
\end{equation}
The corresponding mild solutions belong to the weighted mixed-norm
classes of \cite[Theorem~2.10]{GS2}. The boundary Duhamel term has a strong
trace in the boundary-admissible space $L_t^{q_\infty}L_x^r$ by
Theorem~8.1 of \cite{GS2}; since $q<q_\infty$ and $I$ is finite, this embeds into
$L_t^q(I;L_x^r)$. Therefore the full mild solution has a strong boundary
trace in the space needed below:
\begin{equation}\label{E:negative-trace}
U(\cdot,z,\cdot)
\longrightarrow
U(\cdot,0,\cdot)
\quad\text{in }L_t^q(I;L_x^r(\Rd))
\qquad\text{as }z\to0^+.
\end{equation}

We stress that the notation \(U(x,0,t)\) used below always denotes
the strong boundary trace furnished by
\eqref{E:positive-trace} or \eqref{E:negative-trace}; it is not a
pointwise evaluation at \(z=0\) of an oscillatory bulk integral.

We can now give the notion of solution that will be used for the
memory problem.

\begin{definition}[Mild solution of the memory problem]
\label{D:memory-solution}
Let $u_0\in\mathcal H^s$, let $U_0=\mathcal P_su_0$, and let $T>0$. Let
$\mu\in\C\setminus\{0\}$ and $1<p\le p_c$, set $r=p+1$, and let $q$ be
determined by \eqref{E:admissible-positive} when $0\le a<1$, respectively by
\eqref{E:admissible-negative} when $-1<a<0$.
A mild solution on $(0,T]$ with prescribed history $u_0$ is the pair $(u_0,u)$,
where
\[
u\in L_t^q((0,T);L_x^r(\Rd))
\]
is the strong boundary trace
\[
u(x,t)=U(x,0,t),\qquad 0<t\le T,
\]
of the mild solution $U$ of the half-space problem \eqref{cp02}, in the sense
of \cite{GS2}, with initial datum $U_0$. The trace is understood in the strong
mixed-norm sense of \eqref{E:positive-trace} or \eqref{E:negative-trace}.
The half-space mild solution with datum $U_0$ is unique in the class of
\cite[Definition~2.6]{GS2} in the subcritical regimes, and, at the critical
exponent $p=p_c$, in the ball of that class in which the fixed point of
\cite[Theorems~2.9(1) and~2.10(1)]{GS2} is produced. In either case the trace
$u$ is therefore determined by $u_0$, and uniqueness in
Theorem~\ref{T:memory-wp} is accordingly uniqueness within this
extension-induced class, with the same proviso at $p=p_c$. No claim is made
that a solution of the fractional equation arising in some other sense must be
of this form.
\end{definition}

\begin{remark}\label{R:concrete-history}
If $u_0\in\Sm$, then $u_0$ is an actual past-time function and one may form the
glued function
\[
\widetilde u(x,t)=
\begin{cases}
u_0(x,t),&t\le0,\\
u(x,t),&0<t\le T.
\end{cases}
\]
We do not assert that $\widetilde u$ solves any equation across $t=0$, nor that
it is continuous there: no compatibility between $u_0(\cdot,0)$ and
$u(\cdot,0^+)$ is available, and none is used. The point of the glueing is only
that, in this case, the prescribed past is an actual function and the pair
$(u_0,u)$ may be read as a single function on $\R^{d+1}$. In this sense the
abstract definition above agrees with the usual prescribed-history
interpretation whenever the history has a concrete representative. For a
general element of the completion $\mathcal H^s$, no pointwise evaluation on
$t\le0$ is used.
\end{remark}


This definition makes the role of the extension procedure explicit:
the prescribed past is first lifted to the bulk, the nonlinear
evolution is solved there, and the future evolution of the memory
problem is the boundary trace of the resulting half-space solution.

The next result identifies explicitly the nonlinear boundary term in
this formulation.

\begin{proposition}[Boundary Duhamel term]
\label{P:boundary-Duhamel}
Let \(U\) be a mild solution of \eqref{cp02}, and set
\[
u(x,t)=U(x,0,t),\qquad t>0.
\]
Then the boundary Duhamel contribution in the mild representation of
\(U\) is
\begin{equation}\label{E:boundary-Duhamel}
\frac{\mu}{\G(s)}
\int_0^t
\frac{1}{\sigma^{1-s}}
S(\sigma)
\bigl(
|u(\cdot,t-\sigma)|^{p-1}
u(\cdot,t-\sigma)
\bigr)(x)\,d\sigma .
\end{equation}
The identity holds in the space appearing in \eqref{E:positive-trace} or
\eqref{E:negative-trace}, the boundary value being the one furnished by
Theorem~8.1 of \cite{GS2}; see the proof, where it is stressed that no passage
to the limit under the integral sign is involved.
\end{proposition}

\begin{proof}
The generalized Duhamel formula of \cite{GS2}, with vanishing
interior forcing, gives
\begin{equation}\label{E:bulk-Duhamel}
U
=
\mathbb T_a^\star(U_0)
+
c_s\mu\,
\Theta_a^\star
\bigl(
|U(\cdot,0,\cdot)|^{p-1}U(\cdot,0,\cdot)
\bigr),
\end{equation}
where
\[
c_s
=
-\frac{2^{1-2s}\G(1-s)}
{\G(s)e^{i\pi s/2}}.
\]
By Theorem~\ref{T:id}, the boundary operator is represented by
\[
\Theta_a^\star(\Phi)(x,z,t)
=
-\frac{2^{2s-1}e^{i\frac{\pi s}{2}}}{\G(1-s)}
\int_0^t
\frac{e^{i\frac{z^2}{4\sigma}}}{\sigma^{1-s}}
S(\sigma)
\bigl(\Phi(\cdot,t-\sigma)\bigr)(x)\,d\sigma.
\]
We emphasise that the boundary value is \emph{not} obtained by letting
$z\to0^+$ under the integral sign. Such an interchange is not available here:
the phase $e^{iz^2/4\sigma}$ oscillates without bound as $\sigma\to0^+$, and no
dominating function is at hand. What is used instead is that the right-hand
side above is already meaningful at $z=0$, and that Theorem~8.1 of \cite{GS2}
identifies it with the boundary trace. In detail: for
$\Phi\in\mathscr S(\R^{d+1})$ the integral converges absolutely for every
$z\ge0$, the value $z=0$ included, since $\sigma\mapsto\sigma^{s-1}$ is
integrable at the origin for $s>0$ while
$\sup_{0<\sigma\le t}\|S(\sigma)(\Phi(\cdot,t-\sigma))\|_{L^r_x}<\infty$; the
boundary Strichartz estimate of \cite[Theorem~7.4]{GS2} then extends
$\Theta_a^\star$ by density to every $\Phi\in L^{q'}_tL^{r'}_x$; and
Theorem~8.1 of \cite{GS2} asserts that
\[
\lim_{z\to z_0}\
\bigl\|\Theta_a^\star(\Phi)(\cdot,z,\cdot)-\Theta_a^\star(\Phi)(\cdot,z_0,\cdot)
\bigr\|_{L^q_tL^r_x}=0
\qquad\text{for every } z_0\ge0,
\]
the value at $z_0$ being in each case the one furnished by that extension.
Taking $z_0=0$ therefore gives
\[
\Theta_a^\star(\Phi)(x,0,t)
=
-\frac{2^{2s-1}e^{i\frac{\pi s}{2}}}{\G(1-s)}
\int_0^t
\frac{1}{\sigma^{1-s}}
S(\sigma)
\bigl(\Phi(\cdot,t-\sigma)\bigr)(x)\,d\sigma.
\]
Taking
\[
\Phi=|u|^{p-1}u
\]
and using the normalization in \eqref{cp02} gives
\eqref{E:boundary-Duhamel}.
\end{proof}

\begin{corollary}[The boundary memory equation]\label{C:volterra}
Under the hypotheses of Proposition~\ref{P:boundary-Duhamel}, the boundary
trace $u$ of the mild solution of the memory problem satisfies, for $0<t\le T$,
the Volterra integral equation
\begin{equation}\label{E:volterra}
u(x,t)
=
\mathbb T_a^\star(\mathcal P_su_0)(x,0,t)
\;+\;
\frac{\mu}{\G(s)}
\int_0^t
\frac{1}{\sigma^{1-s}}\,
S(\sigma)\bigl(|u(\cdot,t-\sigma)|^{p-1}u(\cdot,t-\sigma)\bigr)(x)\,d\sigma ,
\end{equation}
the identity holding in $L^q_t(I;L^r_x(\Rd))$ in the sense described in
Proposition~\ref{P:boundary-Duhamel}.
\end{corollary}

\begin{proof}
Take the boundary trace of \eqref{E:bulk-Duhamel}, which exists by
Definition~\ref{D:memory-solution}, and replace the trace of the second
summand by \eqref{E:boundary-Duhamel}.

\end{proof}

\begin{remark}\label{R:volterra}
Equation \eqref{E:volterra} exhibits the memory structure of the problem in a
single formula: the entire influence of the prescribed past is carried by the
first term, the boundary trace of the free half-space evolution of the lifted
history $\mathcal P_su_0$, whereas the second term is a Volterra convolution in
which the weakly singular kernel $\sigma^{s-1}$ couples the solution
to its own past on $(0,t)$. We stress that \eqref{E:volterra} is a
\emph{consequence} of the extension-induced formulation of
Definition~\ref{D:memory-solution}, and is not used as its definition; in
particular we do not assert that \eqref{E:volterra} alone determines $u$.
\end{remark}

\begin{remark}\label{R:potential}
Formula \eqref{E:boundary-Duhamel} is precisely the nonlinear
fractional potential associated with the fundamental solution
\(\mathscr E_s\) of \(\Le^s\). Indeed, Theorem~\ref{T:id} identifies the
kernel of \(\Theta_a^\star\), up to the normalization appearing
above, with the oscillatory Poisson lifting of \(\mathscr E_s\). Thus the
boundary Duhamel term in the half-space theory agrees with the
nonlinear potential constructed in Section~\ref{S:inverse}. We stress that this
identification is available only in the range $0<s<\min\{1,d/2\}$ for which
$\mathscr E_s$ was constructed in \eqref{Es}; in particular, when $d=1$ it
excludes the whole anomalous range $-1<a<0$, that is $\frac12<s<1$. The
well-posedness theory below does not use the present remark.
\end{remark}

\begin{remark}\label{R:graph-formulation}
Definition~\ref{D:memory-solution} does not require a priori that
a concrete glued representative $\widetilde u$, when one is available as in
Remark~\ref{R:concrete-history}, belong to \(\Dom(\Le^s)\). Suppose that this
additional regularity does hold, that
\[
|\widetilde u|^{p-1}\widetilde u\in L^2\bigl(\Rd\times(0,\infty)\bigr),
\]
and, when $0<s<\frac12$, that the stronger graph-domain condition
$\widetilde u\in\Dom(\Le^{\frac14+\frac s2})$ of Theorem~\ref{T:DtNL2} is also
satisfied. Then that theorem identifies the weighted Neumann trace of the
extension with \(\Le^s\widetilde u\), and the mild formulation is consistent
with the graph-space equation
\[
\Le^s\widetilde u=\mu\,|\widetilde u|^{p-1}\widetilde u
\qquad\text{in}\ \ L^2\bigl(\Rd\times(0,\infty)\bigr).
\]
The restriction to $t>0$ is essential, and not a matter of convenience: for
$t\le0$ the function $\widetilde u$ is the prescribed history, and there
$\Le^s\widetilde u$ is an unconstrained element of $L^2$, which the problem
does not determine. The well-posedness result below does not require this
additional graph regularity.
\end{remark}

\subsection{Well-posedness}

We are now ready to transfer the nonlinear theory of \cite{GS2} to
the memory problem. The well-posedness statement below is formulated in the
extension-induced mild sense and does not require the graph regularity needed
for the Dirichlet-to-Neumann identification. When that additional regularity
holds, the mild solution also satisfies the fractional equation in $L^2$ for
$t>0$; for $0<s<\frac12$, this conclusion requires the stronger hypothesis in
Theorem~\ref{T:DtNL2}.

\begin{theorem}[Well-posedness of the nonlinear memory problem]
\label{T:memory-wp}
Let $0<s<1$, set $a=1-2s$, let $\mu\in\C\setminus\{0\}$, and let
$u_0\in\mathcal H^s$ be a prescribed history.

\smallskip
\noindent\emph{(i) The range $0\le a<1$.} Let
\[
p_c=1+\frac{2(1-a)}{d+a+1},
\]
suppose that $1<p\leq p_c$, set $r=p+1$, and let $q$ be determined by
\eqref{E:admissible-positive}. We note that $q>2$ is automatic: since
$2<r\le p_c+1=\frac{2(d+2)}{d+a+1}$, \eqref{E:admissible-positive} gives
$\frac2q\le\frac{d+a+1}{d+2}<1$, the last inequality because $a<1$.
\begin{itemize}
\item[(a)] If $p=p_c$, there exists
$\varepsilon_0=\varepsilon_0(d,a,c_s\mu)>0$ such that
\[
\|u_0\|_{\mathcal H^s}\leq\varepsilon_0
\]
implies that, for every $T<\infty$, the problem has a unique mild solution on
$(0,T]$ with prescribed history $u_0$, in the sense of
Definition~\ref{D:memory-solution}.
\item[(b)] If $1<p<p_c$, there exists
$T=T\bigl(d,a,p,\mu,\|u_0\|_{\mathcal H^s}\bigr)>0$ such that the problem has a
unique mild solution on $(0,T]$ with prescribed history $u_0$. If, in addition,
$\operatorname{Im}(c_s\mu)=0$ --- equivalently, by
Remark~\ref{R:admissible-mu} below, if $\mu\in e^{i\frac{\pi s}{2}}\R$ --- then this
solution extends to a mild solution on $(0,T]$ for every $T<\infty$, by
\cite[Theorem~2.9(2)]{GS2}, applied with vanishing interior forcing and
boundary coefficient $-c_s\mu$.
\end{itemize}

\smallskip
\noindent\emph{(ii) The range $-1<a<0$.} Let
\[
p_c=1+\frac{2}{d+1},
\]
suppose that $1<p\leq p_c$, set $r=p+1$, and let $q,q_\infty$ satisfy
\eqref{E:admissible-negative}.
\begin{itemize}
\item[(a)] If $p=p_c$, then for every $T<\infty$ there exists
$\varepsilon_1=\varepsilon_1(d,a,c_s\mu,T)>0$ such that
\[
\|u_0\|_{\mathcal H^s}\leq\varepsilon_1
\]
implies that the problem has a unique mild solution on $(0,T]$ with prescribed
history $u_0$. One may take
$\varepsilon_1=\varepsilon_0\bigl(C\max\{1,T^{\frac1q-\frac1{q_\infty}}\}\bigr)^{-1}$,
with $\varepsilon_0$ and $C$ as in \cite[Theorem~2.10(1)]{GS2}; since
$\frac1q-\frac1{q_\infty}>0$ in this range, $\varepsilon_1\to0$ as
$T\to\infty$, and no global statement is asserted here.
\item[(b)] If $1<p<p_c$, there exists $T>0$ such that the problem has a unique
mild solution on $(0,T]$ with prescribed history $u_0$.
\end{itemize}

\smallskip
\noindent
In both ranges uniqueness is understood within the extension-induced class of
Definition~\ref{D:memory-solution}; at the critical exponent $p=p_c$, in parts
(i)(a) and (ii)(a), it is uniqueness in the ball of that class in which the
fixed point of \cite[Theorems~2.9(1) and~2.10(1)]{GS2} is produced.
In the range $0\le a<1$ the local Lipschitz
dependence furnished by the fixed-point argument is quantified in
Proposition~\ref{P:contdep} below; no separate continuous-dependence statement
is asserted for $-1<a<0$.
\end{theorem}

\begin{proof}
By Corollary~\ref{C:extends},
\[
U_0=\mathcal P_su_0\in L_a^2(\RN),
\qquad
\|U_0\|_{L_a^2(\RN)}
=
\|u_0\|_{\mathcal H^s}.
\]
Consequently, the hypotheses on the initial datum in Theorems~2.9
and~2.10 of \cite{GS2} translate directly into the corresponding
hypotheses on the prescribed history $u_0$.

Those theorems provide existence and uniqueness of a mild solution
$U$ of the half-space problem \eqref{cp02} in the critical and
subcritical regimes stated above. In the subcritical range $0\le a<1$ and
when $\operatorname{Im}(c_s\mu)=0$, the continuation of $U$ beyond the initial
existence time follows from \cite[Theorem~2.9(2)]{GS2}. The free bulk component
has a strong boundary trace by Theorem~8.2 of \cite{GS2}, equivalently by
Lemma~\ref{L:free-trace}, while the boundary
Duhamel component has a strong trace by Theorem~8.1 of \cite{GS2}. Hence the
full solution has the strong boundary trace required in
Definition~\ref{D:memory-solution}, and its trace on $t>0$ defines the
corresponding mild memory solution. We stress that
Lemma~\ref{L:free-trace} is stated on a compact interval $I=[0,T]$; accordingly,
whenever the half-space solution is global in time, the memory solution is
obtained on $(0,T]$ for each $T<\infty$, which is the form in which the
conclusion is stated above. When $u_0\in\Sm$,
Remark~\ref{R:concrete-history} recovers the usual glued past--future representative.

Uniqueness here is uniqueness within the extension-induced class of
Definition~\ref{D:memory-solution}, and follows directly from uniqueness of
the half-space mild solution. A converse identification with arbitrary
graph-space solutions requires the additional regularity in
Remark~\ref{R:graph-formulation}. In the range $0\le a<1$, continuous dependence
is proved separately in Proposition~\ref{P:contdep}.

\end{proof}

\begin{remark}[The admissible coupling constants]\label{R:admissible-mu}
The hypothesis $\operatorname{Im}(c_s\mu)=0$, which is the condition under
which the $L^2_a$ mass of the extended solution is conserved, is \emph{not} the
familiar requirement that the coupling constant be real. Since
$\G(s),\G(1-s)>0$ for $0<s<1$, we may write
\[
c_s=-A_s\,e^{-i\frac{\pi s}{2}},
\qquad
A_s:=\frac{2^{1-2s}\G(1-s)}{\G(s)}>0,
\]
so that $\arg c_s=\pi-\frac{\pi s}{2}$, and consequently
\begin{equation}\label{E:admissible-mu}
\operatorname{Im}(c_s\mu)=0
\ \Longleftrightarrow\
\mu\in e^{i\frac{\pi s}{2}}\,\R .
\end{equation}
In particular, since $e^{i\frac{\pi s}{2}}\R\cap\R=\{0\}$ for every $0<s<1$, no
\emph{real} coupling constant $\mu\ne0$ satisfies this hypothesis; the whole of
the classical focusing and defocusing range $\mu=\pm c$, with $c>0$, is
excluded. The admissible coupling constants form the
real line rotated by $e^{i\frac{\pi s}{2}}$, the conjugate of the phase
$e^{-i\frac{\pi s}{2}}$ which relates $\Le^s$ to Samko's operator
$(\Delta_x+i\p_t)^s$, see the identity following \eqref{sam}. The rotation
disappears as $s\to0^+$, where one recovers the classical condition
$\mu\in\R$.
\end{remark}

We now prove the continuous dependence asserted in Theorem~\ref{T:memory-wp}.
The relevant half-space statement in the present regime is implicit in the fixed-point argument of
\cite[Theorem~2.9]{GS2}, but it is not stated there as a separate
Lipschitz-continuity theorem; that theorem asserts existence and uniqueness. We therefore include
the short difference estimate below, using only the linear estimate of
\cite[Theorem~2.4]{GS2} and the standard pointwise inequality for the
nonlinearity. Since the lifting \(\mathcal P_s\) is an isometry, this gives the
corresponding statement for histories.

For \(I=[0,T]\), set
\[
\mathcal X_T:=C\bigl(I;L^2_a(\RN)\bigr)\cap
C^b_z\bigl([0,\infty);L^q_t(I;L^r_x(\Rd))\bigr),
\]
with the weighted analogue \(Uk\in\mathcal X_T\),
\(k(z)=\min\{1,z^{a/2}\}\), in the range \(-1<a<0\).

\begin{proposition}[Lipschitz dependence]\label{P:contdep}
Let $0<s\le\frac12$, equivalently $0\le a<1$, and let $p$, $q$, $r$ be
as in part (i) of Theorem~\ref{T:memory-wp}.

\smallskip
\noindent\emph{(1) The subcritical case $1<p<p_c$.} For every $M>0$ there exist
$T=T(d,a,p,\mu,M)>0$ and $C=C(d,a,p,\mu,M)>0$, both independent of the
individual data within the ball of radius $M$, such that the following holds
with $I=[0,T]$.
\begin{itemize}
\item[(i)] \emph{(half-space problem)} if $U,\widetilde U$ are the mild
solutions of \eqref{cp02} on $I$ with initial data
$U_0,\widetilde U_0\in L^2_a(\RN)$ of norm at most $M$, then
\[
\|U-\widetilde U\|_{\mathcal X_T}
\leq C\|U_0-\widetilde U_0\|_{L^2_a(\RN)};
\]
\item[(ii)] \emph{(memory problem)} consequently, if $u,\widetilde u$ are the
mild solutions of \eqref{E:nonlinear-memory} corresponding to histories
$u_0,\widetilde u_0\in\mathcal H^s$ of norm at most $M$, then
\[
\|u-\widetilde u\|_{L^q_t(I;L^r_x(\Rd))}
\leq C\|u_0-\widetilde u_0\|_{\mathcal H^s}.
\]
\end{itemize}

\smallskip
\noindent\emph{(2) The critical case $p=p_c$.} The same two estimates hold for
every $M\le\varepsilon_0$, with $\varepsilon_0$ as in
Theorem~\ref{T:memory-wp}(i)(a), on $I=[0,T]$ for every $T<\infty$, with a
constant $C=C(d,a,p,\mu)$ independent of $T$.

\smallskip
\noindent
The restriction on $M$ in (2) is not technical: for $p=p_c$ and $M>\varepsilon_0$
Theorem~\ref{T:memory-wp} produces no solutions to be compared, and at the
scaling-critical exponent the existence time is in any case not a function of
the size of the datum alone. No separate Lipschitz estimate for the
$-1<a<0$ regime is asserted here; the two-exponent contraction used in
\cite[Theorem~2.10]{GS2} is not rewritten as a separate continuous-dependence
statement.
\end{proposition}

\begin{proof}
Let $U,\widetilde U$ be the corresponding half-space mild solutions, with
initial data $U_0=\mathcal P_su_0$ and
$\widetilde U_0=\mathcal P_s\widetilde u_0$. Subtracting the two Duhamel
representations gives
\[
U-\widetilde U
=\mathbb T_a^\star(U_0-\widetilde U_0)
+c_s\mu\,\Theta_a^\star\bigl(|u|^{p-1}u-|
\widetilde u|^{p-1}\widetilde u\bigr).
\]
By the linear estimates of \cite[Theorem~2.4]{GS2}, applied to the
homogeneous and boundary components of the Duhamel representation,
\[
\|U-\widetilde U\|_{\mathcal X_T}
\leq C\|U_0-\widetilde U_0\|_{L^2_a}
+C|c_s\mu|\,
\bigl\||u|^{p-1}u-|\widetilde u|^{p-1}\widetilde u\bigr\|_{L^{q'}_tL^{r'}_x}.
\]
For the nonlinear term we use
\[
\bigl||w|^{p-1}w-|v|^{p-1}v\bigr|
\leq C_p(|w|+|v|)^{p-1}|w-v|,
\qquad w,v\in\C,
\]
and H\"older's inequality with the exponents used in the self-map estimate
of \cite{GS2}. Since $r=p+1$ gives $r/p=r'$, we obtain
\[
\bigl\||u|^{p-1}u-|\widetilde u|^{p-1}\widetilde u\bigr\|_{L^{q'}_TL^{r'}_x}
\leq C\theta(T)
\bigl(\|u\|_{L^q_TL^r_x}+\|\widetilde u\|_{L^q_TL^r_x}\bigr)^{p-1}
\|u-\widetilde u\|_{L^q_TL^r_x},
\]
where $\theta(T)=1$ in the critical case and $\theta(T)=T^\kappa$ with
$\kappa>0$ in the subcritical case, and the boundary trace is controlled by
$\|U\|_{\mathcal X_T}$. In the subcritical case, the existence time supplied by
\cite[Theorem~2.9(2)]{GS2} in the present range $0\le a<1$ depends on the data
only through $\|U_0\|_{L^2_a}$, so that both solutions live on a common
interval $[0,T]$ with $T=T(M)$, and since $\theta(T)=T^\kappa\to0$ as
$T\to0^+$ we may shrink $T$, still as a function of $M$ alone, until the
nonlinear contribution is at most one half of the left-hand side. In the
critical case $\theta\equiv1$, no such gain in $T$ is available, and one uses
instead the smallness $M\le\varepsilon_0$: by
Theorem~\ref{T:memory-wp}(i)(a) the two solutions exist on $[0,T]$ for every
$T<\infty$ and satisfy $\|U\|_{\mathcal X_T},\|\widetilde U\|_{\mathcal X_T}\le
C\varepsilon_0$, so that the same absorption is achieved, uniformly in $T$, by
taking $\varepsilon_0$ small. Absorbing the nonlinear term yields
\[
\|U-\widetilde U\|_{\mathcal X_T}
\leq 2C\|U_0-\widetilde U_0\|_{L^2_a}.
\]
Finally, since $\mathcal P_s$ is an isometry,
\[
\|U_0-\widetilde U_0\|_{L^2_a}
=\|u_0-\widetilde u_0\|_{\mathcal H^s},
\]
and the boundary trace is bounded by the half-space norm. This proves (ii).

\end{proof}

\begin{remark}
Together with Theorem~\ref{T:memory-wp}, Proposition~\ref{P:contdep} gives
Hadamard well-posedness on $\mathcal H^s$ in the regime $0\le a<1$. In the
range $-1<a<0$ the theorem above gives existence and uniqueness, while this
paper does not assert a separate Lipschitz dependence statement.
\end{remark}

Theorem~\ref{T:memory-wp} completes the passage from the extension
theory developed in the present paper to the nonlinear evolution.
The prescribed history \(u_0\in\mathcal H^s\) is encoded canonically
by its Poisson lifting \(U_0=\mathcal P_su_0\); the nonlinear
evolution is solved in the upper half-space by the theory of
\cite{GS2}; and the strong boundary trace of that evolution defines
the corresponding mild solution of the memory problem. Proposition
\ref{P:boundary-Duhamel} shows, moreover, that the nonlinear boundary
term is exactly the fractional Duhamel potential associated with the
fundamental solution of \(\Le^s\).

%%%%%%%%%%%%%%%%%%%%%%%%%%


%%%%%%%%%%%%%%%%%%%%%%%%%%%%%%%%%%%%%%



%%%%%%%%%%%%%%%%%%%%%%%%%%%%
%%%%%%%%%%%%%%%%%%%%%%%%%%%%%%%%%


\section{Appendix}\label{S:app}

In this section we collect some known results which are useful in the main body of the paper; for the classical theory of Bessel functions we refer to \cite{Wa} and \cite{Le}. 
Consider the Bessel equation of order $\nu\in \mathbb C$
\begin{equation}\label{besseleq}
z^2 \frac{d^2 J}{dz^2} + z \frac{dJ}{dz} + (z^2 - \nu^2)J = 0.
\end{equation}
Assume $\nu\notin\mathbb N\cup\{0\}$. Then, two linearly independent solutions of \eqref{besseleq} are the Bessel function of the first kind and order $\nu$ 
\begin{equation}\label{besseries}
J_\nu(z) = \sum_{k=0}^\infty \frac{(-1)^k}{\G(k+1) \G(\nu+k+1)} \left(\frac z2\right)^{2k+\nu}\ \ \ \ \ \ \ \ |\arg z|<\pi,
\end{equation}
and the function 
\begin{equation}\label{negbes}
J_{-\nu}(z)=\sum_{k=0}^\infty \frac{(-1)^k}{\G(k+1) \G(k+1-\nu)} \left(\frac z2\right)^{2k-\nu}\ \ \ \ \ \ \ \ |\arg z|<\pi.
\end{equation}
From \eqref{besseries}, \eqref{negbes} we have as $z\to 0$, $|\arg z|<\pi$,
\begin{equation}\label{Js}
J_{\nu}(z) \cong \frac 1{\G(\nu+1)} \left(\frac{z}{2}\right)^{\nu},\ \ \ \ J_{-\nu}(z) \cong \frac 1{\G(1-\nu)} \left(\frac{z}{2}\right)^{-\nu}.
\end{equation}
The asymptotic behavior for large $z$ is much more delicate. One has for $\Re\nu>-\dfrac12$, and $0<\delta<\pi$
\begin{align}\label{jnuinfty} J_\nu(z)&=\sqrt{\frac2{\pi
z}}\cos\left(z-\frac{\pi\nu}2-\frac\pi4\right)+
O(z^{-\frac32})\\
&\quad\text{as }|z|\to\infty,\quad-\pi+\delta<\arg z<\pi-\delta.
\notag
\end{align}
We also note that the functions
\begin{equation}\label{FG}
F_\nu(z) = z^\nu J_\nu(z),\ \ \ \ \ G_\nu(z) = z^{-\nu} J_{\nu}(z),
\end{equation}
satisfy the following recurrence relations, see \cite[(5.3.5), p.103]{Le} (note that, to obtain the second relation in \eqref{FG}, we have changed $\nu$ into $-\nu$ in the second relation in (5.3.5). This is allowed since (5.3.5) is valid for arbitrary $\nu\in \C$)
\begin{equation}\label{recu}
\frac{d}{dz} F_\nu(z) = z^\nu J_{\nu-1}(z),\ \ \ \ \ \frac{d}{dz} G_{\nu}(z) = - z^{-\nu} J_{1+\nu}(z).
\end{equation}
If we let $\mu = -\nu$, then the second relation in \eqref{recu} gives
\begin{equation}\label{recubis}
\frac{d}{dz} [G_{-\nu}(z)] = \frac{d}{dz}[z^\nu J_{-\nu}(z)] = - z^\nu J_{1-\nu}(z).
\end{equation}
Consider the generalized Bessel equation
\begin{equation}\label{genbessel}
z^2 \Psi''(z) + (1 - 2\alpha) z \Psi'(z) + \left[\beta^2 \gamma^2
z^{2\gamma} + (\alpha^2 - \nu^2 \gamma^2)\right] \Psi(z) = 0.
\end{equation}
If $\Phi(z)$ is a solution to 
\eqref{besseleq}, then the function defined by the
transformation
\begin{equation}\label{besselcv}
\Psi(z) = z^\alpha \Phi(\beta z^\gamma)
\end{equation}
solves the equation \eqref{genbessel}.

Consider now the modified Bessel equation of order $\nu\in \mathbb C$
\begin{equation}\label{modbesseleq}
z^2 \frac{d^2\Phi}{dz^2} + z \frac{d\Phi}{dz} - (z^2 + \nu^2)\Phi = 0.
\end{equation}
For $\nu\not\in \mathbb N\cup\{0\}$, two linearly independent solutions of \eqref{modbesseleq} are the modified Bessel function of the first kind, 
\begin{equation}\label{Inu}
I_\nu(z) = \sum_{k=0}^\infty \frac{(z/2)^{\nu+2k}}{\G(k+1) \G(k+\nu+1)},\ \ \ \ \ \ \  \ \ |\arg(z)| < \pi,
\end{equation}
and the modified Bessel function of the third kind, or Macdonald function, which is defined by
\begin{equation}\label{Knu}
K_\nu(z) = \frac \pi{2} \frac{I_{-\nu}(z) - I_\nu(z)}{\sin \pi \nu},\ \ \ \ \ \ \ \ \ \ \  \ |\arg(z)| < \pi.
\end{equation}
Notice that 
\begin{equation}\label{same}
K_\nu(z) = K_{-\nu}(z).
\end{equation}
From \eqref{Inu} we have as $z\to 0$, $|\arg z|<\pi$,
\begin{equation}\label{Is}
I_{\nu}(z) \cong \frac 1{\G(\nu+1)} \left(\frac{z}{2}\right)^{\nu},\ \ \ \ I_{-\nu}(z) \cong \frac 1{\G(1-\nu)} \left(\frac{z}{2}\right)^{-\nu}.
\end{equation}
If $\nu> 0$, we infer from \eqref{Knu} and \eqref{Is}
\[
K_{\nu}(z) \cong 2^{\nu-1} \frac{\pi}{\sin(\pi \nu)} \frac{z^{-\nu}}{\G(1-\nu)}.
\]
Combining this asymptotic with the formula, see \cite[(1.2.2) p.3]{Le},  
\begin{equation}\label{sine}
\G(z) \G(1-z) = \frac{\pi}{\sin(\pi z)},
\end{equation} 
we infer when $\nu> 0$
\begin{equation}\label{zeroKnu}
K_{\nu}(z) \cong 2^{\nu-1} \G(\nu) z^{-\nu},\ \ \ \ \ z\to 0,\ \ |\arg z|<\pi.
\end{equation}
The following recurrence relation holds, see \cite[(5.7.9), p.110]{Le}
\begin{equation}\label{recuu}
\frac{d}{dz} [z^\nu K_\nu(z)] = - z^\nu K_{\nu-1}(z) = -  z^\nu K_{1-\nu}(z),
\end{equation}
where in the last equality we have used \eqref{same}.
If we consider the generalized modified Bessel equation
\begin{equation}\label{modgenbessel}
z^2 \Psi''(z) + (1 - 2\alpha) z \Psi'(z) + \left[(\alpha^2 - \nu^2 \gamma^2)- \beta^2 \gamma^2
z^{2\gamma}\right] \Psi(z) = 0,
\end{equation}
if the function $\Phi$ solves \eqref{modbesseleq}, then the function defined by \eqref{besselcv}
is a solution to \eqref{modgenbessel}.
As $z\to \infty$, we have 
\begin{equation}\label{ab}
I_\nu(z) =  \frac{e^z}{(2\pi z)^{1/2}} \left(1+ O(|z|^{-1})\right),
\end{equation}
whereas $K_\nu$ decays exponentially at infinity according to the asymptotic formula
\begin{equation}\label{abK}
K_\nu(z) = \left(\frac{\pi}{2z}\right)^{1/2} e^{-z}  \left(1+ O(|z|^{-1})\right), \ \ \ \ \ |\arg z| \le \pi - \delta,
\end{equation}
see \cite[(5.11.9) on p.123]{Le}. Finally, we use the following integral representation 
\begin{equation}\label{intK}
K_\nu(z) = 2^{-\nu-1} z^\nu \int_0^\infty e^{-\tau} e^{-\frac{z^2}{4\tau}} t^{-\nu-1} d\tau,\ \ \ \ \ |\arg z|<\frac{\pi}4,
\end{equation}  
see \cite[(5.10.25), p.119]{Le}.

%%%%%%%%%%%%%%%%%%%%%%%%%%%%%%%%%%%%%%%
\subsection{Some oscillatory integral identities}\label{S:osc2}

In this section we collect some  oscillatory integrals which are needed in the main body of the paper. Throughout, complex powers and square roots are taken on the principal
branch:
\[
w^\alpha=\exp(\alpha \operatorname{Log} w),
\qquad
\operatorname{Log} w=\log|w|+i\operatorname{Arg} w,
\qquad
-\pi<\operatorname{Arg} w<\pi.
\]
We begin with a basic formula.
 
 \begin{lemma}\label{L:aux}
 Let $A\in \R\setminus\{0\}$ and $0<s<1$. Then
 \begin{equation}\label{As}
 -\frac{s}{\G(1-s)} \int_0^\infty \frac{1}{\sigma^{1+s}} [e^{-i \sigma A} - 1] d\sigma = (i A)^s = |A|^s e^{i\frac{\pi s}{2} \operatorname{sgn}A}.
 \end{equation}
 \end{lemma}
 
 \begin{proof}
We recall the following well-known integral:
\[
-\frac{s}{\G(1-s)} \int_0^\infty \frac{1}{\sigma^{1+s}} [e^{-\sigma L } - 1] d\sigma = L^s,\ \ \ \ \ \ \Re L>0.
\]
Applying this formula with $L = \ve+iA$, $\ve>0$, and letting $\ve\to 0^+$ we reach the desired conclusion \eqref{As}.

 \end{proof}
 
We will also need the following classical oscillatory integral.

\begin{lemma}\label{L:cis}
Let $b>0$ and $0<\Re \mu<1$. Then
\begin{equation}\label{ei}
\int_0^\infty y^{\mu-1} e^{\pm i by} dy = \frac{\G(\mu)}{b^\mu} e^{\pm i\frac{\mu \pi}{2}}.
\end{equation} 
\end{lemma}

\begin{proof}
The proof of \eqref{ei} directly follows from:
\begin{equation}\label{sincosy}
\int_0^\infty y^{\mu-1} \sin(by) dy = \frac{\G(\mu)}{b^\mu} \sin \frac{\mu \pi}{2},\ \ \ \ \ \ \int_0^\infty y^{\mu-1} \cos(by) dy = \frac{\G(\mu)}{b^\mu} \cos \frac{\mu \pi}{2},
\end{equation}
see \cite[formulas 3.761.4 \& 3.761.9, p.420-1]{GR}.

\end{proof}

We also need the following important consequence of Lemma \ref{L:cis}.

\begin{lemma}\label{L:Abel}
Let $0<s<1, t>0$, and $z>0$. Then the following oscillatory identity holds:
\begin{equation}\label{abel}
\int_0^t
\frac{e^{i \frac{z^2}{4\sigma}}}{(t-\sigma)^{1-s}\sigma^{1+s}}\,d\sigma =
4^s\,\G(s)\,e^{i\frac{\pi s}{2}}\,t^{s-1}\,z^{-2s}
e^{i \frac{z^2}{4t}}.
\end{equation}
\end{lemma}

\begin{proof}
Since the factor $\sigma^{-{(1+s)}}$ is not integrable at $\sigma=0$, we interpret the integral $I(t,z)$ in the left-hand side of \eqref{abel}
as a generalized Riemann integral in which we first make the change of variable $\sigma = t u$, and then $\tau = u^{-1}$, obtaining
\begin{align*}
I(t,z) & = \underset{\ve\to 0^+}{\lim} \int_{\ve}^t \frac{e^{i\frac{z^2}{4\sigma}}}
{(t-\sigma)^{1-s}\sigma^{1+s}}\,d\sigma = \frac{1}{t}\underset{\ve\to 0^+}{\lim}  \int_{\ve/t}^1\frac{e^{i\frac{z^2}{4t u}}}
{(1-u)^{1-s} u^{1+s}}\,du
\\
& = \frac{1}{t} \underset{\ve\to 0^+}{\lim} \int_1^{t/\ve}\frac{e^{i\frac{z^2}{4t }\tau}}
{(\tau-1)^{1-s}}\,d\tau = \frac{e^{i\frac{z^2}{4t }}}{t} \underset{\ve\to 0^+}{\lim} \int_0^{t/\ve -1} y^{s-1}e^{i\frac{z^2}{4t } y}\,dy.
\end{align*}
If we now apply \eqref{ei} with $\mu = s$, we reach the desired conclusion.

\end{proof}


%%%%%%%%%%%%%%%%%%%%%%%%%%%%%%%%%%%%%%%%%%%%%%%%%%%%%%%%









%%%%%%%%%%%%%%%%%%%%%%%%%%%%%%%%%%%%%%%%%%%%%%%%%
\bibliographystyle{amsplain}


\begin{thebibliography}{10}

\bibitem{CS}
L. Caffarelli \& L. Silvestre, \emph{An extension problem related to the fractional Laplacian}.
Comm. Partial Differential Equations 32~(2007), no. 7-9, 1245-1260.

\bibitem{CM}
M. Caputo \& F. Mainardi,
\emph{Linear models of dissipation in anelastic solids}. Rivista del Nuovo Cimento \textbf{1} (1971), 161--198.

\bibitem{CN}
A. P. Chegolin \& V. A. Nogin, 
\emph{Integral transforms related to complex powers of the generalized Schr\"odinger operator}.
Integral Transforms Spec. Funct. 17~(2006), no. 6, 409-420.

\bibitem{DGLZ}
Q. Du, M. Gunzburger, R. B. Lehoucq \& K. Zhou,
\emph{Analysis and approximation of nonlocal diffusion problems with volume constraints}. SIAM Review \textbf{54} (2012), 667--696.

\bibitem{EGV}
A. Enciso, M. del Mar Gonz\'alez \& B. Vergara, \emph{
Fractional powers of the wave operator via Dirichlet-to-Neumann maps in anti–de Sitter spaces}.
J. Funct. Anal. 273~(2017), no. 6, 2144-2166.

\bibitem{FD}
 R. Ferreira \& A. de Pablo, \emph{Blow-up for a fully fractional heat equation}.
Discrete Contin. Dyn. Syst. 44~(2024), no. 2, 569-584.

\bibitem{FB}
Y. Four\`es-Bruhat,  \emph{Solution \'el\'ementaire d'\'equations ultrahyperboliques}.
J. Math. Pures Appl. (9) 35~(1956), 277-288.

\bibitem{GMS}
J. E. Gal\'e, P. J. Miana \& P. R. Stinga,
\emph{Extension problem and fractional operators: semigroups and wave equations}.
J. Evol. Equ. 13~(2013), no. 2, 343-368.

\bibitem{GS1}
N. Garofalo \& G. Staffilani, \emph{Strichartz estimates for a Schr\"odinger equation on the half-line with a Neumann boundary condition}. To appear in Journal of Functional Analysis. Preprint available in arXiv:2507.19998

\bibitem{GS2}
N. Garofalo \& G. Staffilani, \emph{Strichartz estimates for Schr\"odinger equations with nonlinear boundary interactions}. Preprint available in arXiv:2605.23229

\bibitem{GR}
I. S. Gradshteyn \& I. M. Ryzhik, \emph{Tables of integrals, series,
and products}, Academic Press, 1980.

\bibitem{Jo}
B. F. Jones, Jr., \emph{Lipschitz spaces and the heat equation}.
J. Math. Mech. 18~(1968/69), 379-409.

\bibitem{Laskin}
N. Laskin, \emph{Fractional Schr\"odinger equation}. Phys. Rev. E (3) \textbf{66} (2002), no. 5, 056108, 7 pp.

\bibitem{Le}
N. N. Lebedev, \emph{Special functions and their applications}, Revised edition, translated from the Russian and edited by R. A. Silverman. Unabridged and corrected republication. Dover Publications, Inc., New York, 1972. xii+308 pp.

\bibitem{Mainardi}
F. Mainardi, \emph{Fractional calculus and waves in linear viscoelasticity: an introduction to mathematical models}. Imperial College Press, London, 2010.

\bibitem{MK}
R. Metzler \& J. Klafter,
\emph{The random walk's guide to anomalous diffusion: a fractional dynamics approach}. Physics Reports \textbf{339} (2000), 1--77.

\bibitem{MS}
B. Muckenhoupt \& E. M. Stein, \emph{Classical expansions and their relation to conjugate harmonic functions}. 
Trans. Amer. Math. Soc. \textbf{118}~(1965), 17-92. 

\bibitem{Naber}
M. Naber, \emph{Time fractional Schr\"odinger equation}. Journal of Mathematical Physics \textbf{45} (2004), 3339--3352.

\bibitem{NS}
V. A. Nogin \& E. V. Sukhinin, \emph{Fractional powers of the Schr\"odinger operator in $L^p$-spaces}. Dokl. Akad. Nauk 360~(1998), no. 5, 608-610.

\bibitem{NyS}
K. Nystr\"om \& O. Sande, \emph{Extension properties and boundary estimates for a fractional heat operator}.
Nonlinear Anal. 140 ~(2016), 29-37.

\bibitem{R}
M. Riesz, \emph{Int\'egrales de Riemann-Liouville et potentiels}, Acta Sci. Math. Szeged, 9~(1938), 1-42.

\bibitem{R2}
M. Riesz, \emph{L'int\'egrale de Riemann-Liouville et le probl\`eme de Cauchy}, (French) Acta Math. 81~(1949), 1-223.

\bibitem{Sa}
S. G. Samko, \emph{Hypersingular integrals and their applications}.
Anal. Methods Spec. Funct., 5
Taylor \& Francis Group, London, 2002, xviii+359 pp.

\bibitem{Samp}
C. H. Sampson, \emph{A characterization of parabolic Lebesgue spaces}. 
Thesis (Ph.D.)-Rice University. 1968. 91 pp.

\bibitem{St}
E. M. Stein, \emph{Singular integrals and differentiability properties of functions}.
Princeton Math. Ser., No. 30
Princeton University Press, Princeton, NJ, 1970, xiv+290 pp.

\bibitem{ST}
P. R. Stinga \& J. L. Torrea, \emph{Regularity theory and extension problem for fractional nonlocal parabolic equations and the master equation}.
SIAM J. Math. Anal. 49 ~(2017), no. 5, 3893-3924.

\bibitem{Wa}
G. N. Watson,  \emph{A treatise on the theory of Bessel functions}.
Cambridge Math. Lib.
Cambridge University Press, Cambridge, 1995, viii+804 pp.

\bibitem{We48}
A. Weinstein,  \emph{Discontinuous integrals and generalized potential theory.} Trans. Amer. Math. Soc. 63 (1948), 342--354.

\bibitem{We2}
A. Weinstein,  \emph{On a singular differential operator}.
Ann. Mat. Pura Appl. (4) 49~(1960), 359--365.



\end{thebibliography}


\end{document}
